% concatenated from main.tex
\documentclass[a4paper,11pt]{article}
\usepackage[utf8]{inputenc}
\usepackage[T1]{fontenc}

\usepackage{boxedminipage}
\usepackage{amsfonts}
\usepackage{amsmath} 
\usepackage{amssymb}
\usepackage{graphicx}
\usepackage{amsthm}
\usepackage{subfig}
\usepackage{cancel}
\usepackage{textcmds}
\usepackage{microtype}
\usepackage{hyperref}
\emergencystretch=2em

% Large asterisk used for forcing iterations.
\DeclareFontFamily{U}{mathb}{\hyphenchar\font45}
\DeclareFontShape{U}{mathb}{m}{n}{
      <5> <6> <7> <8> <9> <10> gen * mathb
      <10.95> mathb10 <12> <14.4> <17.28> <20.74> <24.88> mathb12
      }{}
\DeclareSymbolFont{mathb}{U}{mathb}{m}{n}
\DeclareMathSymbol{\Asterisk}{1}{mathb}{"06}


 

\newtheorem{theorem}{Theorem}[section]
\newtheorem{lemma}[theorem]{Lemma}
\newtheorem{proposition}[theorem]{Proposition}
\newtheorem{corollary}[theorem]{Corollary}
\newtheorem{definition}[theorem]{Definition}
\newtheorem{example}[theorem]{Example}
\newtheorem{fact}[theorem]{Fact}
\newtheorem{question}{Question}
\newtheorem*{claim}{Claim}
\newtheorem*{theorem*}{Theorem}

\newcommand{\forceP}{\mathbb{P}}
\newcommand{\forceQ}{\mathbb{Q}}
\newcommand{\forceR}{\mathbb{R}}
\newcommand{\forceL}{\mathbb{L}}
\newcommand{\forceA}{\mathbb{A}}
\newcommand{\forceB}{\mathbb{B}}
\newcommand{\forceC}{\mathbb{C}}
\newcommand{\forceS}{\mathbb{S}}
\newcommand{\forceT}{\mathbb{T}}
\newcommand{\forceU}{\mathbb{U}}
\newcommand{\key}{\operatorname{key}}
\newcommand{\pcode}{\operatorname{pcode}}


\newcommand{\ZFC}{\mathsf{ZFC}}
\newcommand{\ZF}{\mathsf{ZF}}
\newcommand{\ZFCminus}{\mathsf{ZFC}^{-}}
\newcommand{\AC}{\mathsf{AC}}
\newcommand{\PFA}{\mathsf{PFA}}
\newcommand{\BPFA}{\mathsf{BPFA}}
\newcommand{\MRP}{\mathsf{MRP}}
\newcommand{\MM}{\mathsf{MM}}
\newcommand{\CH}{\mathsf{CH}}
\newcommand{\GCH}{\mathsf{GCH}}
\newcommand{\MA}{\mathsf{MA}_{\omega_1}}
\newcommand{\PD}{\mathsf{PD}}
\newcommand{\mouseM}{\mathcal{M}}
\newcommand{\mouseN}{\mathcal{N}}



\def\undertilde#1{\mathord{\vtop{\ialign{##\crcr
$\hfil\displaystyle{#1}\hfil$\crcr\noalign{\kern1.5pt\nointerlineskip}
$\hfil\tilde{}\hfil$\crcr\noalign{\kern1.5pt}}}}}

%opening
\title{Forcing the ${\Pi^1_n}$-Uniformization Property}
\author{Stefan Hoffelner\footnote{The author's research was funded in parts by the Austrian Science Fund (FWF) Grant-DOI 10.55776/P37228. This research was also funded in parts by the Deutsche Forschungsgemeinschaft (DFG German Research Foundation) under Germany's Excellence Strategy EXC 2044 390685587, Mathematics M\"unster: Dynamics--Geometry--Structure. For the purpose of open access, the author has applied a CC BY public copyright license to any Author Accepted Manuscript version arising from this submission. } }
\date{\today}
\begin{document}

\maketitle

\begin{abstract}
We generically construct a model in which the ${\Pi^1_3}$-uniformization property is true, thus lowering the best known consistency strength from the existence of $M_1^{\#}$ to just $\ZFC$. The forcing construction can be adapted to work over canonical inner models with Woodin cardinals, which yields, for the first time, universes where the $\Pi^1_{2n}$-uniformization property holds, thus producing models which contradict the natural $\PD$-induced pattern.

\end{abstract}

\section{Introduction}
The question of finding nicely definable 
choice functions for a definable family of sets is an old and well-studied subject in descriptive set theory; see, for instance, \cite{Kechris,Moschovakis}. The uniformization problem, first mentioned by N. Lusin in 1930 (see \cite{Lusin}), asks to find choice functions which lie at the same projective level as the set they aim to uniformize.
Recall that for $A\subseteq 2^{\omega}\times 2^{\omega}$, a
\emph{uniformization} of $A$ is a function
\[
   f:\operatorname{pr}_1(A)\longrightarrow 2^\omega
\]
such that
\[
   \operatorname{graph}(f)\subseteq A.
\]
Equivalently, for every $x\in\operatorname{pr}_1(A)$ one has
$(x,f(x))\in A$.


\begin{definition}
 We say that a pointclass $\Gamma$ has the uniformization
 property iff every element of $\Gamma$ admits a uniformization 
 in $\Gamma$.
\end{definition}

It is a classical result due to M. Kondo \cite{Kondo} that lightface $\Pi^1_{1}$-sets do have the 
uniformization property, this also yields the uniformization property for $\Sigma^1_{2}$-sets. This is as much as $\ZFC$ can prove about uniformization.
In the constructible universe $L$, as shown by J. Addison in \cite{Addison}, $\Sigma^1_{n}$ does have the uniformization property for $n \ge 3$, which follows from the existence of a $\Sigma^1_{2}$-good wellorder of the reals; for background on projective wellorders see also \cite{CS,SyVera}. Thus the $\Pi^1_{n}$-uniformization fails for $n \ge 3$. On the other hand, by the celebrated results of Y. Moschovakis (see \cite{Moschovakis2}, Theorem 1), 
$\bf{\Delta}^1_{2n}$-projective determinacy implies ${\Pi}^1_{2n+1}$-uniformization, yet the determinacy assumption exceeds in logical strength $\ZFC$. It is known due to H. W. Woodin (see \cite{MSW},), that $\bf{\Delta}^1_2$-projective determinacy implies the existence of $M_1^{\#}$, hence yields an inner model with a Woodin cardinal. 
As with other regularity properties of the reals like Lebesgue measurability or Baire property, which both follow from $\PD$ as well, it is natural to ask whether the $\Pi^1_3$-uniformization property bears large cardinal strength as well. We shall answer it negatively.


\begin{theorem*}
There is a generic extension of $L$ in which the $\Pi^1_3$-uniformization property is true. 
\end{theorem*} 
The proof can be adapted such that it applies to canonical inner models with Woodin cardinals. This can be used to obtain better lower bounds in terms of consistency strength for the $\Pi^1_n$ uniformization property for odd $n$. For even $n$ we can produce for the first time models where the $\Pi^1_n$ uniformization property holds true.
\begin{theorem*}
Let $M_n$ be the canonical inner model with $n$ Woodin cardinals. Then there is a generic extension of $M_n$ in which the $\Pi^1_{n+3}$ uniformization property holds true.
\end{theorem*}

Since the first version of this paper was written, the method developed here has led to several further constructions.  A variant combines $\Pi$-uniformization with a projective well-order of the reals under $\CH$ \cite{HoffelnerPiUniformizationWellorder}; another combines global $\Sigma$-uniformization with a large continuum and a projective well-order \cite{HoffelnerLargeContinuumGlobalSigma}; and the local Delfino constructions combine related coding and allowability methods with projective regularity, definable well-orders, and adjacent $\Sigma$- or $\Pi$-uniformization \cite{HoffelnerLocalDelfinoSigma,HoffelnerLocalDelfinoPi}.  These later papers should be viewed as refinements and applications of the forcing machinery introduced here.  The present paper isolates the original fixed-point and allowability mechanism needed to force projective $\Pi$-uniformization without assuming the corresponding determinacy strength.

Questions concerning the forceability of (local) consequences of $\PD$ do have a long tradition in set theory. There is a vast body of literature concerning the forcability of local levels of the projective hierarchy satisfying (Boolean combinations of) the Baire property, the perfect set property or Lebesgue measurability; see for instance \cite{FS}. There has been very little progress in the past, however, concerning similar questions for the separation, the reduction and the uniformization property. Indeed, even the question of whether one can force the $\Sigma^1_3$-separation property, which is the weakest of said properties, remained an open problem for 50 years (see \cite{Mathias}, Problem 3029).

This article continues this line of research and introduces the forcing template on which the subsequent work mentioned above builds. It is organized as follows: in the preliminaries section, we briefly introduce the forcings which we will use in the proof and produce a generic extension $W$ of $L$ which will be a well-suited ground model for our needs. We then start to prove the theorems from above. The main idea is to turn the problem of finding a partial order $\forceP$ which forces the $\Pi^1_3$-property into a fixed point problem.  We shall define a derivation operator which acts on a specific set of forcing iterations of length $\omega_1$. This operator will be applied transfinitely often, and will produce better and better approximations to the set of forcings we actually want to use in the end. The process is shown to converge in that eventually a fixed point $\mathcal{P}$, i.e. a suitable set of forcings is reached. Forcings which belong to this fixed point $\mathcal{P}$ allow for a certain, seemingly self-referential  definition of an $\omega_1$-length iteration which forces the $\Pi^1_3$-uniformization property over $W$.
We then follow up, to alter the said process such that it becomes applicable to the  canonical inner models $M_n$ with $n$ Woodin cardinals.


There are some similarities to \cite{H}, in particular the two proofs rely on a similar ground model $W$, which is a generic extension of $L$, and use a similar coding method which relies on a suitably chosen $\omega_1$-sequence of $\omega_1$-Suslin trees. However a more straightforward application of the ideas of \cite{H} will fail badly to produce a model of the $\Pi^1_3$-uniformization property. As a consequence, a solution has to necessarily introduce several new ideas in order to succeed. The presentation of those is the goal of this paper.
\section{Preliminaries}

\subsection{Notation}
The notation we use will be mostly standard, we hope. We write $\forceP=(\forceP_{\alpha} \, : \, \alpha < \gamma)$ for a forcing iteration of length $\gamma$ with initial segments $\forceP_{\alpha}$. The $\alpha$-th factor of the iteration will be denoted with $\forceP(\alpha)$. Note here that we drop the dot on $\forceP(\alpha)$, even though $\forceP(\alpha)$ is in fact a $\forceP_{\alpha}$-name of a partial order. 



\subsection{The forcings which are used}
The forcing notions used in the construction are standard. The preparatory extension is
obtained solely by adding branches through a fixed sequence of Suslin trees,
and the later coding uses the ground-model version of the $\omega_1$-Cohen
forcing together with Jensen--Solovay almost disjoint coding \cite{JensenSolovay}.

Recall that a Suslin tree is a normal tree of height $\omega_1$ with countable
levels and without an uncountable chain or antichain.  Whenever $S$ is a Suslin
tree, we also write $S$ for the forcing whose conditions are the nodes of $S$,
ordered by end-extension.  Thus forcing with $S$ adds a cofinal branch through
$S$.

\begin{definition}
Let $\vec{S}=\langle S_\xi\mid \xi<\kappa\rangle$ be a sequence of Suslin
trees.  We call $\vec{S}$ \emph{independent} if, for every nonempty finite set
$e\subseteq\kappa$, the product
\[
   \prod_{\xi\in e}S_\xi
\]
is again a Suslin tree.
\end{definition}

For $A\subseteq\kappa$ let
\[
   \operatorname{Br}_A(\vec S)=\prod_{\xi\in A}^{\mathrm{fs}}S_\xi
\]
denote the finite-support product of the corresponding branch forcings.  When
the sequence is fixed, we simply write $\operatorname{Br}_A$.

\begin{lemma}\label{finite-support-branch-products}
Let $\vec S=\langle S_\xi\mid\xi<\kappa\rangle$ be an independent sequence of
Suslin trees, where $\kappa\leq\omega_1$.
\begin{enumerate}
\item For every $A\subseteq\kappa$, the forcing $\operatorname{Br}_A$ is c.c.c.
\item If $A\subseteq B\subseteq\kappa$, then $\operatorname{Br}_A$ is a
      complete subforcing of $\operatorname{Br}_B$, and
      \[
         \operatorname{Br}_B\cong
         \operatorname{Br}_A\times\operatorname{Br}_{B\setminus A}.
      \]
\item If $\xi\notin A$, then
      \[
         \operatorname{Br}_A\Vdash ``S_\xi\text{ is a Suslin tree}."
      \]
\end{enumerate}
\end{lemma}

\begin{proof}
For (1), let $\{p_\alpha\mid\alpha<\omega_1\}$ be a family of conditions in
$\operatorname{Br}_A$.  After thinning it out, the supports form a
$\Delta$-system with finite root $r$, and the restrictions
$p_\alpha\upharpoonright r$ belong to the finite product
$\prod_{\xi\in r}S_\xi$.  By independence this finite product is Suslin, hence
c.c.c., so two of the root restrictions are compatible.  Since the remaining
parts of their supports are disjoint, the corresponding two conditions in
$\operatorname{Br}_A$ are compatible.

The factorization in (2) is the usual coordinatewise factorization of a
finite-support product.  For (3), observe that
\[
   \operatorname{Br}_A\times S_\xi
   \cong \operatorname{Br}_{A\cup\{\xi\}},
\]
which is c.c.c. by (1).  The standard product characterization of preservation
of a Suslin tree now yields the assertion.
\end{proof}

We next recall the almost disjoint coding forcing due to Jensen and Solovay \cite{JensenSolovay}.
We identify subsets of $\omega$ with their characteristic functions.  Let
$D=\langle d_\alpha\mid\alpha<\omega_1\rangle$ be an almost disjoint family of
infinite subsets of $\omega$, and let $X\subseteq\omega_1$.

\begin{definition}\label{definitionadcoding}
The almost disjoint coding forcing $\mathbb A_D(X)$ consists of pairs
$(r,R)\in[\omega]^{<\omega}\times[D]^{<\omega}$.  We put
$(s,S)\leq(r,R)$ if and only if
\begin{enumerate}
\item $r\subseteq s$ and $R\subseteq S$;
\item whenever $\alpha\in X$ and $d_\alpha\in R$, then
      $s\cap d_\alpha=r\cap d_\alpha$.
\end{enumerate}
\end{definition}

If $G\subseteq\mathbb A_D(X)$ is generic and
$x_G=\bigcup\{r:(r,R)\in G\}$, then
\[
   \alpha\in X\quad\Longleftrightarrow\quad
   x_G\cap d_\alpha\text{ is finite}.
\]
Throughout the paper, $D$ denotes the canonical $L$-definable almost disjoint
family obtained by recursively choosing the $<_L$-least infinite subset of
$\omega$ almost disjoint from all previously chosen members.

\begin{lemma}\label{a.d.coding preserves Suslin trees}
For every $X\subseteq\omega_1$, the forcing $\mathbb A_D(X)$ is Knaster.  In
particular it is c.c.c. and preserves every ground-model Suslin tree.
\end{lemma}

\begin{proof}
The Knaster property follows from the usual $\Delta$-system argument applied
to the finite side conditions.  If $T$ is a Suslin tree, then the product of a
Knaster forcing with the c.c.c. forcing $T$ is c.c.c.  Hence
$\mathbb A_D(X)\times T$ is c.c.c., which is equivalent to preservation of
$T$ as a Suslin tree.
\end{proof}

Finally, we write
\[
   \mathbb C_L(\omega_1)=({}^{<\omega_1}2)^L
\]
for the $L$-computed $\omega_1$-Cohen forcing, ordered by end-extension.  Thus
conditions are countable binary sequences belonging to $L$.  Whenever this
forcing is used in an outer model of $L$, the notation refers to this fixed
$L$-poset rather than to a recomputation in the ambient universe.

\subsection[The ground model W of the iteration]{The ground model $W$ of the iteration}
We now define the ground model over which the main iteration will be carried
out.  We start with the canonical independent Suslin sequence already present
in $L$ and add a cofinal generic branch through every member of that sequence.

The following lemma gives the simultaneous form of Jensen's diamond
construction which we use throughout the paper; see also \cite{Jech}.  We use
the levelwise product of finitely many trees.  This is a dense suborder of the
corresponding finite forcing product.

\begin{lemma}
\label{canonical-L-Suslin-sequence}
There is a canonical $L$-definable sequence
\[
   \vec S=\langle S_\xi\mid\xi<\omega_1\rangle
\]
of independent normal Suslin trees.  The construction is given by a recursion
which is uniform over $L_{\omega_1}$.
\end{lemma}



\begin{proof}
Work in $L$.  Fix the canonical $L$-diamond sequence
$\langle a_\alpha\mid\alpha<\omega_1\rangle$ and the canonical $L$-coding of
finite sets of ordinals, countable trees, and subsets of their finite products.
We recursively construct coherent countable approximations
$\vec S^{\,\alpha}=\langle S^\alpha_\xi\mid \xi<\alpha\rangle$ such that each
$S^\alpha_\xi$ is a normal splitting tree of height $\alpha$ and
$S^\beta_\xi=S^\alpha_\xi\restriction\beta$ whenever $\xi<\beta<\alpha$.  At
limit stages we take unions.

It remains to describe the passage from $\alpha$ to $\alpha+1$.  For old
coordinates $\xi<\alpha$, successor stages are handled by putting countably
many immediate successors above each node on the top level.  Suppose now that
$\alpha$ is a nonzero limit.  If $a_\alpha$ does not decode as a pair
$(e,A)$, where $e\in[\alpha]^{<\omega}\setminus\{\varnothing\}$ and $A$ is a
maximal antichain in the product
\[
   P_e^\alpha=\bigcup_{\delta<\alpha}
      \prod_{\xi\in e}(S^\alpha_\xi)_\delta,
\]
then we use the $<_L$-least countable family of cofinal branches through each
old tree which covers all old nodes, and put one new top node above each of
these branches.

If $a_\alpha$ does decode such a pair $(e,A)$, we choose the branch families
for the coordinates in $e$ simultaneously so as to seal $A$.  Concretely, by a
countable genericity argument in
$\prod_{(\xi,n)\in e\times\omega}^{\mathrm{fs}}S^\alpha_\xi$, choose cofinal
branches $b_{\xi,n}$, $\xi\in e$, $n<\omega$, such that they cover all old
nodes and such that, for every $\nu:e\to\omega$, the tuple of new top nodes
placed above $\langle b_{\xi,\nu(\xi)}\mid\xi\in e\rangle$ extends some member
of $A$.  The required dense sets are just the usual ones ensuring cofinality,
coverage of old nodes, and the displayed sealing requirement; density of the
last kind uses maximality of $A$ in $P_e^\alpha$.  For $\xi\notin e$ we again
use the default $<_L$-least covering family.  Finally, the new coordinate
$\alpha$ is initialized as the canonical finite-support binary tree of height
$\alpha+1$.  This completes the recursion.  Put
$S_\xi=\bigcup_{\xi<\alpha<\omega_1}S^\alpha_\xi$.  Then each $S_\xi$ is a
normal splitting tree of height $\omega_1$ with countable levels.

We show that the sequence is independent.  Fix a finite nonempty
$e\subseteq\omega_1$ and let
\[
   P_e=\bigcup_{\alpha<\omega_1}
      \prod_{\xi\in e}(S_\xi)_\alpha
\]
be the product.  This is dense in the usual finite product of the
corresponding tree forcings.  Let $A\subseteq P_e$ be a maximal antichain in
$L$.  There is a club of $\alpha>\sup(e)$ such that
$A\cap(P_e\restriction\alpha)$ is maximal in $P_e\restriction\alpha$, and,
by coherence of the coding, a club on which the fixed code $x_{e,A}$ for
$(e,A)$ restricts to a code for
$(e,A\cap(P_e\restriction\alpha))$.  By $\diamondsuit$, choose a limit
$\alpha$ in the intersection of these clubs with
$a_\alpha=x_{e,A}\cap\alpha$.  At stage $\alpha$ the construction seals
$A\cap(P_e\restriction\alpha)$.  Hence every node on level $\alpha$, and
therefore every later node, extends a member of this initial antichain.  Since
$A$ itself is an antichain, $A=A\cap(P_e\restriction\alpha)$, so $A$ is
countable.  Thus $P_e$ has no uncountable antichain.  A cofinal branch through
$P_e$ would, by normal splitting, yield an uncountable antichain by choosing at
successor levels nodes incompatible with the branch.  Hence $P_e$ is Suslin.
Since $e$ was arbitrary, $\vec S$ is independent.  The construction is
canonical from $<_L$, the diamond sequence, and the fixed coding, and is
therefore uniform over $L_{\omega_1}$.
\end{proof}

\begin{lemma}
\label{definabilityofvecS}
Let $M$ be a transitive uncountable model of
\[
   \ZFCminus+``\aleph_1\text{ exists}"
\]
such that
\[
   \omega_1^M=\omega_1^L=\omega_1.
\]
There is a $\Sigma_1$-formula
\[
   \vartheta_{\mathrm{Sus}}(t,\gamma,\eta,\omega_1),
\]
independent of $M$, which defines over $M$ the relation
\[
   \{(t,\gamma,\eta):
      \gamma,\eta<\omega_1\text{ and }t\in(S_\eta)_\gamma\}.
\]
Equivalently, there is a formula
$\operatorname{CanSuslin}_L(X,\omega_1)$, which may be taken to be
$\Sigma_1(\omega_1)$ over $H(\omega_2)$, such that in every
$\omega_1$-preserving outer model $U$ of $L$,
\[
   U\models\operatorname{CanSuslin}_L(X,\omega_1)
   \quad\Longleftrightarrow\quad X=\vec S.
\]
In particular, the canonical sequence $\vec S$ is computed in the same way in
all the outer models considered below.
\end{lemma}

\begin{proof}
By Lemma~\ref{canonical-L-Suslin-sequence}, the sequence \(\vec S\) is
obtained by a canonical recursion over \(L_{\omega_1}\), and the corresponding
level relation is \(\Sigma_1(\omega_1)\)-definable over \(L_{\omega_1}\).
If \(M\) is transitive, satisfies \(\ZFCminus+``\aleph_1\text{ exists}"\),
and has \(\omega_1^M=\omega_1^L\), then \(M\) computes \(L_{\omega_1}\)
correctly.  Hence \(M\) computes the same canonical diamond sequence, the same
coding apparatus, and therefore the same recursive construction of
\(\vec S\).  The same \(\Sigma_1\)-formula
\(\vartheta_{\mathrm{Sus}}(t,\gamma,\eta,\omega_1)\) consequently defines over
\(M\) the relation \(t\in(S_\eta)_\gamma\).

The formulation by \(\operatorname{CanSuslin}_L(X,\omega_1)\) is obtained by
saying that \(X\) is the union of the unique system produced by this
canonical recursion.  Since every \(\omega_1\)-preserving outer model of \(L\)
has the same \(L_{\omega_1}\), it satisfies this formula exactly for the true
sequence \(\vec S\).
\end{proof}

Let
\[
   \operatorname{Br}=\prod_{\xi<\omega_1}^{\mathrm{fs}}S_\xi
\]
be the finite-support branch product, computed in $L$, and let
$H\subseteq\operatorname{Br}$ be $L$-generic.  For every $\xi<\omega_1$ set
\[
   b_\xi=\bigcup\{p(\xi):p\in H\text{ and }\xi\in\operatorname{dom}(p)\}.
\]
Then $b_\xi$ is a cofinal $L$-generic branch through $S_\xi$.  In this paper,
to \emph{destroy} a Suslin tree always means to add such a cofinal generic
branch through it.  We define
\[
   W=L[H].
\]

\begin{theorem}\label{properties-of-W}
The model $W$ has the following properties.
\begin{enumerate}
\item $\operatorname{Br}$ is a c.c.c. forcing of size $\aleph_1$.  Hence $W$
      preserves all cardinals and cofinalities of $L$, and $W\models\CH$.
\item Every $S_\xi$ is destroyed in $W$, witnessed by the branch $b_\xi$.
\item For $A\subseteq\omega_1$ in $L$, let
      $H_A=H\cap\operatorname{Br}_A$.  Then $H_A$ is
      $\operatorname{Br}_A$-generic over $L$, and
      \[
         L[H_A]\models
         ``S_\xi\text{ has a cofinal branch}"
         \quad\Longleftrightarrow\quad \xi\in A.
      \]
\item Every real in $W$ belongs to $L[H_A]$ for some countable
      $A\subseteq\omega_1$ belonging to $L$.
\end{enumerate}
\end{theorem}

\begin{proof}
The first assertion follows from Lemma~\ref{finite-support-branch-products}.
Since $L\models\CH$ and $\operatorname{Br}$ has size $\aleph_1$, there are only
$\aleph_1$ many nice $\operatorname{Br}$-names for reals, and therefore $W$
still satisfies $\CH$.  The definition of $H$ and the dense sets requiring a
condition to reach a prescribed level of a prescribed coordinate give (2).

For (3), use the factorization
\[
   \operatorname{Br}\cong
   \operatorname{Br}_A\times\operatorname{Br}_{\omega_1\setminus A}.
\]
If $\xi\in A$, then $b_\xi\in L[H_A]$.  If $\xi\notin A$, then
Lemma~\ref{finite-support-branch-products} shows that
$\operatorname{Br}_A$ preserves $S_\xi$ as a Suslin tree, so no cofinal branch
through $S_\xi$ belongs to $L[H_A]$.

Finally, let $\dot x$ be a nice $\operatorname{Br}$-name for a real.  The
antichains occurring in $\dot x$ are countable, and the union of the finite
supports of all conditions appearing in them is therefore a countable set
$A\subseteq\omega_1$ in $L$.  Thus $\dot x$ is already a
$\operatorname{Br}_A$-name, proving (4).
\end{proof}

The relevant inner models of $W$ used below are obtained by retaining
selected branch coordinates from $H$.

\paragraph{Coding forcing}
Let $x\in W$ be a real and let $\eta<\omega_1$.  Write
$\mathbb C_L(\omega_1)_\eta$ for a tagged copy of
$\mathbb C_L(\omega_1)$, and let $\dot g_\eta$ be its canonical name for the
generic subset of $\omega_1$.  We now describe a canonical
$\mathbb C_L(\omega_1)_\eta$-name
$\dot Y(\check x,\dot g_\eta)$ for a subset of $\omega_1$.  The coding
forcing is the two-step iteration
\[
   \operatorname{Code}(x,\eta)
   :=\mathbb C_L(\omega_1)_\eta
      *\dot{\mathbb A}_D\bigl(\dot Y(\check x,\dot g_\eta)\bigr).
\]

For the definition of $Y$, let $g\subseteq\omega_1$ be
$\mathbb C_L(\omega_1)_\eta$-generic over $W$, and identify $g$ with its
characteristic function in ${}^{\omega_1}2$.  Let
\[
   \varrho_L:({}^{<\omega_1}2)^L\longrightarrow\omega_1
\]
be the canonical $L$-definable bijection.  We use $\varrho_L$ and $g$ to
form the set
$h\subseteq\omega_1$ of indices of the $\omega$-blocks of $\vec S$ on which
the characteristic function of $x$ will be written.  Put
\[
   h:=\{\varrho_L(g\restriction\alpha):\alpha<\omega_1\}
\]
and
\begin{align*}
   A:={}&\{\omega\gamma+2n:\gamma\in h,\ n\notin x\}\\
       &\mathbin{\cup}\{\omega\gamma+2n+1:\gamma\in h,\ n\in x\}.
\end{align*}

Let $X\subseteq\omega_1$ be chosen canonically so that it codes the following
objects:
\begin{itemize}
   \item the set $A\subseteq\omega_1$;
   \item the family $\langle b_\beta\mid\beta\in A\rangle$, where $b_\beta$
         is the branch through $S_\beta$ added by the $\beta$-th coordinate
         of $H$.  In particular, $b_\beta$ is $L$-generic for the forcing
         $S_\beta$.
\end{itemize}

If $L_{\theta}[X]$ is some uncountable $\ZFCminus$-model, and if we assume that $\gamma \in h$ then
 $L_{\theta}[X]$ can read off $x$ via looking at the $\omega$-block of $\vec{S}$-trees starting at $\gamma$ and determine which tree has an $\omega_1$-branch in $L_{\theta}[X]$. Thus  we say that $x$ is coded into $\vec{S}$ at the $\omega$-block starting at $\gamma$ in such a situation. To be more specific $L_{\theta}[X]$ satisfies the following formula:
\begin{itemize}
 \item[$(\ast)_{\gamma, x}$]  $n \in x$ if and only if $S_{\omega \cdot \gamma +2n+1}$ has an $\omega_1$-branch, and $n \notin x$ if and only if $S_{\omega \cdot \gamma +2n}$ has an $\omega_1$-branch.
\end{itemize}
Indeed if $n \notin x$ then we added a cofinal branch through $S_{\omega \cdot \gamma+ 2n}$. If on the other hand $S_{\omega \cdot\gamma +2n}$ does not have an $\omega_1$-branch in $L_{\theta}[X]$ then we must have added an $\omega_1$-branch through $S_{\omega \cdot \gamma +2n+1}$ as we always add an $\omega_1$-branch through either $S_{\omega \cdot \gamma +2n+1}$ or $S_{\omega \cdot \gamma +2n}$ and adding branches through some $S_{\alpha}$'s  will not affect that some $S_{\beta}$ remain Suslin in $L[X]$, as $\vec{S}$ is independent.



We now apply a localization argument inspired by David's trick\footnote{See \cite{David} for the original argument, where the strings in Jensen's coding machinery are altered so that certain unwanted universes are destroyed.  This destruction is emulated in our context below.} to this situation.  We reorganize the information coded by $X \subset \omega_1$ as a set $Y \subset \omega_1$ as follows.
It is clear that any transitive, $\aleph_1$-sized $\ZFCminus$ model $N$ of the form $N=L_{\theta}[X]$  will be able to first define $\vec{S}$ correctly  and also correctly decode out of $X$ all the information regarding $x$ being coded at each $\omega$-block of $\vec{S}$ starting at every $\gamma \in h$. 
Consequently, if we code the model $(N,\in)$ which contains $X$, as a set $X_N \subset \omega_1$, then for any uncountable $\beta$ such that $L_{\beta}[X_N] \models \ZFCminus$:
\[L_{\beta}[X_N] \models \text{\ldq The model decoded out of }X_N \text{ satisfies $(\ast)_{\gamma,x}$ for every $\gamma \in h$\rdq.} \]
In particular there will be an $\aleph_1$-sized ordinal $\beta$ as above and we can fix a club $C \subset \omega_1$ and a sequence $(M_{\alpha} \, : \, \alpha \in C)$ of countable elementary submodels  of $L_{\beta} [X_N]$ such that
\[\forall \alpha \in C (M_{\alpha} \prec L_{\beta}[X_N] \land M_{\alpha} \cap \omega_1 = \alpha)\]
Now let the set $Y\subset \omega_1$ code the pair $(C, X_N)$ such that the odd entries of $Y$ should code $X_N$ and if $E(Y)$ denotes  the set of even entries of $Y$ and $\{c_{\alpha} \, : \, \alpha < \omega_1\}$ is the enumeration of $C$ then
\begin{enumerate}
\item $E(Y) \cap \omega$ codes a well-ordering of type $c_0$.
\item $E(Y) \cap [\omega, c_0) = \emptyset$.
\item For all $\beta$, $E(Y) \cap [c_{\beta}, c_{\beta} + \omega)$ codes a well-ordering of type $c_{\beta+1}$.
\item For all $\beta$, $E(Y) \cap [c_{\beta}+\omega, c_{\beta+1})= \emptyset$.
\end{enumerate}
The following formula $\sigma(v_0)$ is now true for our real $x$:
\begin{itemize}
\item[$\sigma(x) \equiv$] For any countable transitive model $M$ of  ``$\ZFCminus$ and $\aleph_1$ exists$"$ such that $\omega_1^M=(\omega_1^L)^M$ and $ Y \cap \omega_1^M \in M$, $M$ can construct its version of the universe $L[Y \cap \omega_1^M]$, and the latter will see that there is an $\aleph_1^M$-sized transitive model $N \in L[Y \cap \omega_1^M]$ which models $(\ast)_{\gamma,x}$ for $\aleph_1^M$-many $\gamma \in M$.
\end{itemize}
Thus we have a local version of the property $(\ast)_{\gamma,x}$.

We have finally defined the desired set $Y$.  The second step of
$\operatorname{Code}(x,\eta)$ is
\[
   \mathbb A_D(Y),
\]
the Jensen--Solovay almost disjoint coding forcing from
Definition~\ref{definitionadcoding}, computed with the fixed canonical family
$D\in L$.  It codes the set $Y$ into one real $r$.  This forcing only depends
on the subset of $\omega_1$ we code, thus $\mathbb A_D(Y)$ will be independent
of the surrounding universe in which we define it, as long as it has the right
$\omega_1$ and contains the set $Y$.

The effect of the coding forcing $\operatorname{Code}(x,\eta)$ is that it generically adds a real $r$ such that
\begin{itemize}
\item[$ \Psi (r, x)\equiv $] For any countable, transitive model $M$ of ``$\ZFCminus$ and $\aleph_1$ exists$"$, such that $\omega_1^M=(\omega_1^L)^M$ and $ r  \in M$, $M$ can construct its version of $L[r]$, denoted by $L[r]^M$, which in turn thinks that there is a transitive $\ZFCminus$-model $N$ of size $\aleph_1^M$  such that $N$ believes $(\ast)_{\gamma,x}$ for an $\aleph_1^M$-sized set of ordinals $\gamma$.
\end{itemize}

Indeed, if $r$ and $M$ are as above, then $M$ and $L[r]^M$ compute the canonical almost disjoint family $D$ correctly below $\omega_1^M$.  As a consequence, $L[r]^M$ will contain the set $Y \cap \omega_1^M$, where $Y \subset \omega_1$ is as in $\sigma(x)$.  So in $L[Y \cap \omega_1^M]$, there is an $\aleph_1^M$-sized, transitive $N$ which models $(\ast)_{\gamma,x}$ for $\aleph_1^M$-many $\gamma  \in M$, as claimed. 

Note that $\Psi (r,x)$ is a $\Pi^1_2$-formula in the parameters $r$ and $x$, as the set $h \cap M \subset \omega_1^M$ is coded into $r$. We say in the above situation that $r$ witnesses that the real $x$ \emph{ is written into $\vec{S}$}, or that $x$ \emph{is coded into} $\vec{S}$.


To summarize, for every real $x$ and every $\eta<\omega_1$, the forcing $\operatorname{Code}(x,\eta)$ adds a real $r$ such that the $\Pi^1_2$-formula $\Psi(r,x)$ holds in the resulting extension of $W$.

We let $\Phi(x)$ denote the $\Sigma^1_3$-statement ``there is a real witnessing that $x$ is coded into $\vec{S}"$ that is
\[ \Phi(x) := \exists r \Psi (r,x). \] 
Via spelling out $\Psi(r,x)$ we see that $\Phi(x)$ is the following formula:
\begin{align*}
\Phi(x)\equiv \exists r\,\forall M\bigl(&M\text{ is countable and transitive, }M\models\ZFCminus,\\
&\omega_1^M=(\omega_1^L)^M\rightarrow M\models ``L[r]\text{ sees that }r\text{ codes}\\
&\text{a transitive }\aleph_1^M\text{-sized }\ZFCminus\text{-model }N'\\
&\text{such that }N'\text{ satisfies }(\ast)_{\gamma,x}
  \text{ for }\aleph_1^M\text{-many }\gamma"\bigr).
\end{align*}



 The projective and local statement $\Phi(x)$, if true,  will determine how certain inner models of the surrounding universe will look like with respect to branches through $\vec{S}$.
That is to say, if we assume that $\Psi(r,x) $ holds then $r$ also witnesses that the assertion of $\Psi(r,x)$ remains true even if we replace the  countable, transitive models $M$ in the statement of $\Psi (r,x)$ with any transitive  model $M$ of  the theory ``$\ZFCminus +$ $\aleph_1$ exists and $\aleph_1= \aleph_1^L"$. In other words  we can drop the assumption on the countability of $M$ in $\Psi (r,x)$, provided $\Psi(r,x)$ is true.
Indeed if we assume 
that there would be an uncountable, transitive $M$ which models $\ZFCminus$ and ``$
\aleph_1$ exists$"$ and $\omega_1= \omega_1^L$, which additionally contains $r$ and which witnesses that the conclusion of $\Psi(r,x)$ is false. Then by L\"owenheim-Skolem, there would be a countable $N\prec M$, $r\in N$ which we can transitively collapse to obtain the transitive $\bar{N}$. But $\bar{N}$ would witness that the conclusion of $\Psi(r,x)$ is not true for every countable, transitive model, which is a contradiction to our assumption that $\Psi(r,x)$ is true. 

Consequently, the real $r$ carries enough information that
the universe $L[r]$ will see that certain trees from $\vec{S}$ have branches in that $L[r]$ will decode, with the help of the almost disjoint family $D \in L$ a set $h \subset \omega_1$ such that
\begin{align*}
n \in x \Rightarrow L[r] \models  \forall \gamma \in h \, (``S_{\omega \gamma + 2n+1} \text{ has an $\omega_1$-branch}").
\end{align*}
and
\begin{align*}
n \notin x \Rightarrow L[r] \models \forall \gamma \in h \, ( ``S_{\omega \gamma + 2n} \text{ has an $\omega_1$-branch}").
\end{align*}
Indeed, the universe $L[r]$ will see that there is a transitive model $N$ of ``$\ZFCminus+$ $\aleph_1$ exists and $\aleph_1=\aleph_1^L"$ which believes $(\ast)_{\gamma,x}$ for every $\gamma \in h \subset \omega_1$, the latter being coded into $r$. But by upwards $\Sigma_1$-absoluteness, and the fact that $N$ can compute $\vec{S}$ correctly, if $N$ thinks that some tree in $\vec{S}$ has a branch, then $L[r]$ must think so as well. 


The definition of the forcing $\operatorname{Code}(x,\eta)$ exhibits a significant degree of absoluteness.  Specifically, its definition is independent of the universe in which it is computed, provided that the universe has the correct $\omega_1$ and contains the real $x$.  We will leverage this property shortly.

We do not iterate the full forcings $\operatorname{Code}(x,\eta)$. Instead, we separate their two components. We first add one global $\omega_1$-Cohen reservoir and then, over the resulting extension, form a finite-support iteration of the almost disjoint coding forcings.


\subsection{Allowable forcings}
\label{subsection-allowable-forcings}

The forcing $\operatorname{Code}(u,\eta)$ described above has two components. Its $\omega_1$-Cohen part chooses a coding area, while its Jensen--Solovay part writes the resulting set $Y\subseteq\omega_1$ into a real. In the applications below, $u$ will be the fixed recursive real code of a triple $(x,y,m)$. We write
\[
   \dot Y(\dot x,\dot y,\dot m,\dot g)
\]
for the canonical name of the set $Y$ obtained by applying the preceding coding construction to the recursive code of $(\dot x,\dot y,\dot m)$ and the Cohen generic $\dot g$.

In an allowable forcing the two components of the coding forcing are organized separately. Put
\[
   \forceC^{\mathrm{area}}
   :=\prod_{\eta<\omega_1}^{\mathrm{cs}}\mathbb C_L(\omega_1).
\]
After passing to the standard dense partial-function presentations and fixing a bijection between $\omega_1$ and $\omega_1\times\omega_1$, the forcing $\forceC^{\mathrm{area}}$ is isomorphic to one copy of $\mathbb C_L(\omega_1)$. We retain the product presentation because its coordinate generics
\[
   \langle\dot g_\eta\mid\eta<\omega_1\rangle
\]
supply the coding areas used by the later iteration. For every $\eta<\omega_1$, the generic $\dot g_\eta$ determines the area
\[
   \dot h_\eta
   :=\{\varrho_L(\dot g_\eta\restriction\xi):\xi<\omega_1\}.
\]

As computed in $L$, the forcing $\forceC^{\mathrm{area}}$ is $\sigma$-closed. Since $W=L[H]$ is obtained from $L$ by the c.c.c. forcing $\operatorname{Br}$, Easton's lemma implies that
\[
   W\models
   ``\forceC^{\mathrm{area}}\text{ is }\omega\text{-distributive}."
\]
See \cite[Chapter~15]{Jech}. Consequently, $\forceC^{\mathrm{area}}$ adds no reals over $W$. The same remains true after permuting its coordinates. Thus the product presentation provides $\omega_1$ many coding areas without adding reals at the reservoir step.

Over the reservoir extension we form a finite-support iteration of almost disjoint coding forcings. Schematically, at a stage $\beta$, the bookkeeping supplies names
\[
   \dot x_\beta,\quad
   \dot y_\beta,\quad
   \dot m_\beta.
\]
Independently of the generic filter, we fix in $L$ an injective assignment
\[
   \vec\eta=\langle\eta_\beta\mid\beta<\delta\rangle
\]
of reservoir coordinates to the stages of the iteration.  For every
$\beta\leq\delta$, put
\[
   U_\beta:=\operatorname{ran}(\vec\eta\upharpoonright\beta).
\]
The bookkeeping data used at stage $\beta$ will be represented by names whose
reservoir support is contained in $U_\beta$.  Since $\vec\eta\in L$, the same
assignment is available in every branch inner model of $W$.  The restriction
on the names does not remove any possible real value: after the c.c.c. forcing
which produces $W$ and the c.c.c. almost disjoint tail below $\beta$, the
unused part of the $L$-computed $\sigma$-closed reservoir is
$\omega$-distributive by Easton's lemma.  The next factor is
\[
   \dot{\forceA}_\beta
   :=\mathbb A_D\bigl(
      \dot Y(\dot x_\beta,\dot y_\beta,\dot m_\beta,
             \dot g_{\eta_\beta})
   \bigr).
\]
If
\[
   \dot{\forceR}_\delta
   =\Asterisk_{\beta<\delta}^{\mathrm{fs}}
      \dot{\forceA}_\beta,
\]
then the full forcing is
\[
   \forceP_\delta
   =\forceC^{\mathrm{area}}*\dot{\forceR}_\delta.
\]
Thus the Cohen reservoir is added only once. After this first step, only the almost disjoint coding forcings are iterated, with finite support.

Let $K*G$ be $\forceP_\delta$-generic over $W$. For $\beta<\delta$, put
\[
   u_\beta
   :=\bigl\langle
      \dot x_\beta^{K*G},
      \dot y_\beta^{K*G},
      \dot m_\beta^{K*G}
   \bigr\rangle.
\]
We call $\{u_\beta:\beta<\delta\}$ the set of reals intentionally coded by the iteration. It is essential that this iterative process does not inadvertently code reals which were not explicitly intended. Lemma~\ref{no-unwanted-allowable-codes} will show that, in the final extension,
\[
   \Phi(u)
   \quad\Longleftrightarrow\quad
   u=u_\beta\text{ for some }\beta<\delta.
\]

\begin{definition}\label{definition-allowable-iteration}
Let $\delta<\omega_1$, let
\[
   F=\langle
      (\dot x_\beta,\dot y_\beta,\dot m_\beta)
      \mid\beta<\delta
   \rangle\in W
\]
be a bookkeeping function, and let
\[
   \vec\eta=\langle\eta_\beta\mid\beta<\delta\rangle\in L
\]
be an injective sequence of ordinals below $\omega_1$.  For
$\beta\leq\delta$, put
\[
   U_\beta:=\operatorname{ran}(\vec\eta\upharpoonright\beta).
\]
An \emph{allowable iteration over the reservoir of length $\delta$ relative
to $(F,\vec\eta)$} is a sequence
\[
   \forceP=\langle\forceP_\beta\mid\beta\leq\delta\rangle,
   \qquad
   \forceP_\beta=\forceC^{\mathrm{area}}*\dot{\forceR}_\beta,
\]
defined recursively as follows.
\begin{enumerate}
   \item $\dot{\forceR}_0$ is the trivial forcing.  Thus
   \[
      \forceP_0=\forceC^{\mathrm{area}}.
   \]

   \item Suppose that $\dot{\forceR}_\beta$ has been defined.  By induction,
   $\dot{\forceR}_\beta$ is the canonical lift of a
   $\forceC^{\mathrm{area}}_{U_\beta}$-name.  Put
   \[
      \forceP_\beta^{U_\beta}
      :=\forceC^{\mathrm{area}}_{U_\beta}*\dot{\forceR}_\beta.
   \]
   Then $\forceP_\beta^{U_\beta}$ is a complete subforcing of
   $\forceP_\beta$.  The entries of $F(\beta)$ are required to be canonical
   lifts of $\forceP_\beta^{U_\beta}$-names such that
   \[
      \forceP_\beta^{U_\beta}\Vdash
      ``\dot x_\beta,\dot y_\beta\in2^\omega
      \text{ and }\dot m_\beta\in\omega."
   \]
   Earlier names are identified with their canonical lifts to
   $\forceP_\beta$.  Since $\eta_\beta\notin U_\beta$, the coordinate used
   for the new coding area does not occur in the data which determine the
   real coded at this stage.  Let
   \[
      \dot Y_\beta
      :=\dot Y(\dot x_\beta,\dot y_\beta,\dot m_\beta,
               \dot g_{\eta_\beta})
      \quad\text{and}\quad
      \dot{\forceA}_\beta:=\mathbb A_D(\dot Y_\beta).
   \]
   Then
   \[
      \dot{\forceR}_{\beta+1}
      :=\dot{\forceR}_\beta*\dot{\forceA}_\beta.
   \]

   \item If $\lambda\leq\delta$ is a limit ordinal, then
   $\dot{\forceR}_\lambda$ is the finite-support limit of
   $\langle\dot{\forceR}_\beta\mid\beta<\lambda\rangle$.
\end{enumerate}
The final forcing of the iteration is $\forceP_\delta$.  A forcing is
\emph{allowable} if it is isomorphic to the final forcing of an allowable
iteration over the reservoir.

Suppose that $\forceP$ and $\forceQ$ are allowable iterations over the
reservoir of lengths $\delta\leq\varepsilon$, relative to
$(F^\forceP,\vec\eta^{\forceP})$ and
$(F^\forceQ,\vec\eta^{\forceQ})$, respectively.  We say that $\forceQ$ is an
\emph{allowable extension} of $\forceP$, and write
\[
   \forceQ\triangleright\forceP,
\]
if
\[
   F^\forceQ\upharpoonright\delta=F^\forceP,
   \qquad
   \vec\eta^{\forceQ}\upharpoonright\delta=\vec\eta^{\forceP},
\]
and the iteration
$\langle\dot{\forceR}^{\forceQ}_\beta\mid\beta\leq\varepsilon\rangle$
end-extends
$\langle\dot{\forceR}^{\forceP}_\beta\mid\beta\leq\delta\rangle$.
\end{definition}

Thus an allowable forcing has the schematic form
\[
   \left(\prod_{\eta<\omega_1}^{\mathrm{cs}}
      \mathbb C_L(\omega_1)\right)
   *
   \left(\Asterisk_{\beta<\delta}^{\mathrm{fs}}
      \mathbb A_D(\dot Y_\beta)\right).
\]
A condition may be represented by a pair $(c,\dot p)$, where the product
support of $c$ is countable and the iteration support of $\dot p$ is finite.
The first coordinate is the reservoir of coding areas; only the almost
disjoint forcings are iterated.  The injectivity of
$\vec\eta=\langle\eta_\beta\mid\beta<\delta\rangle$ ensures that each coding
area is used at most once.  The requirement $\vec\eta\in L$ ensures that every
branch inner model used in the localization arguments computes the same area
assignment.  At stage $\beta$, all bookkeeping names use only coordinates in
$U_\beta$, while $\eta_\beta\notin U_\beta$ supplies the next coding area.

We identify a triple $(x,y,m)$ with its fixed recursive real code.  For a real
$u$, $\gamma<\omega_1$, and $n<\omega$, put
\[
   \tau_u(\gamma,n)=
   \begin{cases}
      \omega\gamma+2n+1,& n\in u,\\
      \omega\gamma+2n,& n\notin u,
   \end{cases}
\]
and, for $h\subseteq\omega_1$, let
\[
   \operatorname{Sel}(h,u)
   :=\{\tau_u(\gamma,n):\gamma\in h,\ n<\omega\}.
\]
Thus a stage which codes $u$ on the area $h$ uses precisely the branches
$\langle b_\xi\mid\xi\in\operatorname{Sel}(h,u)\rangle$.


\begin{lemma}
\label{basic-properties-allowable-forcings}
Let $\forceP_\delta=\forceC^{\mathrm{area}}*\dot{\forceR}_\delta$ be an
allowable forcing of length $\delta<\omega_1$.
\begin{enumerate}
   \item For every $\beta<\delta$,
   \[
      \forceP_\beta\Vdash
      ``\dot{\forceA}_\beta\text{ is Knaster and }
      |\dot{\forceA}_\beta|=\aleph_1."
   \]
   Consequently,
   $\forceC^{\mathrm{area}}\Vdash
   ``\dot{\forceR}_\delta\text{ is c.c.c. and has size at most }
   \aleph_1."$

   \item $\forceP_\delta$ preserves $\omega_1$ and $\CH$.

   \item If $\forceP^0$ and $\forceP^1$ are allowable, then
   \[
      \forceP^0\times\forceP^1
   \]
   is isomorphic to an allowable forcing.  More precisely, after an
   $L$-definable reindexing of the two reservoir factors onto disjoint
   sets of coordinates, the product admits an allowable iteration over the
   global reservoir.
\end{enumerate}
\end{lemma}

\begin{proof}
Clause (1) follows from Lemma~\ref{a.d.coding preserves Suslin trees} and the
definition of $\mathbb A_D(Y)$.  Indeed, $D$ has size $\aleph_1$, so every
almost disjoint coding forcing has size $\aleph_1$ and is Knaster.  A
finite-support iteration of c.c.c. forcings is c.c.c.; since
$\delta<\omega_1$ and $\CH$ holds after the reservoir forcing, the iteration
has size at most $\aleph_1$.

For (2), note first that $\forceC^{\mathrm{area}}$, as computed in $L$, is
$\sigma$-closed and has size $\aleph_1$.  Since $W=L[H]$ is obtained from $L$
by the c.c.c. branch product, Easton's lemma yields
\[
   W\models
   ``\forceC^{\mathrm{area}}\text{ is }\omega\text{-distributive}."
\]
Hence the reservoir forcing preserves $\omega_1$, adds no reals, and preserves
$\CH$.  In the reservoir extension, $\dot{\forceR}_\delta$ is a c.c.c.
forcing of size at most $\aleph_1$ by (1).  It therefore preserves
$\omega_1$, and, since the ground model of this second step satisfies $\CH$,
it adds at most $\aleph_1$ many reals.  Thus the final extension still
satisfies $\CH$.

For (3), write
\[
   \forceP^i
   =\forceC_i^{\mathrm{area}}*\dot{\forceR}^{i}_{\delta_i}
   \qquad(i<2).
\]
Fix in $L$ a partition $\omega_1=E_0\mathbin{\dot\cup}E_1$ and bijections
$e_i:\omega_1\longrightarrow E_i$.  Since the reservoirs are countable-support
products of the same $L$-computed forcing,
\[
   \forceC_0^{\mathrm{area}}\times\forceC_1^{\mathrm{area}}
   \cong\forceC^{\mathrm{area}}
\]
via the reindexing $e_0,e_1$.  Under this isomorphism the coding areas used by
the two factors are moved to disjoint sets of reservoir coordinates.

Over this common reservoir, the product
\[
   \dot{\forceR}^{0}_{\delta_0}	\times
   \dot{\forceR}^{1}_{\delta_1}
\]
of the two finite-support almost-disjoint coding iterations is isomorphic to
their concatenation: first use the translated factors from
$\dot{\forceR}^{0}_{\delta_0}$ and then the translated factors from
$\dot{\forceR}^{1}_{\delta_1}$.  
\end{proof}

\begin{corollary}\label{allowable-forcings-preserve-ch}
Every allowable forcing preserves $\CH$.  More generally, suppose that
\[
   \forceP_{\omega_1}
   =\forceC^{\mathrm{area}}*\dot{\forceR}_{\omega_1}
\]
is an $\omega_1$-length allowable iteration over the reservoir whose restriction to every
$\alpha<\omega_1$ is allowable.  Then $\forceP_{\omega_1}$ preserves
$\omega_1$ and $\CH$.
\end{corollary}

\begin{proof}
Only the second assertion requires comment.  In the reservoir extension,
$\dot{\forceR}_{\omega_1}$ is the finite-support iteration of the
$\omega_1$ many c.c.c. factors appearing in the iteration.  It is
c.c.c. and, by $\CH$, has size $\aleph_1$.  The argument from
Lemma~\ref{basic-properties-allowable-forcings}(2) now applies verbatim.
\end{proof}

Let $K*G$ be $\forceP_\delta$-generic over $W$.  For $\beta<\delta$ set
\[
   u_\beta
   :=\bigl\langle
      \dot x_\beta^{K*G},
      \dot y_\beta^{K*G},
      \dot m_\beta^{K*G}
   \bigr\rangle
   \quad\text{and}\quad
   h_\beta:=\dot h_{\eta_\beta}^{K}.
\]
We call $\{u_\beta:\beta<\delta\}$ the set of reals intentionally coded by
the iteration.

\begin{lemma}
\label{allowable-localization-lemma}
Let
\[
   \forceP_\delta=\forceC^{\mathrm{area}}*\dot{\forceR}_\delta
\]
be an allowable forcing over $W=L[H]$, and let $K*G$ be
$\forceP_\delta$-generic over $W$.  If $x\in W[K][G]\cap2^\omega$,
then there are countable sets
\[
   E_x\subseteq\omega_1,\qquad I_x\subseteq\delta,
   \qquad B_x\subseteq\omega_1
\]
such that, putting
\[
   \operatorname{Sel}_{I_x}
   :=\bigcup_{\beta\in I_x}\operatorname{Sel}(h_\beta,u_\beta),
\]
we have
\[
   x\in
   L[H_{B_x\cup\operatorname{Sel}_{I_x}}]
    [K\upharpoonright E_x]
    [G\upharpoonright I_x].
\]
Here $G\upharpoonright I_x$ denotes the generic sequence for the induced
closed subiteration on $I_x$.  Since $\vec\eta\in L$, the same area assignment
belongs to every model $L[H_B]$ occurring in the localization argument; no
branch support is needed to recover it.  The set $I_x$ may be chosen closed
under the bookkeeping names of the iteration, and $E_x$ may be chosen so that
\[
   \{\eta_\beta:\beta\in I_x\}\subseteq E_x.
\]
Thus, whenever a stage belongs to $I_x$, the reservoir coordinate assigned to
that stage also belongs to the localized reservoir support.
\end{lemma}

\begin{proof}
We prove the assertion by induction on the length $\delta$ of the allowable
iteration.  The statement is clear at length $0$: the reservoir forcing adds
no reals over $W$, and every real of $W=L[H]$ belongs to $L[H_B]$ for some
countable $B\subseteq\omega_1$ by Theorem~\ref{properties-of-W}.

Assume the lemma is known for all shorter allowable iterations.  Let
$x\in W[K][G]\cap2^\omega$, and fix a nice $\forceP_\delta$-name $\dot x$ for
$x$.  A condition in $\forceP_\delta$ has the form
$(c,\dot p)$, where $c$ is a condition in the reservoir product and
$\dot p$ is a finite-support iteration condition.  Let
\[
   I_0=\bigcup\{\operatorname{supp}_{\mathrm{it}}(p):
        p\text{ occurs in }\dot x\}
\]
and let $E_0$ be the union of the reservoir supports of the conditions
occurring in $\dot x$.  Since $\dot x$ is a nice name for a real,
$I_0\subseteq\delta$ and $E_0\subseteq\omega_1$ are countable.

We close $I_0$ under the names which define the iterands.  Thus, whenever a
stage $\beta$ is retained, we also retain nice names for
\[
   \dot x_\beta,\quad \dot y_\beta,
   \quad \dot m_\beta,
\]
and all iteration coordinates and reservoir coordinates occurring in those
names.  We also add $\eta_\beta$ to the reservoir support.  Repeating this
operation only countably many times gives countable sets
$I\subseteq\delta$ and $E\subseteq\omega_1$ such that
\[
   \{\eta_\beta:\beta\in I\}\subseteq E.
\]
By construction, the name $\dot x$ and all names needed to define the
iterands at coordinates in $I$ are read from the induced closed subiteration
on $I$ together with the reservoir coordinates in $E$.

For each $\beta\in I$, the real $u_\beta$ belongs to the extension by the
initial iteration of length $\beta$.  By the induction hypothesis, applied
to that shorter iteration, there are countable sets
$B_\beta\subseteq\omega_1$, $E_\beta\subseteq\omega_1$, and
$I_\beta\subseteq\beta$ such that
\[
   u_\beta\in
   L[H_{B_\beta\cup\operatorname{Sel}_{I_\beta}}]
    [K\upharpoonright E_\beta]
    [G\upharpoonright I_\beta].
\]
Enlarge $I$ and $E$ by all $I_\beta$ and $E_\beta$ obtained in this way, and
iterate the closure.  Since only countably many stages are involved, the final
sets, still denoted by $I$ and $E$, are countable.

Let
\[
   B=\bigcup_{\beta\in I} B_\beta\cup B_0,
\]
where $B_0$ is a countable branch support for the ground-model parameters
which occur in the closed system of names.  Again
Theorem~\ref{properties-of-W} gives such a $B_0$, and $B$ is countable.  Put
\[
   \operatorname{Sel}_I=
   \bigcup_{\beta\in I}\operatorname{Sel}(h_\beta,u_\beta).
\]
For every $\beta\in I$, the set $Y_\beta$ used by the $\beta$-th almost
disjoint coding factor is computed from the localized data for $u_\beta$, the
reservoir coordinate $g_{\eta_\beta}$, and precisely the branches with indices
in $\operatorname{Sel}(h_\beta,u_\beta)$.  Therefore the whole closed system of
names, and in particular $\dot x$, is interpreted in
\[
   L[H_{B\cup\operatorname{Sel}_I}]
    [K\upharpoonright E]
    [G\upharpoonright I].
\]
Hence $x$ belongs to this model.  This proves the induction step and the
lemma.
\end{proof}

\begin{lemma}\label{no-unwanted-allowable-codes}
Let $\forceP_\delta$ be allowable, and let $K*G$ be generic over $W$.  Then
\[
   W[K][G]\models
   \Phi(u)
   \quad\Longleftrightarrow\quad
   u=u_\beta\text{ for some }\beta<\delta.
\]
The same conclusion holds for an $\omega_1$-length iteration over the reservoir
all of whose proper initial segments are allowable.
\end{lemma}

\begin{proof}
If $u=u_\beta$, the real added by $\mathbb A_D(\dot Y_\beta)$ witnesses
$\Psi(r,u)$ immediately after stage $\beta$.  Since $\Psi$ is $\Pi^1_2$, this
remains true in every later forcing extension by Shoenfield absoluteness.
Hence every intentionally coded real satisfies $\Phi$.

Conversely, work in $W[K][G]$ and suppose that $r$ witnesses $\Phi(u)$.  By
the coding analysis of the preceding subsection, $r$ decodes an unbounded
coding area $h\subseteq\omega_1$, of the form
\[
   h=\{\varrho_L(a\restriction\xi):\xi<\omega_1\}
\]
for a function $a\in{}^{\omega_1}2$ all of whose proper initial segments
belong to $L$, and
\[
   L[r]\models
   ``S_{\tau_u(\gamma,n)}\text{ has an }\omega_1\text{-branch}"
   \qquad(\gamma\in h,\ n<\omega).
\]

We use the following elementary separation fact.  If $a\neq a'$, then the
corresponding areas $h(a)$ and $h(a')$ have bounded intersection.  Indeed,
after the first coordinate at which $a$ and $a'$ differ, no sufficiently long
initial segment of one equals an initial segment of the other; injectivity of
$\varrho_L$ leaves only countably many common values.  Hence the areas
$\langle h_\beta\mid\beta<\delta\rangle$ are pairwise almost disjoint.

Assume towards a contradiction that $u\neq u_\beta$ for all $\beta<\delta$.
Apply Lemma~\ref{allowable-localization-lemma} to the two reals $r$ and $u$,
and take the union of the resulting supports.  Thus there are countable sets
\[
   E\subseteq\omega_1,
   \qquad I\subseteq\delta,
   \qquad B\subseteq\omega_1
\]
such that, with
\[
   \operatorname{Sel}_I=
      \bigcup_{\beta\in I}\operatorname{Sel}(h_\beta,u_\beta),
\]
we have
\[
   r,u\in
   M:=L[H_{B\cup\operatorname{Sel}_I}]
      [K\upharpoonright E]
      [G\upharpoonright I].
\]
We may assume that $I$ is closed under the bookkeeping names of the allowable
iteration.

We now choose a tree coordinate required by the alleged code for $u$ but not
present in the localized branch support.  Put
\[
   B^{\mathrm{blk}}
   :=\{\gamma<\omega_1:
       [\omega\gamma,\omega\gamma+\omega)\cap B\neq\emptyset\}.
\]
This set is countable.

First suppose that $h\neq h_\beta$ for every $\beta\in I$.  Since $I$ is
countable, choose
\[
   \gamma\in h\setminus
      \left(B^{\mathrm{blk}}\cup
      \bigcup_{\beta\in I}(h\cap h_\beta)\right).
\]
Then for any $n<\omega$,
\[
   \tau_u(\gamma,n)\notin
   B\cup\operatorname{Sel}_I.
\]
Fix $n=0$ and set $\xi=\tau_u(\gamma,0)$.

Now suppose that $h=h_{\beta_0}$ for some $\beta_0\in I$.  Since
$u\neq u_{\beta_0}$, choose $n<\omega$ with
$u(n)\neq u_{\beta_0}(n)$.  Choose
\[
   \gamma\in h_{\beta_0}\setminus
      \left(B^{\mathrm{blk}}\cup
      \bigcup_{\beta\in I\setminus\{\beta_0\}}
         (h_{\beta_0}\cap h_\beta)
      \right),
\]
and put $\xi=\tau_u(\gamma,n)$.  At the $n$-th pair of the $\gamma$-block, the
stage $\beta_0$ selects the opposite tree, and no other stage in $I$ uses this
block.  Hence again
\[
   \xi\notin B\cup\operatorname{Sel}_I.
\]

In either case, $\xi=\tau_u(\gamma,n)$ for some $\gamma\in h$ and $n<\omega$,
while $\xi\notin B\cup\operatorname{Sel}_I$.  By
Lemma~\ref{finite-support-branch-products}, the tree $S_\xi$ remains Suslin in
$L[H_{B\cup\operatorname{Sel}_I}]$.  The reservoir part
$K\upharpoonright E$ is an $L$-computed product of $\omega_1$-Cohen forcings,
and the almost disjoint part $G\upharpoonright I$ is a finite-support
iteration of Knaster forcings whose defining sets use only branches in
$B\cup\operatorname{Sel}_I$.  Therefore the preservation facts from
Lemma~\ref{a.d.coding preserves Suslin trees}, together with the preservation
of unused Suslin trees by the reservoir forcing, give
\[
   M\models ``S_\xi\text{ is Suslin}."
\]

But $r\in M$, and the displayed decoding property of $r$ gives
\[
   L[r]\models
   ``S_\xi\text{ has an }\omega_1\text{-branch}."
\]
Since $L[r]\subseteq M$, this branch belongs to $M$, contradicting that
$M\models``S_\xi$ is Suslin".  Hence $u=u_\beta$ for some $\beta<\delta$.

Finally, let the iteration have length $\omega_1$, and suppose that
$\Phi(u)$ holds in the final extension, witnessed by $r$.  Nice names for $u$
and $r$ involve only countably many conditions, and each condition has finite
iteration support.  Hence these names are supported by some
$\alpha<\omega_1$.  Since $\Psi(r,u)$ is $\Pi^1_2$, Shoenfield absoluteness
implies that it already holds after the $\alpha$-th initial segment.  Applying
the first part of the lemma to that allowable initial segment shows that $u$
was intentionally coded before $\alpha$.
\end{proof}




\subsection{Motivation for the allowability hierarchy}
\label{subsection-motivation-allowability}

Given the notion of allowable forcing, we now form successive approximations to the class of forcings which we eventually want to use.  We first describe the idea leading to these approximations.

Fix a recursive enumeration
\[
   \langle \varphi_m(v_0,v_1)\mid m\in\omega\rangle
\]
of the $\Pi^1_3$-formulas with two free real variables, and put
\[
   A_m:=\{(x,y)\in(2^\omega)^2:\varphi_m(x,y)\}.
\]
We write $f_m$ for the uniformizing function which we intend to define for $A_m$, and write $f(m,x)$ for $f_m(x)$.  The goal is to choose, for every $m\in\omega$ and every real $x$ for which
\[
   A_{m,x}:=\{y\in2^\omega:(x,y)\in A_m\}
\]
is nonempty, a real $f(m,x)$ such that
\[
   (x,f(m,x))\in A_m,
\]
and to code every triple $(x,y,m)$ such that
\[
   y\neq f(m,x)\quad\text{and}\quad(x,y)\in A_m
\]
into $\vec S$.  The final forcing will have length $\omega_1$, and each of its countable initial segments will be allowable.  Since the statement
\[
   ``(x,y,m)\text{ is coded into }\vec S"
\]
is $\Sigma^1_3$, the requirement
\[
   (x,y)\in A_m
   \quad\text{and}\quad
   (x,y,m)\text{ is not coded into }\vec S
\]
is $\Pi^1_3$.  The bookkeeping of the final iteration will ensure that, whenever $A_{m,x}\neq\varnothing$, exactly one such $y$ remains.

The idea behind the definition can be seen in the following toy example.  Work in $W$, let $x\in W$, and suppose that
\[
   A_{m,x}=\{y_0,y_1\}.
\]
Assume in addition that no allowable forcing over $W$ adds a new member to $A_{m,x}$.  We consider only the problem of defining $f(m,x)$ and omit, for the moment, the interaction with the other sections.

Call $y\in A_{m,x}$ \emph{eliminable} if there is an allowable forcing $\forceP$ such that
\[
   \forceP\Vdash (x,y)\notin A_m.
\]
Call $y$ \emph{persistent} if it is not eliminable.  Thus, if $y_0$ is persistent, we may reserve $y_0$ as the value of $f(m,x)$ and code $(x,y_1,m)$ into $\vec S$.  If $y_0$ is eliminable, witnessed by an allowable forcing $\forceP$, we consider $y_1$.  If $y_1$ is persistent, we reserve $y_1$ and code $(x,y_0,m)$.

It remains to consider the case in which both values are eliminable.  Let $\forceP$ and $\forceQ$ witness, respectively, that
\[
   \forceP\Vdash (x,y_0)\notin A_m
   \qquad\text{and}\qquad
   \forceQ\Vdash (x,y_1)\notin A_m.
\]
By Lemma~\ref{basic-properties-allowable-forcings}, the product
\[
   \forceP\times\forceQ
\]
is isomorphic to an allowable forcing.  Identifying the product with its
normalized allowable presentation, upward absoluteness gives
\[
   \forceP\times\forceQ\Vdash
   (x,y_0)\notin A_m\wedge(x,y_1)\notin A_m.
\]
By the assumption that no new member is added to the section, the product forces
\[
   A_{m,x}=\varnothing.
\]
In this case there is no value to uniformize at $x$.

This argument isolates the two operations which are needed later:
\begin{equation}
\label{two-operations-uniformization}
\begin{split}
   &\text{reserve a value which is persistent under the forcings under consideration;}\\
   &\text{combine eliminating forcings by taking products.}
\end{split}
\end{equation}
The difficulty is that an eliminating forcing may add codes which interfere with other sections.  Thus it is not enough to use arbitrary allowable forcings as witnesses.  We need a class $\mathcal P$ of allowable forcings such that the following rule produces forcings in $\mathcal P$ again.  For every task $(x,m)$ considered by the bookkeeping, either
\[
   \exists y\in A_{m,x}\ \forall\forceQ\in\mathcal P
   \bigl(\forceQ\text{ extending the current forcing}\rightarrow
   \forceQ\Vdash(x,y)\in A_m\bigr),
\]
in which case one such $y$ is reserved and the other values are coded, or every value under consideration can be eliminated by a forcing in $\mathcal P$.

This turns the construction into a fixed point problem.  Let
\[
   \mathcal A^W_0
\]
be the class of allowable iterations.  The first refinement searches for values which are persistent under $\mathcal A^W_0$-extensions; the resulting iterations form $\mathcal A^W_1$.  The next refinement searches also for values which are persistent under $\mathcal A^W_1$-extensions, and so on.  We obtain a decreasing hierarchy
\[
   \mathcal A^W_0\supseteq\mathcal A^W_1\supseteq\cdots
   \supseteq\mathcal A^W_\alpha\supseteq\cdots.
\]
The case $\alpha=1$ is the first approximation described above. 

The reason for continuing the construction is the usual moving-target problem.  In the definition of $1$-allowability we test persistence under $0$-allowable extensions, but the iteration which we produce is itself $1$-allowable.  We should therefore also test persistence under $1$-allowable extensions.  Repeating this argument leads to the transfinite hierarchy defined in the next subsection.


\subsection[The hierarchy of alpha-allowable iterations]{The hierarchy of $\alpha$-allowable iterations}
\label{subsection-alpha-allowability}

We formulate the definition for iterations over the fixed reservoir.  A forcing is called $\alpha$-allowable if it is isomorphic to the final forcing of an $\alpha$-allowable iteration.

We use the canonical $<_L$-coding convention for objects in $H(\omega_2)^W$.  Every object used below is represented by its $<_L$-least $\operatorname{Br}$-name in $L$, and codes are compared by $<_L$.

We also fix a support convention which takes account of both parts of an allowable iteration.  For $E\subseteq\omega_1$, put
\[
   \forceC^{\mathrm{area}}_E
   :=\prod_{\eta\in E}^{\mathrm{cs}}\mathbb C_L(\omega_1).
\]
Suppose that
\[
   \forceP_\beta
   =\forceC^{\mathrm{area}}*\dot{\forceR}_\beta
\]
is an initial segment of an allowable iteration relative to
$(F,\vec\eta)$, and put
\[
   U_\beta:=\operatorname{ran}(\vec\eta\upharpoonright\beta).
\]
Let $\vec\tau$ be a finite tuple of $\forceP_\beta$-names for reals or
natural numbers.  A pair
\[
   (E,A)\in[U_\beta]^{\leq\omega}\times[\beta]^{\leq\omega}
\]
is \emph{admissible for $\vec\tau$} if the induced iteration on the stages in
$A$, using only the reservoir coordinates in $E$, is recursively defined, its
final forcing
\[
   \forceP_{E,A}
   :=\forceC^{\mathrm{area}}_E*\dot{\forceR}_{E,A}
\]
is a complete subforcing of $\forceP_\beta$, and every member of $\vec\tau$
is the canonical lift of a $\forceP_{E,A}$-name.  In particular, the pair is
closed under the names defining the retained iterands: if $\gamma\in A$, then
the coding coordinate $\eta_\gamma$ belongs to $E$, and the names which define
the $\gamma$-th iterand are already names over the induced iteration on
$A\cap\gamma$ and the reservoir coordinates in $E$.

The reservoir support $E$ is countable and is contained in $U_\beta$.  Its
role is to record the coordinates which occur in the names defining the
relevant part of the almost disjoint iteration.  It does not record new real
information.  Indeed, the reservoir and each of its coordinate factors are
$L$-computed $\sigma$-closed forcings.  After the c.c.c. forcing which
produces $W$ and the c.c.c. almost disjoint tail below $\beta$, Easton's lemma
implies that the unused reservoir factor on
$\omega_1\setminus U_\beta$ is $\omega$-distributive.  Thus every real in the
full intermediate extension has an equivalent name which uses only
coordinates in $U_\beta$.  The same observation applies after a permutation
of the reservoir coordinates.  Combining this with the usual nice-name
argument for the c.c.c. tail, and then closing recursively under the names
defining the retained iterands, shows that every finite tuple of reals and
natural numbers has a tuple of equivalent names with an admissible support pair whose two coordinates are countable.

We put
\[
   \operatorname{supp}_{\mathrm{AP}}(\vec\tau,\forceP_\beta)
   :=\text{the }<_L\text{-least pair }(E,A)
      \text{ which is admissible for }\vec\tau.
\]
The iteration support $A$ need not be an ordinal.

For an admissible support pair $\mathfrak a=(E,A)$ occurring in this sense, put
\[
   \forceP_{\mathfrak a}:=\forceP_{E,A}.
\]
Whenever an ambient generic $K*G_\beta$ is fixed, write
\[
   G_{\mathfrak a}:=K_E*G_A,
   \qquad
   M_{\mathfrak a}:=W[K_E][G_A].
\]
If $\mathfrak b=(F,B)$ and $\mathfrak a=(E,A)$ are admissible support pairs,
$F\subseteq E$, $B\subseteq A$, and $\forceP_{\mathfrak b}$ is a complete
subforcing of $\forceP_{\mathfrak a}$, we write
\[
   \mathfrak b\preccurlyeq\mathfrak a.
\]

For the persistence relation we distinguish literal end-extensions from complete subiterations obtained after reindexing the reservoir.  The notation
\[
   \forceQ\triangleright\forceP
\]
continues to mean that $\forceQ$ literally end-extends $\forceP$ over the same reservoir, as in Definition~\ref{definition-allowable-iteration}.

We introduce a second relation which is invariant under the normalization of products.  A \emph{reservoir block} is a pair $(B,e)$ such that
\[
   B\in L,\qquad B\subseteq\omega_1,\qquad |B|=\omega_1,
\]
and
\[
   e:\omega_1\longrightarrow B
\]
is a bijection belonging to $L$.  If $\forceP$ is an allowable iteration over the reservoir, let $e_*\forceP$ denote its reindexed copy obtained by replacing
\[
   \forceC^{\mathrm{area}}
   \quad\text{with}\quad
   \forceC^{\mathrm{area}}_B,
\]
replacing every canonical coordinate name $\dot g_\eta$ by $\dot g_{e(\eta)}$, and translating recursively all names in the finite-support tail.

Let $\forceP$ and $\forceQ$ be allowable iterations of lengths $\delta$ and $\varepsilon$, respectively.  We write
\[
   \forceQ\succeq_{\mathrm{res}}\forceP
\]
and say that $\forceQ$ is a \emph{complete reservoir extension} of $\forceP$ if there are a reservoir block $(B,e)$ and a set $A\in W$, $A\subseteq\varepsilon$, of order type $\delta$ such that the following hold.  Writing
\[
   \pi:\delta\longrightarrow A
\]
for the increasing enumeration of $A$, the stages in $A$ are closed under the names defining their iterands when only the reservoir coordinates in $B$ are retained.  Hence the induced iteration
\[
   \forceQ_{B,A}
   :=\forceC^{\mathrm{area}}_B*\dot{\forceR}^{\forceQ}_{B,A}
\]
is recursively defined and is a complete subforcing of $\forceQ$.  Moreover, reindexing the reservoir by $e$ and the stages by $\pi$ induces an isomorphism
\[
   \iota_{e,\pi}:\forceP\cong\forceQ_{B,A}.
\]
The reservoir block, the set of stages, and the isomorphism are regarded as part of the witness to
$\forceQ\succeq_{\mathrm{res}}\forceP$.

The relation $\succeq_{\mathrm{res}}$ is reflexive and transitive.  Reflexivity is witnessed by the identity reindexing and the full set of stages.  For transitivity, compose the two reservoir reindexings and the two increasing maps between the retained stages; complete subforcings compose.  In particular,
\[
   \forceQ\triangleright\forceP
   \quad\Longrightarrow\quad
   \forceQ\succeq_{\mathrm{res}}\forceP.
\]
Further, if an allowable iteration is the normalized reservoir presentation of
\[
   \forceP^0\times\forceP^1,
\]
obtained by moving the two copies of the reservoir to disjoint blocks of coordinates and concatenating their almost disjoint tails, then each factor occurs in this iteration as a complete reservoir subiteration.

If $G$ is $\forceP$-generic and
\[
   \forceQ\succeq_{\mathrm{res}}\forceP
\]
is witnessed by $\iota_{e,\pi}$, we identify $G$ with its image under this isomorphism and write
\[
   \forceQ/G
\]
for the quotient forcing over the chosen complete reservoir subiteration.

We need a way how to single out reals which should become the values of our uniformizing functions. Suppose we are given a real $x$ and a $\Pi^1_3$ set $A_m$ and we know that there is a real $y$ such that $(x,y)$ is persistent in $A_m$, that is it will remain in $A_m$ in all further allowable extensions. The task is to determine which real $z$ to pick as a value for $f(m,x)$. A first naive attempt is to simply pick the $<_L$-least name $\sigma$ for such a real and let $z=\sigma^G$, for $G$ the generic filter, but this does not work for technical reasons. 

How to do this is described now. Let $\mathfrak a=(E,A)$ be an
admissible support pair and let $z\in M_{\mathfrak a}$ be a real, or a fixed
finite tuple of reals and natural numbers.  In the application
$z=\langle x,y,m\rangle$.  A \emph{presentation of $z$ below $\mathfrak a$}
is a quadruple
\[
   c=(\mathfrak b,\forceR,p,\dot\sigma)
\]
such that
\[
   \mathfrak b\preccurlyeq\mathfrak a,
   \qquad
   z\in M_{\mathfrak b},
\]
$\forceR\in M_{\mathfrak b}$ is allowable over $M_{\mathfrak b}$, the coding
areas used by $\forceR$ are disjoint from those already used by $\mathfrak b$,
$p\in\forceR$, and
\[
   M_{\mathfrak b}\models
   p\Vdash_{\forceR}\dot\sigma=\check z .
\]
Using the fixed coding convention, $c$ determines the pointed allowable
presentation
\[
   d_c=
   (\forceP_{\mathfrak b}*\dot\forceR,
      (1_{\forceP_{\mathfrak b}},\dot p),\dot\sigma),
\]
where $\dot\forceR$ and $\dot p$ are the canonical $\forceP_{\mathfrak b}$-names
for $\forceR$ and $p$, and $\dot\sigma$ is lifted to the two-step forcing.

For pointed allowable presentations
\[
   d=(\forceS,q,\dot\tau),
   \qquad
   e=(\forceT,r,\dot\sigma),
\]
let $\equiv_{\mathrm{pres}}$ be the equivalence relation generated by the
following relation: $d$ and $e$ are related if there are a normalized
allowable forcing $\forceU$, allowable complete embeddings
\[
   i:\forceS\longrightarrow\forceU,
   \qquad
   j:\forceT\longrightarrow\forceU,
\]
and a condition $u\in\forceU$ such that
\[
   u\leq i(q),j(r)
   \qquad\text{and}\qquad
   u\Vdash_{\forceU} i(\dot\tau)=j(\dot\sigma).
\]
Here allowable complete embedding means that the quotient forcing is
allowable.  Define
\[
   \pcode(d)
   :=\min_{<_L}
      \{\operatorname{code}(e):e\equiv_{\mathrm{pres}} d\}.
\]
Thus $\pcode(d)$ is the canonical code of the presentation class of $d$.  The
definition is made in the ground model from the fixed coding of allowable
forcings, embeddings, conditions, names, and the forcing relation.

For $z\in M_{\mathfrak a}$, we define the orbit of $z$ relative to $\mathfrak a$ to be
\[
   \mathcal O_{\mathfrak a}(z)
   :=
   \{\pcode(d_c):c\text{ is a presentation of }z
      \text{ below }\mathfrak a\}
\]
and
\[
   \key_{\mathfrak a}(z)
   :=\min_{<_L}
      \mathcal O_{\mathfrak a}(z).
\]
The orbit is nonempty, since we may take $\mathfrak b=\mathfrak a$ and $\forceR$
trivial.  It remembers only the allowable presentation class of $z$, not the
support or factorization in which $z$ was first seen.

\begin{lemma}
\label{lemma-semantic-orbit-support-independence}
Suppose that
\[
   \mathfrak b\preccurlyeq\mathfrak a
   \qquad\text{and}\qquad
   z\in M_{\mathfrak b}.
\]
Then
\[
   \mathcal O_{\mathfrak b}(z)=\mathcal O_{\mathfrak a}(z),
\]
and therefore
\[
   \key_{\mathfrak b}(z)=\key_{\mathfrak a}(z).
\]
\end{lemma}

\begin{proof}
The inclusion $\mathcal O_{\mathfrak b}(z)\subseteq\mathcal O_{\mathfrak a}(z)$ is immediate.
Conversely, let
\[
   c=(\mathfrak c,\forceR,p,\dot\sigma)
\]
be a presentation of $z$ below $\mathfrak a$.  Let $\mathfrak d$ be the admissible
closure of $\mathfrak b\cup\mathfrak c$ inside $\mathfrak a$.  Since
$z\in M_{\mathfrak b}\cap M_{\mathfrak c}$, the lifted $\mathfrak b$-presentation
of $z$ and the lifted $\mathfrak c$-presentation of $z$ have the same value in
$M_{\mathfrak d}$.  By the truth lemma, after strengthening in the quotient to
$\mathfrak d$, this equality is forced.

Thus $d_c$ and the presentation obtained from $\mathfrak b$, followed by the
quotient to $\mathfrak d$ and then by $\forceR$, have a common allowable
refinement below which their names are forced to be equal.  They therefore
determine the same presentation class.  Hence
$\pcode(d_c)\in\mathcal O_{\mathfrak b}(z)$.  This proves the reverse
inclusion, and the equality of keys follows by taking the $<_L$-least
element of the common orbit.
\end{proof}

At a stage with local admissible support $\mathfrak a$ and given a pair $(x,m)$, where $x$ is a real and $m \in \omega$,
we first minimize the persistence rank.  Among the reals of least rank we compare
\[
   \key_{\mathfrak a}(\langle x,y,m\rangle),
\]
and choose the real with $<_L$-least key.  This is the canonical selection rule
for the value of $f(m,x)$.


We have assembled all the relevant notions to define the main concept of this work.
We define $\rho$-allowability by induction on $\rho$.

\begin{definition}
\label{definition-alpha-allowability}
A $0$-allowable iteration is an allowable iteration over the reservoir in the sense of Definition~\ref{definition-allowable-iteration}.  For purposes of the hierarchy, it is paired with the canonical auxiliary name for the empty set.

Let $0<\rho$, and assume that $\zeta$-allowability has been defined for every $\zeta<\rho$.  Let $\delta<\omega_1$, let
\[
   F=\left\langle
      (\dot x_\beta,\dot z_\beta,\dot m_\beta)
      \mid\beta<\delta
   \right\rangle\in W
\]
be a bookkeeping function, and let
\[
   \vec\eta=\langle\eta_\beta\mid\beta<\delta\rangle\in L
\]
be an injective sequence of ordinals below $\omega_1$.  For
$\beta\leq\delta$, put
\[
   U_\beta:=\operatorname{ran}(\vec\eta\upharpoonright\beta).
\]
Here $\dot z_\beta$ is the bookkeeping guess for a value in the
$\dot x_\beta$-section of $A_{\dot m_\beta}$.  We require that, whenever
$\forceP_\beta$ has been defined and
\[
   (E_\beta,A_\beta)
   :=\operatorname{supp}_{\mathrm{AP}}
      \bigl(F(\beta),\forceP_\beta\bigr),
\]
then
\[
   \forceP_{E_\beta,A_\beta}\Vdash
   ``\dot x_\beta,\dot z_\beta\in2^\omega
     \text{ and }\dot m_\beta\in\omega,"
\]
and
\[
   E_\beta\subseteq U_\beta.
\]
Thus the local bookkeeping data at stage $\beta$ use only coding-area
coordinates assigned before stage $\beta$.  Since $\vec\eta$ is injective,
\[
   \eta_\beta\notin U_\beta,
\]
and therefore $\eta_\beta\notin E_\beta$.  The entire assignment belongs to
$L$, so it is computed in the same way in every branch inner model of $W$.
By Easton's lemma, requiring stage-respecting names does not remove any real
which can occur in the construction.  We shall consider only pairs
$(F,\vec\eta)$ satisfying these requirements.

A $\rho$-allowable iteration relative to $(F,\vec\eta)$ is a sequence
\[
   \forceP=\langle\forceP_\beta\mid\beta\leq\delta\rangle,
   \qquad
   \forceP_\beta=\forceC^{\mathrm{area}}*\dot{\forceR}_\beta,
\]
together with auxiliary names $\langle\dot I_\beta\mid\beta\leq\delta\rangle$, defined as follows.

At stage $0$,
\[
   \dot{\forceR}_0=\mathbf 1
   \qquad\text{and}\qquad
   \dot I_0=\varnothing.
\]
Suppose that $\beta<\delta$, that $\forceP_\beta$ and $\dot I_\beta$ have been defined, and let
\[
   K*G_\beta
\]
be $\forceP_\beta$-generic over $W$.  Put
\[
   E:=E_\beta,
   \qquad
   A:=A_\beta,
   \qquad
   K_E:=K\upharpoonright E,
   \qquad
   G_A:=G_\beta\upharpoonright A,
\]
so that
\[
   G_{E,A}:=K_E*G_A
\]
is $\forceP_{E,A}$-generic over $W$.  Let
\[
   x:=\dot x_\beta^{G_{E,A}},\qquad
   z:=\dot z_\beta^{G_{E,A}},\qquad
   m:=\dot m_\beta^{G_{E,A}}.
\]
The coding area assigned to stage $\beta$ is the reservoir coordinate
$\eta_\beta$.  Since
\[
   E\subseteq U_\beta
   \qquad\text{and}\qquad
   \eta_\beta\notin U_\beta,
\]
the area used at this stage is outside the local reservoir support.
Earlier names are identified with their canonical lifts.

Put
\[
   \mathfrak a_\beta:=(E,A).
\]
For $\zeta<\rho$ and $y\in M_{\mathfrak a_\beta}\cap2^\omega$, say that $y$ is \emph{$\zeta$-persistent for $(x,m)$ over $\forceP_\beta$} if
\[
   W[K][G_\beta]\models (x,y)\in A_m
\]
and, for every $\zeta$-allowable iteration $\forceQ\in W$ together with a fixed witness to
\[
   \forceQ\succeq_{\mathrm{res}}\forceP_\beta,
\]
we have
\[
   W[K][G_\beta]\models
   ``\forceQ/(K*G_\beta)\Vdash(x,y)\in A_m".
\]
Equivalently, using the standing convention of restricting a forcing below a condition, no $\zeta$-allowable complete reservoir extension of $\forceP_\beta$ forces
\[
   (x,y)\notin A_m.
\]

Let
\[
\begin{aligned}
   C_\beta^\rho(K*G_\beta)
   :=\bigl\{(\zeta,y)\in
      \rho\times\bigl(W[K_E][G_A]\cap2^\omega\bigr):{}&
      y\text{ is }\zeta\text{-persistent}\\
      &\text{for }(x,m)\text{ over }\forceP_\beta
   \bigr\}.
\end{aligned}
\]
We distinguish two cases.

\medskip
\noindent
\textup{(a) The first uniformization case.}
Assume that $C_\beta^\rho(K*G_\beta)\neq\varnothing$.  Choose
\[
   \zeta_\beta
   :=\min\{\zeta:\exists y\ ((\zeta,y)\in C_\beta^\rho(K*G_\beta))\},
\]
and then choose $y_\beta$ so that $(\zeta_\beta,y_\beta)\in C_\beta^\rho(K*G_\beta)$ and
\[
   \key_{\mathfrak a_\beta}(\langle x,y_\beta,m\rangle)
\]
is $<_L$-least among the keys of all such reals.  Let $z_\beta^*=z$ if
$z\neq y_\beta$.  If $z=y_\beta$, let $z_\beta^*$ be the real
$w\neq y_\beta$ in $W[K_E][G_A]$ for which
\[
   \key_{\mathfrak a_\beta}(\langle x,w,m\rangle)
\]
is $<_L$-least. 

The next almost disjoint coding factor is
\[
   \dot{\forceA}_\beta
   :=\mathbb A_D\bigl(
      \dot Y(\dot x_\beta,\dot z_\beta^*,\dot m_\beta,
             \dot g_{\eta_\beta})
   \bigr),
\]
where $\dot z_\beta^*$ is the name obtained from the preceding uniform choice.  We put
\[
   \dot{\forceR}_{\beta+1}
   :=\dot{\forceR}_\beta*\dot{\forceA}_\beta
\]
and, for $K*G_{\beta+1}$ extending $K*G_\beta$, define
\[
   \dot I_{\beta+1}^{K*G_{\beta+1}}
   :=\dot I_\beta^{K*G_\beta}
     \cup\{(x,y_\beta,m,\zeta_\beta)\}.
\]
In this case $y_\beta$ is declared to be a tentative value for $f_m(x)$ of rank $\zeta_\beta$.  The triple $(x,z_\beta^*,m)$ is coded instead.

\medskip
\noindent
\textup{(b) The second uniformization case.}
Assume that $C_\beta^\rho(K*G_\beta)=\varnothing$.  In this case no real in $M_{\mathfrak a_\beta}\cap2^\omega$ passes the test at any rank below $\rho$, and the construction uses the bookkeeping guess.  We put
\[
   \dot{\forceA}_\beta
   :=\mathbb A_D\bigl(
      \dot Y(\dot x_\beta,\dot z_\beta,\dot m_\beta,
             \dot g_{\eta_\beta})
   \bigr),
\]
\[
   \dot{\forceR}_{\beta+1}
   :=\dot{\forceR}_\beta*\dot{\forceA}_\beta,
\]
and
\[
   \dot I_{\beta+1}^{K*G_{\beta+1}}
   :=\dot I_\beta^{K*G_\beta}.
\]
Thus no new tentative value is added at stage $\beta$.

If $\lambda\leq\delta$ is a limit ordinal, let $\dot{\forceR}_\lambda$ be the finite-support limit of
\[
   \langle\dot{\forceR}_\beta\mid\beta<\lambda\rangle
\]
and define, for $K*G_\lambda$ generic,
\[
   \dot I_\lambda^{K*G_\lambda}
   :=\bigcup_{\beta<\lambda}
      \dot I_\beta^{K*G_\beta}.
\]
All choices above are made from the fixed order $<_L$ on codes and are uniform in $K*G_\beta$.  Hence they determine forcing names in $W$.  The notation $K*G_\beta$ is only semantic shorthand; no generic filter is chosen during the recursion.

If the recursion produces $(\forceP_\delta,\dot I_\delta)$, we say that this pair is $\rho$-allowable relative to $(F,\vec\eta)$.  We call $\forceP_\delta$ $\rho$-allowable if it is $\rho$-allowable relative to some bookkeeping function $F$ and some injective area assignment $\vec\eta\in L$ satisfying the displayed requirements.
\end{definition}

For a generic filter $K*G$ for the final forcing, we write
\[
   I^{K*G}:=\dot I_\delta^{K*G}.
\]
The interpreted auxiliary set $I^{K*G}\in W[K][G]$ is the set of tentative values for the uniformizing functions.  If $(x,y,m,\xi)\in I^{K*G}$, then $y$ is a tentative value for $f_m(x)$ of rank $\xi$.  The same value may occur with more than one rank.  The auxiliary set does not enter the definition of the forcing factors; it records the tentative values selected by the recursion.

The case $\rho=1$ is the notion of $1$-allowability.  Since the only ordinal below $1$ is $0$, clause~\textup{(a)} asks whether there is a value which is persistent under all allowable, that is, $0$-allowable, continuations.  Thus Definition~\ref{definition-alpha-allowability} specializes to the intended notion of $1$-allowability when $\rho=1$.

The restriction to reals in
\[
   W[K_E][G_A]
\]
is part of the definition.  We do not scan all reals in $W[K][G_\beta]$.  The reservoir forcing itself adds no reals, by Easton's lemma, but the coordinate set $E$ is still needed to record the parameters on which the induced almost disjoint iteration depends.  This support stratification is used in the product argument: after the reservoir coordinates and iteration stages of two factors are moved to disjoint blocks, the names belonging to one factor are evaluated in the corresponding translated complete subiteration.

At a successor level $\rho+1$, the definition permits the additional rank $\rho$ and therefore tests persistence under $\rho$-allowable continuations.  At a limit $\lambda$, it tests all ranks $\zeta<\lambda$ at once.  No separate limit-stage definition of allowability is required.

For formulas whose relevant instances are preserved by every forcing under consideration, clause~\textup{(a)} must apply whenever the section is nonempty.  Thus tentative values already occur at the first derivative.  We will not need a separate assertion that the inclusions between consecutive classes are proper.

\begin{lemma}
\label{shrinkinglemma}
Let $\xi<\alpha$.  Every $\alpha$-allowable forcing is $\xi$-allowable.  Consequently,
\[
   \mathcal A^W_0\supseteq\mathcal A^W_1\supseteq\cdots
   \supseteq\mathcal A^W_\alpha\supseteq\cdots,
\]
where $\mathcal A^W_\alpha$ denotes the class of $\alpha$-allowable forcings over $W$.
\end{lemma}

\begin{proof}
The case $\xi=0$ is immediate, since $0$-allowability is just allowability.
Let $0<\xi<\alpha$, and let $(\forceP,\dot I)$ be $\alpha$-allowable
relative to $(F,\vec\eta)$.  We define, by recursion on the stages, a
bookkeeping function $F'$ such that the $\xi$-construction relative to
$(F',\vec\eta)$ produces the same forcing iteration.

Assume that the two constructions agree below a stage $\beta$.  Let
$\mathfrak a_\beta=(E_\beta,A_\beta)$ be the local admissible support used by
the $\alpha$-construction at this stage, and let the corresponding local data
be $(x,z,m)$.

Suppose first that the $\alpha$-construction is in clause~\textup{(a)} and
selects a real $y_\beta$ of rank $\rho<\xi$.  For every $\zeta<\xi$, the
persistence test is the same in the two constructions, since it only refers to
$\zeta$-allowable complete reservoir extensions of the common initial forcing.
The local support is the same, and the key
\[
   \key_{\mathfrak a_\beta}(\langle x,y,m\rangle)
\]
is computed from the same orbit.  Hence the two constructions have the same
least rank below $\xi$ and then choose the same real by the $<_L$-least key.
We put $F'(\beta)=F(\beta)$, and the same coding factor is used.

Suppose next that the $\alpha$-construction is in clause~\textup{(a)} but
selects a real of rank $\rho\geq\xi$.  Then no real in
$M_{\mathfrak a_\beta}$ passes the persistence test at any rank below $\xi$.
Thus the $\xi$-construction must be in clause~\textup{(b)}.  We choose
$F'(\beta)$ to be a stage-respecting name, over the same local support, for
the triple which the $\alpha$-construction codes at stage $\beta$.  Clause
\textup{(b)} of the $\xi$-construction therefore adds exactly the same almost
disjoint coding factor.

Finally, if the $\alpha$-construction is already in clause~\textup{(b)}, then
no real passes the test at any rank below $\alpha$, hence none passes at a
rank below $\xi$.  We put $F'(\beta)=F(\beta)$, so the $\xi$-construction is
also in clause~\textup{(b)} and uses the same factor.

The choices made in the recursion are uniform in the generic filter.  The new
bookkeeping names are chosen in the local models used by the
$\alpha$-construction, whose reservoir supports are contained in the
corresponding $U_\beta$.  Hence $(F',\vec\eta)$ satisfies the stage-respecting
support requirement.  At limit stages both constructions take the same
finite-support limit.  Therefore the final forcing of the given
$\alpha$-allowable iteration is also $\xi$-allowable.
\end{proof}

\begin{lemma}
\label{1allowableclosedunderproducts}
\label{alphaallowableclosedunderproducts}
Let $\alpha$ be an ordinal.  If $\forceP^0$ and $\forceP^1$ are
$\alpha$-allowable, then
\[
   \forceP^0\times\forceP^1
\]
is isomorphic to an $\alpha$-allowable forcing.  More precisely, a normalized
$\alpha$-allowable presentation of the product is obtained from the two given
presentations by an $L$-definable reindexing of the two reservoirs onto
disjoint blocks and by concatenating the translated almost disjoint tails.  The
resulting bookkeeping function belongs to $W$, and the resulting injective area
assignment belongs to $L$.
\end{lemma}

\begin{proof}
We argue by induction on $\alpha$, proving the normalized form stated above.
For $\alpha=0$, this is exactly
Lemma~\ref{basic-properties-allowable-forcings}(3).

Assume that the assertion holds for every $\xi<\alpha$.  Let
$(F_i,\vec\eta^{\,i})$ witness the $\alpha$-allowability of
$\forceP^i$, and let $\delta_i$ be the length of $\forceP^i$
$(i<2)$.  Fix in $L$ a partition
\[
   \omega_1=B_0\mathbin{\dot\cup}B_1
\]
and bijections
\[
   e_i:\omega_1\longrightarrow B_i
   \qquad(i<2).
\]
Move the reservoir of $\forceP^i$ to $B_i$ by $e_i$.  We present the product
by first taking the translated almost disjoint tail of $\forceP^0$ and then
the translated tail of $\forceP^1$.  Let
\[
   \forceS=\langle\forceS_\gamma\mid
      \gamma\leq\delta_0+\delta_1\rangle
\]
be this normalized product presentation.  Its final forcing is isomorphic to
$\forceP^0\times\forceP^1$.  Put
\[
   \vec\eta^{\,\otimes}
   =
   \langle e_0(\eta^{\,0}_\beta)\mid\beta<\delta_0\rangle
    {}^\frown
    \langle e_1(\eta^{\,1}_\beta)\mid\beta<\delta_1\rangle.
\]
Then $\vec\eta^{\,\otimes}\in L$ is injective.  Let
$F^\otimes$ be obtained by translating $F_0$ and $F_1$ by $e_0$ and
$e_1$, respectively, and concatenating the translated sequences.  Then
$F^\otimes\in W$.

It remains to check the stage rule.  The first block is handled by the same
argument as below, with only the reindexing $e_0$, so we consider a stage
$\delta_0+\beta$ in the second block.  Let $\beta<\delta_1$, and let
\[
   \mathfrak a^1_\beta=(E,A)
\]
be the local admissible support of the corresponding stage in the construction
of $\forceP^1$.  In the product presentation the translated support is
\[
   \mathfrak a^\otimes_{\delta_0+\beta}
   :=(e_1[E],\delta_0+A),
   \qquad
   \delta_0+A=\{\delta_0+\gamma:\gamma\in A\}.
\]
It is admissible, since
\[
   E\subseteq
   \operatorname{ran}(\vec\eta^{\,1}\upharpoonright\beta)
\]
implies
\[
   e_1[E]\subseteq
   \operatorname{ran}
      (\vec\eta^{\,\otimes}\upharpoonright(\delta_0+\beta)).
\]
The local models
\[
   W[K\upharpoonright e_1[E]]
    [G\upharpoonright(\delta_0+A)]
   \quad\text{and}\quad
   W[K^1\upharpoonright E][G^1\upharpoonright A]
\]
are canonically isomorphic.  We identify their reals through this isomorphism.

Fix $\xi<\alpha$ and a local real $y$.  We claim that $y$ is
$\xi$-persistent for $(x,m)$ at the stage $\beta$ of $\forceP^1$ iff its
translated copy is $\xi$-persistent at the stage $\delta_0+\beta$ of
$\forceS$.  If $y$ is not persistent in $\forceP^1$, let $\forceQ$ be a
$\xi$-allowable complete reservoir extension of $\forceP^1_\beta$ whose
quotient forces
\[
   (x,y)\notin A_m.
\]
By Lemma~\ref{shrinkinglemma}, $\forceP^0$ is $\xi$-allowable, and by the
induction hypothesis $\forceP^0\times\forceQ$ has a normalized
$\xi$-allowable presentation.  This presentation contains
$\forceS_{\delta_0+\beta}$ as a complete reservoir subiteration, and its
quotient still forces
\[
   (x,y)\notin A_m.
\]
Hence the translated $y$ is not persistent in the product presentation.

Conversely, if a $\xi$-allowable complete reservoir extension of
$\forceS_{\delta_0+\beta}$ forces
\[
   (x,y)\notin A_m,
\]
then, since $\forceS_{\delta_0+\beta}$ contains the translated copy of
$\forceP^1_\beta$ as a complete reservoir subiteration, the same forcing is
also a $\xi$-allowable complete reservoir extension of $\forceP^1_\beta$ by
transitivity of $\succeq_{\mathrm{res}}$.  Thus $y$ is not persistent in the
original construction of $\forceP^1$.  This proves the claim.

The key comparison is also preserved.  The reindexing $e_1$ sends presentations
below $\mathfrak a^1_\beta$ bijectively to presentations below
$\mathfrak a^\otimes_{\delta_0+\beta}$, and the corresponding pointed
presentations have the same presentation class.  Hence, for every local real
$y$,
\[
   \key_{\mathfrak a^1_\beta}
      (\langle x,y,m\rangle)
   =
   \key_{\mathfrak a^\otimes_{\delta_0+\beta}}
      (\langle x,y,m\rangle).
\]
Equivalently, the two constructions compute the same orbits and therefore the
same $<_L$-least keys.  Thus they have the same least persistence rank and,
at that rank, choose the same real.  If the selected value equals the
bookkeeping value, the same key comparison also chooses the same real to be
coded instead.  Therefore the translated stage of $\forceS$ uses exactly the
translated coding factor from $\forceP^1$.

The same verification applies to every stage of the second block, and the
first block was already covered by the reindexing argument.  At limits the
finite-support limits agree.  Hence $\forceS$ is an $\alpha$-allowable
iteration relative to $(F^\otimes,\vec\eta^{\,\otimes})$.  Since the final
forcing of $\forceS$ is isomorphic to $\forceP^0\times\forceP^1$, the product
is isomorphic to an $\alpha$-allowable forcing.
\end{proof}

\begin{lemma}
\label{fvaluesremain2}
Let $(\forceP,\dot I)$ be $\alpha$-allowable, let $G$ be $\forceP$-generic over $W$, and suppose that
\[
   (x,y,m,\xi)\in\dot I^G
   \qquad\text{for some }\xi<\alpha.
\]
Then
\[
   W[G]\models(x,y)\in A_m.
\]
Moreover, if $\forceQ$ is $\xi$-allowable together with a fixed witness to
\[
   \forceQ\succeq_{\mathrm{res}}\forceP,
\]
then
\[
   W[G]\models
   ``\forceQ/G\Vdash(x,y)\in A_m".
\]
\end{lemma}

\begin{proof}
Let $\beta$ be the stage at which the occurrence $(x,y,m,\xi)$ is first
entered into the auxiliary name.  At that stage the construction is in
clause~\textup{(a)}, and by definition the selected real $y$ is
$\xi$-persistent for $(x,m)$ over $\forceP_\beta$.

By Lemma~\ref{shrinkinglemma}, the final forcing $\forceP$ is
$\xi$-allowable.  Moreover
\[
   \forceP\succeq_{\mathrm{res}}\forceP_\beta.
\]
Applying the definition of $\xi$-persistence to this complete reservoir
extension gives
\[
   W[G]\models (x,y)\in A_m.
\]

If $\forceQ$ is as in the second assertion, then transitivity of
$\succeq_{\mathrm{res}}$ gives
\[
   \forceQ\succeq_{\mathrm{res}}\forceP_\beta.
\]
The same persistence statement at stage $\beta$ therefore yields
\[
   W[G]\models ``\forceQ/G\Vdash (x,y)\in A_m'',
\]
which is the desired conclusion.
\end{proof}

\begin{lemma}
\label{nonempty-alpha-allowable}
For every ordinal $\alpha$, the class $\mathcal A^W_\alpha$ is nonempty.
\end{lemma}

\begin{proof}
We argue by induction on $\alpha$.  The class $\mathcal A^W_0$ is nonempty.  Suppose that $0<\alpha$ and that $\xi$-allowability has been defined for every $\xi<\alpha$.  Choose a bookkeeping function of positive countable length satisfying the displayed requirements, and use, for example, the canonical injective area assignment $\eta_\beta=\beta$, which belongs to $L$.  The bookkeeping may be chosen to use check-names, so its reservoir support is empty and hence contained in $U_\beta$ at every stage.  At each stage, either the set $C_\beta^\alpha(K*G_\beta)$ is nonempty and clause~\textup{(a)} applies, or it is empty and clause~\textup{(b)} applies.  In either case the next factor is an almost disjoint coding forcing of the form required in Definition~\ref{definition-allowable-iteration}.  The recursion therefore produces an $\alpha$-allowable iteration.
\end{proof}

There are only set many isomorphism types of allowable forcings in $W$: every such forcing has size at most $\aleph_1$ and belongs to $H(\omega_2)^W$.  Since the sequence
\[
   \langle\mathcal A^W_\alpha\mid\alpha\in\mathrm{Ord}\rangle
\]
is decreasing and each $\mathcal A^W_\alpha$ is nonempty, it eventually stabilizes.  Let
\[
   \alpha_0
   :=\min\left\{\alpha:
      \forall\beta\geq\alpha\ \bigl(\mathcal A^W_\beta=\mathcal A^W_\alpha\bigr)
   \right\}.
\]

Set
\[
   \alpha^\ast:=\alpha_0+1.
\]

\begin{definition}
An \emph{$\infty$-allowable iteration} is an $\alpha^\ast$-allowable iteration.  A forcing is \emph{$\infty$-allowable} if it is isomorphic to the final forcing of such an iteration.
\end{definition}

Since
\[
   \mathcal A^W_{\alpha^\ast}=\mathcal A^W_{\alpha_0},
\]
the class of $\infty$-allowable forcings is the fixed class $\mathcal A^W_{\alpha_0}$.  The use of the successor index $\alpha^\ast$ ensures that the stage rule also tests persistence under the fixed class itself.  Equivalently, in clause~\textup{(a)} of an $\infty$-allowable iteration, we minimize over
\[
   \xi<\alpha^\ast.
\]
Thus every tentative value entered into the auxiliary set of an $\infty$-allowable iteration has rank
\[
   \xi<\alpha^\ast.
\]




\section[Definition of the universe with Pi13 uniformization]{Definition of the universe in which the ${\Pi^1_3}$ uniformization property holds}
\label{section-L-final-uniformization}
The notion of $\infty$-allowability will now be used to define the universe in which the ${\Pi^1_3}$-uniformization property holds.  Let $W$ be our ground model.  We construct an $\omega_1$-length iteration whose countable initial segments are $\infty$-allowable over $W$.

Fix a bookkeeping function
\[
   F:\omega_1\longrightarrow H(\omega_1)^W
\]
such that
\begin{itemize}
\item $F$ is surjective;
\item for every $a\in H(\omega_1)^W$, the set $F^{-1}(\{a\})$ is unbounded in $\omega_1$.
\end{itemize}

Fix the canonical area assignment
\[
   \vec\eta=\langle\eta_\beta\mid\beta<\omega_1\rangle
   =\langle\beta\mid\beta<\omega_1\rangle\in L.
\]
Thus
\[
   U_\beta
   =\operatorname{ran}(\vec\eta\upharpoonright\beta)
   =\beta,
\]
and the coordinate $\eta_\beta=\beta$ supplies the coding area used at stage
$\beta$.

We define in $W$, by recursion on $\beta<\omega_1$, an iteration $\forceP_{\omega_1}$ together with a sequence of auxiliary names
\[
   \langle\dot I_\beta\mid\beta\leq\omega_1\rangle.
\]
No generic filter is chosen during this recursion.  The notation below involving a generic $K*G_\beta$ is only a semantic description of the canonical names selected in $W$.  At a limit stage $\lambda\leq\omega_1$, we keep the fixed reservoir $\forceC^{\mathrm{area}}$ and take the finite-support limit of the almost disjoint iteration.  We let $\dot I_\lambda$ be the canonical name forced to be the union of the earlier interpreted auxiliary names.

Suppose that $\beta<\omega_1$ and that $(\forceP_\beta,\dot I_\beta)$ has already been constructed.  Assume that the value $F(\beta)$ decodes as a pair $(F(\beta)_0,F(\beta)_1)$ and that
\[
   F(\beta)_0=(\eta_1,\eta_2,m),
   \qquad \eta_1\leq\beta,
   \qquad m\in\omega.
\]
Let $(\dot x,\dot y)$ be the $\eta_2$-th pair, in the canonical
$<_L$-order fixed above, of nice $\forceP_{\eta_1}$-names in
stage-respecting normal form; equivalently, their canonical reservoir support
is contained in
\[
   U_{\eta_1}=\eta_1.
\]
Lift these names canonically to $\forceP_\beta$.  By Easton's lemma, every
pair of reals in the $\forceP_{\eta_1}$-extension has stage-respecting names in this normal form.

For the semantic description, let $K*G_\beta$ be $\forceP_\beta$-generic over
$W$.  Let
\[
   \mathfrak a_\beta=(E_\beta,A_\beta)
   :=\operatorname{supp}_{\mathrm{AP}}
      ((\dot x,\dot y,\check m),\forceP_\beta),
\]
and put
\[
   G_{\mathfrak a_\beta}:=K_{E_\beta}*G_{A_\beta},
   \qquad
   M_{\mathfrak a_\beta}:=W[K_{E_\beta}][G_{A_\beta}].
\]
Set
\[
   x=\dot x^{G_{\mathfrak a_\beta}},
   \qquad
   y=\dot y^{G_{\mathfrak a_\beta}}.
\]
Assume that
\[
   W[K][G_\beta]\models(x,y)\in A_m.
\]
We distinguish two cases.

\begin{enumerate}
\item[(1)] \emph{The first uniformization case.}
There is a least ordinal $\zeta<\alpha^\ast$ for which some real
\[
   y_0\in M_{\mathfrak a_\beta}\cap2^\omega
\]
satisfies
\[
   W[K][G_\beta]\models(x,y_0)\in A_m,
\]
and no $\zeta$-allowable complete reservoir extension of $\forceP_\beta$
forces $(x,y_0)$ out of $A_m$.  Among all such values, choose $y_0$ so that
\[
   \key_{\mathfrak a_\beta}(\langle x,y_0,m\rangle)
\]
is $<_L$-least.  Let $\dot y_0$ be the $<_L$-least stage-respecting
$\forceP_{\mathfrak a_\beta}$-name for $y_0$, lifted to $\forceP_\beta$.
This name is used only to form the auxiliary name; the choice of $y_0$ was
made by the semantic key.

In this case $y_0$ is declared to be a tentative value for $f_m(x)$ of rank
$\zeta$.  If $y\neq y_0$, let
\[
   z_\beta^*:=y.
\]
If $y=y_0$, let $z_\beta^*$ be the real $w\neq y_0$ in
$M_{\mathfrak a_\beta}$ for which
\[
   \key_{\mathfrak a_\beta}(\langle x,w,m\rangle)
\]
is $<_L$-least.  Let $\dot z_\beta^*$ be the $<_L$-least stage-respecting
$\forceP_{\mathfrak a_\beta}$-name for $z_\beta^*$, lifted to
$\forceP_\beta$.

We let
\[
   \forceP(\beta)
   :=\mathbb A_D\bigl(
      \dot Y(\dot x,\dot z_\beta^*,\check m,\dot g_{\eta_\beta})
   \bigr)
\]
be the corresponding coding forcing over the fixed reservoir and put
\[
   \forceP_{\beta+1}=\forceP_\beta*\forceP(\beta).
\]
The auxiliary name $\dot I_{\beta+1}$ is the canonical name forced to satisfy
\[
   \dot I_{\beta+1}^{\dot G_{\beta+1}}
   =
   \dot I_\beta^{\dot G_\beta}
   \cup
   \{(\dot x^{\dot G_\beta},
       \dot y_0^{\dot G_\beta},
       \check m,\check\zeta)\}.
\]
Thus the tentative value $y_0$ is not coded, while
$(x,z_\beta^*,m)$ is coded.

\item[(2)] \emph{The second uniformization case.}
Suppose that no real in $M_{\mathfrak a_\beta}\cap2^\omega$ passes the test
at any rank $\zeta<\alpha^\ast$.  In particular, the current real $y$ can be
forced out of $A_m$ by an $\alpha^\ast$-allowable complete reservoir
continuation of the current iteration.  We choose the $<_L$-least such
continuation, restrict its quotient below a condition forcing
\[
   (x,y)\notin A_m,
\]
and append this quotient to the current iteration.  The auxiliary name is
continued by the auxiliary name supplied by this $\alpha^\ast$-allowable
continuation.  No new tentative value is selected at this stage.
\end{enumerate}


This defines a pair $(\forceP_{\omega_1},\dot I_{\omega_1})\in W$.  The bookkeeping can be recoded into a function $F'\in W$ such that, for every $\delta<\omega_1$, the initial pair $(\forceP_\delta,\dot I_\delta)$ is $\alpha^\ast$-allowable over $W$ relative to
\[
   (F'\upharpoonright\delta,\vec\eta\upharpoonright\delta).
\]
Here $\vec\eta\upharpoonright\delta\in L$, and every bookkeeping name used
before $\delta$ has reservoir support contained in the corresponding
$U_\beta=\beta$.

\begin{fact}
The pair $(\forceP_{\omega_1},\dot I_{\omega_1})$ belongs to $W$, and every countable initial segment is $\alpha^\ast$-allowable over $W$ relative to a fixed bookkeeping function $F'\in W$ and the corresponding initial segment of the canonical area assignment $\vec\eta\in L$.
\end{fact}

Let $G_{\omega_1}$ be $\forceP_{\omega_1}$-generic over $W$.  For every $\beta\leq\omega_1$, let $G_\beta=G_{\omega_1}\cap\forceP_\beta$ and put
\[
   I_\beta:=\dot I_\beta^{G_\beta},
   \qquad
   I:=I_{\omega_1}.
\]

\begin{lemma}
\label{lemma-stability-of-eliminability}
Let $\beta\leq\gamma<\omega_1$, and suppose that
$\forceP_\gamma$ is an initial segment of the final iteration extending
$\forceP_\beta$.  Let $G_\gamma$ be $\forceP_\gamma$-generic over $W$, and
let
\[
   x,y\in W[G_\beta]\cap2^\omega,
   \qquad m<\omega,
   \qquad \xi\in\mathrm{Ord}.
\]
If $y$ is not $\xi$-persistent for $(x,m)$ over $\forceP_\beta$, then $y$ is
not $\xi$-persistent for $(x,m)$ over $\forceP_\gamma$.  Equivalently, if
$y$ is $\xi$-persistent for $(x,m)$ over $\forceP_\gamma$, then $y$ was
already $\xi$-persistent for $(x,m)$ over $\forceP_\beta$.
\end{lemma}

\begin{proof}
Suppose that $y$ is not $\xi$-persistent over $\forceP_\beta$.  After
restricting below a condition, there is a $\xi$-allowable complete reservoir
extension
\[
   \forceQ\succeq_{\mathrm{res}}\forceP_\beta
\]
such that, over the chosen copy of $\forceP_\beta$,
\[
   \forceQ/G_\beta\Vdash (x,y)\notin A_m .
\]
Since $\forceP_\gamma$ is an initial segment of the final
$\infty$-allowable iteration, it is $\alpha^\ast$-allowable.  If
$\xi<\alpha^\ast$, Lemma~\ref{shrinkinglemma} implies that
$\forceP_\gamma$ is $\xi$-allowable.  If $\xi\geq\alpha^\ast$, then
\[
   \mathcal A^W_\xi=\mathcal A^W_{\alpha^\ast}
\]
by the stabilization of the allowability hierarchy.  Hence, in all cases,
$\forceP_\gamma$ is $\xi$-allowable.  There is a $\forceP_\beta$-name
$\dot{\forceR}_{\beta,\gamma}$ such that
\[
   \forceP_\gamma
   \cong
   \forceP_\beta*\dot{\forceR}_{\beta,\gamma}.
\]
After identifying $\forceP_\beta$ with its chosen complete reservoir copy in
$\forceQ$, there is also a $\forceP_\beta$-name $\dot{\mathbb T}_Q$ such that
\[
   \forceQ
   \cong
   \forceP_\beta*\dot{\mathbb T}_Q,
\]
and, whenever $G_\beta$ is $\forceP_\beta$-generic,
\[
   \dot{\forceR}_{\beta,\gamma}^{G_\beta}
      \cong\forceP_\gamma/G_\beta,
   \qquad
   \dot{\mathbb T}_Q^{G_\beta}
      \cong\forceQ/G_\beta.
\]
Over the common initial forcing form the fibered product
\[
   \forceS_0
   :=
   \forceP_\beta*
   \bigl(
      \dot{\forceR}_{\beta,\gamma}
      \times
      \dot{\mathbb T}_Q
   \bigr).
\]
Apply the product construction from
Lemma~\ref{alphaallowableclosedunderproducts} relative to the common complete
reservoir subiteration $\forceP_\beta$: after moving the fresh reservoir
coordinates of the two tails to disjoint blocks, the amalgam has a normalized
$\xi$-allowable presentation.  Thus there are a normalized $\xi$-allowable
iteration $\forceS$ and an isomorphism
\[
   \pi:\forceS_0\cong\forceS.
\]

Under $\pi$, the canonical copy of $\forceP_\gamma$ in $\forceS_0$ becomes a
complete reservoir subiteration of $\forceS$.  The quotient of $\forceS$ over
this copy contains the translated copy of $\dot{\mathbb T}_Q$.  Since
\[
   \dot{\mathbb T}_Q^{G_\beta}
   \Vdash (x,y)\notin A_m
\]
and $(x,y)\notin A_m$ is a $\Sigma^1_3$ statement, the same witness remains
valid after adjoining the generic for
$\dot{\forceR}_{\beta,\gamma}^{G_\beta}$, by Shoenfield absoluteness applied
to the $\Pi^1_2$ matrix.  Consequently,
\[
   \forceS/G_\gamma\Vdash (x,y)\notin A_m.
\]
Thus $\forceS$ is a $\xi$-allowable complete reservoir extension of
$\forceP_\gamma$ witnessing that $y$ is not $\xi$-persistent over
$\forceP_\gamma$, as required.
\end{proof}

By Lemma~\ref{fvaluesremain2}, every tentative value for $f_m(x)$ of rank below $\alpha^\ast$ belongs to $A_m$ in the final model.  It remains to show that, for every $m\in\omega$ and every real $x$ with $A_{m,x}\neq\varnothing$, there is exactly one $y$ such that
\[
   (x,y)\in A_m
   \qquad\text{and}\qquad
   (x,y,m)\text{ is not coded into }\vec S.
\]
The next lemma gives the required dichotomy.

\begin{lemma}
In $W[G_{\omega_1}]$ the following dichotomy holds true:
\begin{enumerate}
\item The auxiliary set contains a tentative value for $(x,m)$, that is,
\[
   \exists y\in2^\omega\ \exists\xi<\alpha^\ast
   \bigl((x,y,m,\xi)\in I\bigr).
\]
In this case there is a unique real $y_0$ such that
\[
   W[G_{\omega_1}]\models
   (x,y_0)\in A_m
   \ \wedge\ 
   (x,y_0,m)\text{ is not coded into }\vec S.
\]

\item The auxiliary set contains no tentative value for $(x,m)$, that is,
\[
   \forall y\in2^\omega\ \forall\xi<\alpha^\ast
   \bigl((x,y,m,\xi)\notin I\bigr).
\]
In this case
\[
   W[G_{\omega_1}]\models A_{m,x}=\varnothing.
\]
\end{enumerate}
\end{lemma}
\begin{proof}
Assume first that the auxiliary set contains a tentative value for $(x,m)$.
Let $\xi_0<\alpha^\ast$ be the least rank for which some stage selects a
real $y$ for $(x,m)$ and enters $(x,y,m,\xi_0)$ into the auxiliary set.
Among all such stages of rank $\xi_0$, choose one for which the key used at
that stage is $<_L$-least.  Let $y_0$ be the value selected there, let
$\kappa_0$ be this key, and let $\beta$ be the least stage at which
$y_0$ is selected with rank $\xi_0$ and key $\kappa_0$.
\par \medskip

\begin{claim}
$W[G_{\omega_1}] \models (x,y_0) \in A_m$.
\end{claim}
\begin{proof}[Proof of the first claim]
This follows immediately from Lemma~\ref{fvaluesremain2}.
\end{proof}

\par \medskip

\begin{claim}[No reversal]
Let $\eta$ be a stage at which the local datum is again $(x,m)$ and whose
local admissible support is $\mathfrak a_\eta$.  Suppose that
$y_0\in M_{\mathfrak a_\eta}$ and that $y_0$ is $\xi_0$-persistent for
$(x,m)$ over $\forceP_\eta$.  Then the construction cannot select a value
$v\neq y_0$ for $(x,m)$ at stage $\eta$.
\end{claim}
\begin{proof}[Proof of the no-reversal claim]
Since $y_0$ passes the persistence test at rank $\xi_0$, the second
uniformization case does not apply at stage $\eta$.  Let $v$ be the value
selected in the first uniformization case and let $\rho$ be its rank.  Then
$\rho\leq\xi_0$.

If $\rho<\xi_0$, then $(x,v,m,\rho)$ is entered into the auxiliary set,
contradicting the minimality of $\xi_0$.  Hence $\rho=\xi_0$.
Assume that $v\neq y_0$.  At stage $\eta$, the real $y_0$ is one of the
rank-$\xi_0$ reals compared by key.  Therefore the selection of $v$ gives
\[
   \key_{\mathfrak a_\eta}(\langle x,v,m\rangle)
   <_L
   \key_{\mathfrak a_\eta}(\langle x,y_0,m\rangle).
\]
Passing to a common admissible closure of the support used at stage $\eta$
and the support used when $y_0$ was selected, Lemma~\ref{lemma-semantic-orbit-support-independence}
shows that
\[
   \key_{\mathfrak a_\eta}(\langle x,y_0,m\rangle)=\kappa_0.
\]
The same support-independence identifies the key of $v$ at stage $\eta$
with the key carried by the occurrence $(x,v,m,\xi_0)$ in the auxiliary
set.  Thus there is an occurrence of rank $\xi_0$ with key $<_L\kappa_0$,
contradicting the choice of $\kappa_0$.  Hence $v=y_0$.
\end{proof}

\par \medskip

\begin{claim}
\begin{align*}
 W[G_{\omega_1}] \models``y_0 & \text{ is the unique real $y$ such that }  \\& (x,y,m) \text{ is not coded somewhere in the $\vec{S}$-sequence.$"$ }
\end{align*}
\end{claim}
\begin{proof}[Proof of the second claim]
We first show that $(x,y_0,m)$ is not coded into $\vec S$.  Suppose towards a
contradiction that stage $\eta$ codes $(x,y_0,m)$.  Then $y_0$ belongs to
the local model used at that stage.

If $\eta<\beta$, then $y_0$ is $\xi_0$-persistent over $\forceP_\eta$.
Indeed, otherwise Lemma~\ref{lemma-stability-of-eliminability} would imply
that $y_0$ is not $\xi_0$-persistent over $\forceP_\beta$, contradicting
its selection at stage $\beta$.  If $\eta\geq\beta$, then
$(x,y_0,m,\xi_0)$ is already in the auxiliary set of the initial segment
$\forceP_\eta$, and Lemma~\ref{fvaluesremain2} gives that $y_0$ is
$\xi_0$-persistent over $\forceP_\eta$.

Thus the no-reversal claim applies at stage $\eta$.  The stage cannot be in
the second uniformization case, and in the first uniformization case it must
select $y_0$.  But a stage which selects $y_0$ does not code $(x,y_0,m)$: if
the bookkeeping value is $y_0$, the definition uses the alternate value
$z_\eta^*\neq y_0$, and otherwise it codes the bookkeeping value.  This is a
contradiction.  Hence $(x,y_0,m)$ is not coded into $\vec S$.

We now prove uniqueness.  Let $y\neq y_0$ and suppose that
\[
   W[G_{\omega_1}]\models (x,y)\in A_m.
\]
Choose a stage $\eta\geq\beta$ at which the bookkeeping presents
stage-respecting names for $x$, $y$, and $m$, and such that the local
admissible support also contains a support for $y_0$.  By the downward form of
Shoenfield absoluteness applied to the $\Pi^1_3$ definition of $A_m$, the
model $W[G_\eta]$ already satisfies $(x,y)\in A_m$.  Since
$(x,y_0,m,\xi_0)$ belongs to the auxiliary set of the initial segment
$\forceP_\eta$, Lemma~\ref{fvaluesremain2} shows that $y_0$ is
$\xi_0$-persistent over $\forceP_\eta$.  The no-reversal claim therefore
implies that the stage selects $y_0$.  Since the bookkeeping value at this
stage is $y\neq y_0$, the construction codes $(x,y,m)$ into $\vec S$.

This proves that $y_0$ is the unique real $y$ such that $(x,y)\in A_m$ and
$(x,y,m)$ is not coded into $\vec S$, and finishes the proof in the first case
of the dichotomy.
\end{proof}


\medskip
We now assume that the auxiliary set contains no tentative value for $(x,m)$.  We must show that
$W[G_{\omega_1}] \models `` \text{The $x$-section of $A_m$ is empty}".$

Under this assumption, whenever we are at a stage $\beta$ such that there is a $y$ such that $F(\beta)=(x,y,m)$, then the second uniformization case in the definition of $\forceP_{\omega_1}$ must apply.  For every such $y$, the stage $\beta$ construction appends an $\infty$-allowable eliminating tail which forces
\[
   (x,y)\notin A_m.
\]
By Lemma~\ref{lemma-stability-of-eliminability}, this eliminability is stable under all later $\infty$-allowable continuations.  Since $(x,y)\notin A_m$ is $\Sigma^1_3$, upward absoluteness, equivalently Shoenfield absoluteness for its $\Pi^1_2$ matrix, gives
\[
   W[G_{\omega_1}]
   \models \lnot\exists y\,((x,y)\in A_m).
\]
This finishes the proof of the lemma.


\end{proof}
\begin{corollary}
In $W[G_{\omega_1}]$, the $\Pi^1_3$-uniformization property holds.  More
precisely, for every $m\in\omega$, define a relation $U_m$ on
$2^\omega\times2^\omega$ by
\[
   U_m(x,y)
   \quad\Longleftrightarrow\quad
   (x,y)\in A_m
   \ \wedge\
   \neg\Phi\bigl(\langle x,y,m\rangle\bigr).
\]
Then $U_m$ is the graph of a $\Pi^1_3$ uniformizing function for $A_m$.
\end{corollary}

\begin{proof}
The relation $(x,y)\in A_m$ is $\Pi^1_3$, while
$\Phi(\langle x,y,m\rangle)$ is $\Sigma^1_3$.  Hence $U_m$ is
$\Pi^1_3$.  By the preceding dichotomy, whenever the section $A_{m,x}$ is
nonempty, there is exactly one $y\in A_{m,x}$ for which
$\Phi(\langle x,y,m\rangle)$ fails; if the section is empty, there is no such
$y$.  Therefore $U_m$ is the graph of a function with domain
$\operatorname{pr}_1(A_m)$, and this graph is contained in $A_m$.
\end{proof}



\section[Forcing over the canonical inner models Mn]{Forcing over the canonical inner models $M_n$}
\label{section-forcing-over-Mn}

Fix for the remainder of this section a natural number $n$ with
$1\leq n<\omega$, and assume that  $M_n^{\#}$ exists. As always we let $M_n$ be the resulting inner model one obtains when iterating the sharp out of the universe.  Our goal is to prove the following theorem.

\begin{theorem*}
There is a set-generic extension of $M_n$ in which the
$\Pi^1_{n+3}$-uniformization property holds.
\end{theorem*}

The proof follows the construction over $L$ from the preceding sections, but
three ingredients require $M_n$-analogues: the recovery of the canonical lower
part $M_n|\omega_1$, the definition of the coding predicate, and
the absoluteness argument which keeps tentative uniformizing values in their
projective sections.  We begin with the inner-model-theoretic background needed
for these analogues.

\subsection{Inner-model-theoretic background and generic absoluteness}
\label{subsection-Mn-background-and-absoluteness}

We use the notation and comparison conventions of Steel's analysis of mice
with finitely many Woodin cardinals; see \cite{Steel2,Steel3}.  An
\emph{ordinary premouse} means a premouse in the parameter-free premouse
language.  In particular, no real or other base set is included as a predicate.
Every proper initial segment of $M_n$ is $n$-small and $\omega$-sound, and the
background construction of $M_n$ supplies the realizability strategy used on
such an initial segment in a comparison.

\begin{definition}
\label{definition-n-smallness}
Let $\mathcal M$ be a premouse and let
$\eta<\operatorname{Ord}^{\mathcal M}$.  We say that $\mathcal M$ is
\emph{$n$-small above $\eta$} if, whenever $E$ is an extender on the
$\mathcal M$-sequence and
\[
   \eta<\operatorname{crit}(E),
\]
the initial segment
$\mathcal J^{\mathcal M}_{\operatorname{crit}(E)}$ does not have $n$ Woodin
cardinals above $\eta$.  We say that $\mathcal M$ is \emph{$n$-small} if it
is $n$-small above $0$.
\end{definition}

The weak iterability notion used at the finite projective levels is Steel's
$\Pi_n$-iterability.  Its recursive definition depends on the parity of $n$.
At even levels one uses the corresponding finite weak iteration game, while at
odd levels the last round includes the appropriate goodness requirement for
the branch selected by player~II.  We use the standard game from
\cite{Steel2,Steel3} and isolate the notation needed below.

\begin{definition}
\label{definition-Steel-Pin-iterability}
Let $\mathcal M$ be a countable premouse and let
$\eta<\operatorname{Ord}^{\mathcal M}$ be a cutpoint.  We say that
$\mathcal M$ is \emph{$\Pi_n$-iterable above $\eta$} if player~II has a
winning strategy in Steel's game
\[
   \mathcal G(\mathcal M,\eta,n)
\]
for $\Pi_n$-iterability above $\eta$.  We say simply that $\mathcal M$ is
\emph{$\Pi_n$-iterable} if it is $\Pi_n$-iterable above $0$.
\end{definition}
Steel's analysis gives the following complexity statement.

\begin{fact}
\label{fact-projective-complexity-Pin-iterability}
For each fixed $1\leq n<\omega$, the relation
\[
   ``x\text{ codes a countable premouse which is }
   \Pi_n\text{-iterable}"
\]
is $\Pi^1_{n+1}$ in the real code $x$.  Moreover, every countable realizable
$n$-small premouse is $\Pi_n$-iterable.
\end{fact}

We shall use $\Pi_n$-iterability only through the following comparison
consequence.

\begin{lemma}
\label{lemma-comparison-with-Pin-iterable-premouse}
Let $\mathcal M$ and $\mathcal N$ be countable ordinary premice.  Assume that
both are $\omega$-sound and satisfy
\[
   \rho_\omega(\mathcal M)=\rho_\omega(\mathcal N)=\omega,
\]
that $\mathcal M\triangleleft M_n$, and that $\mathcal N$ is $n$-small and
$\Pi_n$-iterable.  Then the comparison of $\mathcal M$ with $\mathcal N$ is
successful.  Consequently,
\[
   \mathcal M\trianglelefteq\mathcal N
   \quad\text{or}\quad
   \mathcal N\trianglelefteq\mathcal M.
\]
\end{lemma}

\begin{proof}
This is the parameter-free, lower-part case of Steel's comparison theorem for
$n$-small mice.  Run the usual least-disagreement coiteration.  The
$\mathcal M$-side is governed by the realization strategy inherited from the
background construction of $M_n$, while the $\mathcal N$-side is governed by
the strategy witnessing $\Pi_n$-iterability.  Steel's branch uniqueness and
$\mathcal Q$-structure analysis show that the required branches at limit
stages are well-founded and that the coiteration terminates.  Since both
premice are $\omega$-sound and project to $\omega$, the standard no-drop
pullback argument gives the displayed initial-segment relation between the
original premice.  See \cite{Steel2,Steel3}.
\end{proof}

We next specify the generic extensions in which the comparison lemma will be
used.  All cardinalities in the following definition are computed in $M_n$.

\begin{definition}
\label{definition-small-Mn-generic-extension}
A model $N$ is a \emph{small $M_n$-generic extension} if
\[
   N=M_n[G]
\]
for some cardinal-preserving forcing $\forceP\in M_n$ such that
$M_n\models|\forceP|\leq\aleph_1$ and some $M_n$-generic filter
$G\subseteq\forceP$.  A \emph{small two-step extension} is an extension
\[
   M_n[G]\subseteq M_n[G][H]
\]
obtained from a two-step forcing
$\forceP*\dot{\forceQ}\in M_n$ which is forced to be cardinal preserving and
of size at most $\aleph_1$.
\end{definition}

The preparatory branch forcing and every forcing iteration used later will be
shown to have size $\aleph_1$ and to preserve cardinals.  Hence all models
which occur in the proof are small $M_n$-generic extensions in the sense of
Definition~\ref{definition-small-Mn-generic-extension}.

\begin{lemma}
\label{lemma-preservation-Steel-Pin-iterability}
Let $N$ be a small $M_n$-generic extension.  If
$\mathcal M\triangleleft M_n$ is countable in $N$, then
\[
   N\models ``\mathcal M\text{ is }\Pi_n\text{-iterable}."
\]
\end{lemma}

\begin{proof}
The background construction of $M_n$ realizes every proper initial segment of
$M_n$.  The standard preservation argument for Steel's weak iteration games
shows that this realization continues to supply the required branches in the
small, $\omega_1$-preserving generic extensions considered here.  Thus the
$\Pi_n$-iterability of a genuine initial segment of $M_n$ is preserved.  We
will use only this preservation statement; we do not claim that the class of
all real codes for $\Pi_n$-iterable premice is absolute between arbitrary
outer models with the same $\omega_1$.
\end{proof}

The second background ingredient is the mouse-operator closure which yields
projective generic absoluteness.  Recall that a set is
\emph{self-wellordered}, abbreviated \emph{swo}, if it codes a wellorder of
itself.  For a self-wellordered set $X$, let $M_k^{\#}(X)$ denote the canonical
active $X$-mouse with $k$ Woodin cardinals.  Thus $M_0^{\#}(X)$ is the usual
sharp over $X$.

\begin{fact}
\label{fact-Mn-closure-under-lower-mouse-operator}
The model $M_n$ is closed under the operator
\[
   X\longmapsto M_{n-1}^{\#}(X)
\]
below its least Woodin cardinal.  In particular, for every self-wellordered
$X\in H(\omega_2)^{M_n}$, the mouse $M_{n-1}^{\#}(X)$ exists and is normally
$\omega_2$-iterable in $M_n$.
\end{fact}

This closure is preserved by the small forcing notions which occur in the
construction.

\begin{lemma}
\label{lemma-small-forcing-preserves-lower-mouse-operator}
Let $N=M_n[G]$ be a small $M_n$-generic extension.  Then
\[
\begin{aligned}
   N\models{}&``M_{n-1}^{\#}(X)\text{ exists and is normally }
      \omega_2\text{-iterable}\\
   &\text{for every swo }X\in H(\omega_2)."
\end{aligned}
\]
Moreover, whenever $X$ belongs to both $M_n$ and $N$, the two models compute
the same mouse $M_{n-1}^{\#}(X)$.
\end{lemma}

\begin{proof}
Apply the generic preservation theorem for mouse operators to
$M_{n-1}^{\#}$.  In the notation of \cite[Lemma~3.4]{Schlicht}, take
$\kappa=\omega_1$ and replace the index occurring there by $n-1$.
Fact~\ref{fact-Mn-closure-under-lower-mouse-operator} supplies the hypothesis
that $M_{n-1}^{\#}(X)$ exists for every self-wellordered
$X\in H(\omega_2)^{M_n}$.  The cited lemma gives existence, normal
$\omega_2$-iterability, and correctness of the operator in the forcing
extension.
\end{proof}

We can now state the precise generic absoluteness principle used later.  It is
the two-step form which permits real parameters introduced by the first
forcing.

\begin{theorem}
\label{theorem-small-two-step-projective-absoluteness}
Let
\[
   M_n[G]\subseteq M_n[G][H]
\]
be a small two-step extension.  Then
\[
   M_n[G]\prec_{\Sigma^1_{n+2}} M_n[G][H].
\]
Equivalently, for every real parameter $a\in M_n[G]$ and every
$\Sigma^1_{n+2}$ formula $\vartheta(v)$,
\[
   M_n[G]\models\vartheta(a)
   \quad\Longleftrightarrow\quad
   M_n[G][H]\models\vartheta(a).
\]
The same equivalence holds for $\Pi^1_{n+2}$ formulas.
\end{theorem}

\begin{proof}
By Lemma~\ref{lemma-small-forcing-preserves-lower-mouse-operator}, the
intermediate model $M_n[G]$ is closed under
$X\mapsto M_{n-1}^{\#}(X)$ for self-wellordered
$X\in H(\omega_2)^{M_n[G]}$.  Apply the Martin--Solovay--Schindler generic
absoluteness theorem in the form of \cite[Lemma~3.7]{Schlicht}, with the mouse
index $n-1$ and $\kappa=\omega_1$, to the second forcing.  This yields
$\Sigma^1_{n+2}$-elementarity between $M_n[G]$ and $M_n[G][H]$.  The
$\Pi^1_{n+2}$ assertion follows by taking complements.
\end{proof}

The point at which Theorem~\ref{theorem-small-two-step-projective-absoluteness}
enters the uniformization proof is recorded in the following consequence.

\begin{corollary}
\label{corollary-persistence-Sigma-n-plus-3-witnesses}
Let
\[
   M_n[G]\subseteq M_n[G][H]
\]
be a small two-step extension, let $a\in M_n[G]$ be a real, and let
$\varphi(v)$ be a $\Sigma^1_{n+3}$ formula.  If
\[
   M_n[G]\models\varphi(a),
\]
then
\[
   M_n[G][H]\models\varphi(a).
\]
Consequently, if $A\subseteq(2^\omega)^2$ is $\Pi^1_{n+3}$ and
\[
   M_n[G]\models (x,y)\notin A,
\]
then $(x,y)\notin A$ remains true in $M_n[G][H]$.
\end{corollary}

\begin{proof}
Write
\[
   \varphi(a)\equiv\exists z\,\psi(z,a),
\]
where $\psi$ is $\Pi^1_{n+2}$.  Choose a witness $z\in M_n[G]$.  By
Theorem~\ref{theorem-small-two-step-projective-absoluteness}, the statement
$\psi(z,a)$ remains true in $M_n[G][H]$, so the same witness establishes
$\varphi(a)$ there.  The final assertion is the special case in which
$\varphi(x,y)$ defines the complement of $A$.
\end{proof}


\subsection[The definable lower part of Mn]{The definable lower part of $M_n$}
\label{subsection-definable-lower-part-Mn}

For the coding construction we do not need to recover the whole proper class
model $M_n$.  It is enough to recover the canonical initial segment
$M_n|\omega_1$ uniformly in the generic extensions which occur in
the proof.  To this end we isolate a projectively definable cofinal system of
countable lower-part initial segments of $M_n$.

The lower-part requirement in the next definition is important.  At the
higher projective levels there can be nonstandard $\Pi_n$-iterable premice
which do not occur on the $M_n$-sequence.  Since all ordinals below
$\omega_1^{M_n}$ lie below the least Woodin cardinal of $M_n$, the genuine
initial segments relevant to the present coding argument have no Woodin
cardinal.  Restricting to such premice removes the unwanted objects.

\begin{definition}
\label{definition-Mn-lower-part-approximations}
Let $N=M_n[G]$ be a small $M_n$-generic extension in the sense of
Definition~\ref{definition-small-Mn-generic-extension}.  In $N$, let
$\mathcal I_n^N$ be the class of all countable premice $\mathcal M$ such that
\begin{enumerate}
   \item $\mathcal M$ is a passive ordinary premouse;
   \item $\mathcal M$ is a lower-part premouse, that is,
         $\mathcal M$ has no Woodin cardinal;
   \item $\mathcal M$ is $n$-small and $\omega$-sound, and
         \[
            \rho_\omega(\mathcal M)=\omega;
         \]
   \item $\mathcal M$ is $\Pi_n$-iterable in the sense of
         Definition~\ref{definition-Steel-Pin-iterability}.
\end{enumerate}
We identify a real coding such a premouse with its transitive collapse.  When
there is no danger of confusion about the ambient extension, we simply write
$\mathcal I_n$.
\end{definition}

\begin{lemma}
\label{lemma-definable-Mn-initial-segments}
Let $N=M_n[G]$ be a small $M_n$-generic extension.  Then the set of real codes
for members of $\mathcal I_n^N$ is $\boldsymbol{\Pi}^1_{n+1}$-definable,
uniformly in $N$.  Moreover, every member of $\mathcal I_n^N$ is of the form
\[
   \mathcal J^{M_n}_\eta
\]
for some $\eta<\omega_1$, and the set
\[
   C_n:=\{\eta<\omega_1:
      \mathcal J^{M_n}_\eta\in\mathcal I_n^N\}
\]
is cofinal in $\omega_1$.
\end{lemma}

\begin{proof}
The assertion that a real codes a well-founded countable passive ordinary
premouse has complexity at most $\Pi^1_1$.  Once the premouse has been
decoded, the requirements that it be lower-part, $n$-small,
$\omega$-sound, and have $\rho_\omega=\omega$ are given by the usual
fine-structural conditions and do not increase the complexity beyond
$\Pi^1_1$.  By
Fact~\ref{fact-projective-complexity-Pin-iterability}, the remaining
iterability requirement is $\Pi^1_{n+1}$.  Since $n\geq1$, the conjunction is
$\boldsymbol{\Pi}^1_{n+1}$.  The same formula is used in every small
$M_n$-generic extension, which gives uniformity.

We next show that no nonstandard lower-part premouse belongs to
$\mathcal I_n^N$.  Let $\mathcal M\in\mathcal I_n^N$.  Since $\mathcal M$ is
countable in $N$ and the forcing from $M_n$ to $N$ preserves $\omega_1$, we
have
\[
   \operatorname{Ord}^{\mathcal M}<\omega_1^{M_n}.
\]
Choose an ordinal $\eta$ such that
\[
   \operatorname{Ord}^{\mathcal M}<\eta<\omega_1
\]
and such that $\mathcal J^{M_n}_\eta$ is passive, lower-part,
$\omega$-sound, and projects to $\omega$.  Standard fine structure gives
cofinally many such ordinals $\eta<\omega_1$.  The premouse
$\mathcal J^{M_n}_\eta$ is realizable in the background construction of
$M_n$, and it remains $\Pi_n$-iterable in $N$ by
Lemma~\ref{lemma-preservation-Steel-Pin-iterability}.

Apply Lemma~\ref{lemma-comparison-with-Pin-iterable-premouse} to
$\mathcal J^{M_n}_\eta$ and $\mathcal M$.  We obtain
\[
   \mathcal J^{M_n}_\eta\trianglelefteq\mathcal M
   \quad\text{or}\quad
   \mathcal M\trianglelefteq\mathcal J^{M_n}_\eta.
\]
The first alternative is impossible because
$\operatorname{Ord}^{\mathcal M}<\eta$.  Hence
\[
   \mathcal M\trianglelefteq\mathcal J^{M_n}_\eta,
\]
so $\mathcal M$ is itself an initial segment of $M_n$.  Thus
$\mathcal M=\mathcal J^{M_n}_\xi$ for some $\xi<\omega_1$.

Conversely, suppose that $\eta<\omega_1$ and that
$\mathcal J^{M_n}_\eta$ is passive, lower-part, $\omega$-sound, and projects
to $\omega$.  It is $n$-small and realizable in the construction of $M_n$.
By Lemma~\ref{lemma-preservation-Steel-Pin-iterability}, it is
$\Pi_n$-iterable in $N$, and hence
\[
   \mathcal J^{M_n}_\eta\in\mathcal I_n^N.
\]
Since the ordinals $\eta$ with these fine-structural properties are cofinal in
$\omega_1$, the set $C_n$ is cofinal in $\omega_1$.
\end{proof}

We can now recover $M_n|\omega_1$ without referring to the proper
class $M_n$ in the defining formula.  In a small generic extension $N$, let
$\Theta_n(X)$ be the following assertion:
\begin{align*}
\Theta_n(X)\ \Longleftrightarrow\ {}&X\text{ is transitive and }
   \operatorname{Ord}\cap X=\omega_1,\\
&\text{and for every }z,\\[-2mm]
&\qquad z\in X
   \ \Longleftrightarrow\
   \exists\mathcal M\,
      \bigl(\mathcal M\in\mathcal I_n^N\ \wedge\ z\in\mathcal M\bigr).
\end{align*}
Here the quantifier over $\mathcal M$ may equivalently be read as a quantifier
over real codes for countable premice satisfying
Definition~\ref{definition-Mn-lower-part-approximations}.

\begin{definition}
\label{definition-recover-Mn-omega1}
Let $N=M_n[G]$ be a small $M_n$-generic extension.  We say in $N$ that
\[
   X=M_n|\omega_1
\]
if $\Theta_n(X)$ holds.  Equivalently,
\[
   X=\bigcup\mathcal I_n^N.
\]
\end{definition}

\begin{lemma}
\label{lemma-correctness-recovery-Mn-omega1}
Let $N=M_n[G]$ be a small $M_n$-generic extension.  Then there is a unique
$X$ such that $N\models\Theta_n(X)$, and this set is the true initial segment
\[
   X=M_n|\omega_1
    =\mathcal J^{M_n}_{\omega_1}.
\]
Consequently, the same formula $\Theta_n$ recovers the same set
$M_n|\omega_1$ in every small $M_n$-generic extension.
\end{lemma}

\begin{proof}
By Lemma~\ref{lemma-definable-Mn-initial-segments}, every member of
$\mathcal I_n^N$ is an initial segment
$\mathcal J^{M_n}_\eta$ for some $\eta<\omega_1$.  Therefore
\[
   \bigcup\mathcal I_n^N
   \subseteq\mathcal J^{M_n}_{\omega_1}.
\]
The same lemma states that the set of $\eta<\omega_1$ for which
$\mathcal J^{M_n}_\eta\in\mathcal I_n^N$ is cofinal in $\omega_1$.  Hence
\[
   \mathcal J^{M_n}_{\omega_1}
   \subseteq\bigcup\mathcal I_n^N.
\]
Thus
\[
   \bigcup\mathcal I_n^N=\mathcal J^{M_n}_{\omega_1}.
\]
This union is transitive and has ordinal height $\omega_1$, so it satisfies
$\Theta_n$.  Conversely, the last clause of $\Theta_n$ uniquely determines
$X$ to be this union.  Since the argument applies in every small generic
extension, the recovery is uniform.
\end{proof}

\begin{corollary}
\label{corollary-projective-recovery-reals-Mn-omega1}
In every small $M_n$-generic extension and for every real $x$,
\[
   x\in M_n|\omega_1
   \quad\Longleftrightarrow\quad
   \exists\mathcal M\in\mathcal I_n\ (x\in\mathcal M).
\]
In particular, the set of reals belonging to
$M_n|\omega_1$ has a uniform $\boldsymbol{\Sigma}^1_{n+2}$
definition in these extensions.
\end{corollary}

\begin{proof}
The equivalence follows immediately from
Lemma~\ref{lemma-correctness-recovery-Mn-omega1}.  In terms of real codes, the
right-hand side begins with an existential real quantifier and then asserts a
$\boldsymbol{\Pi}^1_{n+1}$ property of the code, together with the arithmetic
assertion that the decoded premouse contains $x$.  It is therefore
$\boldsymbol{\Sigma}^1_{n+2}$.
\end{proof}

\paragraph{Local use of the recovery formula.}
The uniformity asserted above concerns the actual outer models of $M_n$ used
in the construction.  It does not say that an arbitrary countable transitive
model which internally evaluates the formula $\Theta_n$ must recover the true
$M_n|\omega_1$.  Nor do we claim that the collection of real codes
belonging to $\mathcal I_n$ is literally the same in every outer model.  New
reals may provide new codes and new putative iteration trees.  What is
absolute for the extensions under consideration is the union of the premice
which pass the test: it is always the true
$\mathcal J^{M_n}_{\omega_1}$.

For this reason, the localized coding predicate introduced below will not ask
a countable auxiliary model to reconstruct the global union on its own.
Instead, it will be supplied with a countable premouse $m$ which, from the
outside, belongs to $\mathcal I_n$ and satisfies
\[
   \omega_1^m=\omega_1^M,
\]
where $M$ is the auxiliary countable transitive model.  The premouse $m$ then
provides the correct local approximation to all canonical objects constructed
from $M_n|\omega_1$.

\subsection[Canonical parameters from Mn restricted to omega1]{Canonical parameters from $M_n|\omega_1$}
\label{subsection-canonical-parameters-Mn}

We next isolate the canonical objects from the lower part of $M_n$ which will
be used by the coding construction.  They are the $M_n$-analogue of the
canonical constructible parameters used in Sections~2 and~3: a diamond
sequence, an independent sequence of Suslin trees, an almost disjoint family,
and a map which turns proper initial segments of an $M_n$-Cohen subset of
$\omega_1$ into coding areas.  The point of the present subsection is not only
that these objects exist in $M_n$, but that they are defined uniformly from the
lower part recovered in Subsection~\ref{subsection-definable-lower-part-Mn}.
Consequently the same definitions are available in every small
$M_n$-generic extension occurring below.

We first record the canonical order which will be used to make the choices in
all the ensuing recursions.

\begin{fact}
\label{fact-canonical-construction-order-Mn}
The order of construction of $M_n$ induces a canonical well-order
$<_{M_n}$ of $M_n$.  Its restriction to the reals of $M_n$ is a
$\Delta^1_{n+2}$ well-order.  Moreover, the restriction of this order to
$M_n|\omega_1$ is uniformly definable over the structure
$M_n|\omega_1$.
\end{fact}

Every real of $M_n$ belongs to $M_n|\omega_1$.  Hence, whenever a
subset of a countable ordinal or a countable structure is to be compared using
$<_{M_n}$, we use the fixed canonical real code of that object and compare
those codes.  We shall use $<_{M_n}$ only to make canonical choices among
objects already belonging to $M_n$; it will never be used to choose a member
of a projective section in the final generic extension.

The proof of the diamond principle uses the following form of condensation.
We state precisely the consequence needed for the countable hulls appearing
below.

\begin{theorem}
\label{theorem-condensation-lower-part-hulls-Mn}
Let $\mathcal P\trianglelefteq M_n$ be an $\omega$-sound initial segment and
let
\[
   \pi:\bar{\mathcal N}\longrightarrow\mathcal P
\]
be the inverse of the transitive collapse of a sufficiently elementary
substructure of $\mathcal P$.  Suppose that
\[
   \operatorname{crit}(\pi)
   =\rho_\omega(\bar{\mathcal N})
   =\omega_1^{\bar{\mathcal N}}.
\]
Then one of the following alternatives holds.
\begin{enumerate}
   \item $\bar{\mathcal N}\trianglelefteq M_n$.
   \item There are an initial segment $\mathcal Q\trianglelefteq M_n$ and an
         extender $E$ on the $M_n$-sequence such that
         \[
            \operatorname{lh}(E)=\omega_1^{\bar{\mathcal N}}
         \]
         and $\bar{\mathcal N}$ is an initial segment of the degree-zero
         ultrapower $\operatorname{Ult}_0(\mathcal Q,E)$.
\end{enumerate}
For the hulls used in the proof of
Lemma~\ref{lemma-canonical-Mn-diamond-sequence}, the second alternative is
impossible.  Hence their transitive collapses are initial segments of $M_n$.
\end{theorem}

\begin{proof}
The dichotomy is the corresponding instance of Steel's condensation theorem
for the fine-structural construction of $M_n$; see \cite{Steel3}.  We explain
why the ultrapower alternative is excluded in the application below.  The
hulls are chosen so that the standard projectum of the collapse is its
internal first uncountable cardinal
\[
   \lambda=\omega_1^{\bar{\mathcal N}}.
\]
The lower part $\bar{\mathcal N}\restriction\lambda$ therefore computes
$\lambda$ as its first uncountable cardinal.  If the second alternative held,
then an extender on the relevant $M_n$-sequence would be indexed at
$\lambda$.  In the corresponding extender model the index of such an
extender is a cardinal with the closure properties supplied by the extender;
in particular it cannot simultaneously be computed by the lower part as its
first uncountable cardinal.  This is the standard contradiction which rules
out the ultrapower alternative in Jensen's diamond argument.  Thus only the
initial-segment alternative remains.
\end{proof}

We can now carry out Jensen's canonical diamond construction inside $M_n$.

\begin{lemma}
\label{lemma-canonical-Mn-diamond-sequence}
There is a sequence
\[
   \vec D^{\,n}
   =\langle D^n_\alpha\mid\alpha<\omega_1\rangle\in M_n
\]
such that $D^n_\alpha\subseteq\alpha$ for every $\alpha<\omega_1$ and
\[
   M_n\models
   ``\vec D^{\,n}\text{ is a }\diamondsuit_{\omega_1}\text{-sequence}."
\]
The sequence is obtained by a recursion which is uniform over
$M_n|\omega_1$.
\end{lemma}

\begin{proof}
Work in $M_n$.  We recursively define
$\langle D^n_\beta\mid\beta<\alpha\rangle$ for
$\alpha<\omega_1$.  At successor stages put $D^n_\alpha=\varnothing$.
Suppose that $\alpha$ is a nonzero limit ordinal and the sequence below
$\alpha$ has already been defined.  Call a pair $(A,C)$ a
\emph{counterexample at $\alpha$} if
\[
   A\subseteq\alpha,
   \qquad C\subseteq\alpha\text{ is club in }\alpha,
\]
and
\[
   \forall\beta\in C\,
      \bigl(D^n_\beta\neq A\cap\beta\bigr).
\]
If such a pair exists in $M_n$, let $(A_\alpha,C_\alpha)$ be the pair whose
canonical real code is $<_{M_n}$-least, and put
\[
   D^n_\alpha=A_\alpha.
\]
If no such pair exists, put $D^n_\alpha=\varnothing$.

The recursion is uniform over $M_n|\omega_1$.  Indeed, every
object considered at stage $\alpha$ is coded by a real of $M_n$, and all
reals of $M_n$, together with the restriction of the construction order to
those reals, belong to or are uniformly definable over
$M_n|\omega_1$.

Suppose toward a contradiction that $\vec D^{\,n}$ is not a diamond sequence
in $M_n$.  Let $(A,C)$ be the $<_{M_n}$-least counterexample, so that
$A\subseteq\omega_1$, $C\subseteq\omega_1$ is club, and
\[
   \forall\alpha\in C\,
      \bigl(D^n_\alpha\neq A\cap\alpha\bigr).
\]
Choose an $\omega$-sound initial segment
$\mathcal J^{M_n}_\theta$ containing
$A$, $C$, $\vec D^{\,n}$, and the relevant part of the canonical construction
order.  Take a countable elementary substructure
\[
   X\prec\mathcal J^{M_n}_\theta
\]
containing these parameters and such that
\[
   \alpha:=X\cap\omega_1\in C.
\]
Let
\[
   \pi:\bar X\longrightarrow X
\]
be the inverse of the transitive collapse.  The hull may be chosen so that
\[
   \operatorname{crit}(\pi)
   =\rho_\omega(\bar X)
   =\omega_1^{\bar X}=\alpha.
\]
By Theorem~\ref{theorem-condensation-lower-part-hulls-Mn}, the collapse
$\bar X$ is an initial segment of $M_n$.

The canonical construction order and the recursion defining
$\vec D^{\,n}$ are therefore computed correctly in $\bar X$.  The collapse
fixes all ordinals below $\alpha$ and sends $A$ and $C$ to
$A\cap\alpha$ and $C\cap\alpha$, respectively.  By elementarity,
$\bar X$ regards $(A\cap\alpha,C\cap\alpha)$ as the least counterexample to
the sequence below $\alpha$.  Consequently the recursion at stage $\alpha$
gives
\[
   D^n_\alpha=A\cap\alpha.
\]
This contradicts $\alpha\in C$.  Hence $\vec D^{\,n}$ is a
$\diamondsuit_{\omega_1}$-sequence.
\end{proof}

We fix from now on the sequence supplied by
Lemma~\ref{lemma-canonical-Mn-diamond-sequence}.  The phrase
\emph{the canonical $M_n$-diamond sequence} always refers to this fixed
sequence.

\begin{lemma}
\label{lemma-uniform-recovery-Mn-diamond}
There is a first-order formula
$\operatorname{CanDiamond}_n(X,M)$ such that, whenever
$N=M_n[G]$ is a small $M_n$-generic extension and
\[
   M=M_n|\omega_1
\]
is the set recovered in $N$ by
Definition~\ref{definition-recover-Mn-omega1}, there is a unique $X$ for
which
\[
   N\models\operatorname{CanDiamond}_n(X,M),
\]
and this $X$ is the true sequence $\vec D^{\,n}$.
\end{lemma}

\begin{proof}
The formula says that $X$ is obtained by the recursion in the proof of
Lemma~\ref{lemma-canonical-Mn-diamond-sequence}, using the construction order
on the reals of $M$.  By
Lemma~\ref{lemma-correctness-recovery-Mn-omega1}, every small generic
extension recovers the same transitive structure
$M_n|\omega_1$.  Fact~\ref{fact-canonical-construction-order-Mn}
then gives the same construction order, and transfinite recursion gives a
unique sequence.  Since the defining recursion is the one carried out in
$M_n$, that sequence is $\vec D^{\,n}$.
\end{proof}

The diamond sequence is used to construct the sequence of trees into which
reals will later be written.

\begin{lemma}
\label{lemma-canonical-independent-Mn-Suslin-sequence}
There is a sequence
\[
   \vec S^{\,n}
   =\langle S^n_\xi\mid\xi<\omega_1\rangle\in M_n
\]
of normal Suslin trees such that, for every nonempty finite set
$e\subseteq\omega_1$, the finite product
\[
   \prod_{\xi\in e}S^n_\xi
\]
is a Suslin tree.  The sequence is obtained by Jensen's simultaneous tree
construction from $\vec D^{\,n}$, with every choice made $<_{M_n}$-least.
In particular, it is uniformly definable from
$M_n|\omega_1$ and $\vec D^{\,n}$.
\end{lemma}

\begin{proof}
This is Jensen's simultaneous Suslin-tree construction from the proof of
Lemma~\ref{canonical-L-Suslin-sequence}, carried out inside $M_n$ and using
the canonical $M_n$-diamond sequence $\vec D^{\,n}$ from
Lemma~\ref{lemma-canonical-Mn-diamond-sequence}.  At successor stages we
preserve normality and splitting, and at limit stages the diamond guesses are
used to seal maximal antichains in finite products.  All choices are made
$<_{M_n}$-least.

The same argument as in the $L$-case shows that every finite product
\[
   \prod_{\xi\in e}S^n_\xi
\]
is Suslin.  Uniformity from $M_n|\omega_1$ follows because $\vec D^{\,n}$ and
the restriction of the construction order are uniformly recoverable from
$M_n|\omega_1$.
\end{proof}

We identify every $S^n_\xi$ with a subset of $\omega_1$ using the canonical
$M_n$-coding of countable normal trees.  As in the constructible case, we
shall say that a tree $S^n_\xi$ is \emph{destroyed} when a cofinal generic
branch through it is added.

\begin{corollary}
\label{corollary-uniform-recovery-Mn-Suslin-sequence}
There is a first-order formula
$\operatorname{CanSuslin}_n(X,M)$ such that, in every small
$M_n$-generic extension, the parameter
$M=M_n|\omega_1$ uniquely determines the true sequence
$X=\vec S^{\,n}$.
\end{corollary}

\begin{proof}
First recover $\vec D^{\,n}$ by
Lemma~\ref{lemma-uniform-recovery-Mn-diamond}, and then perform the canonical
simultaneous tree recursion from
Lemma~\ref{lemma-canonical-independent-Mn-Suslin-sequence}.  Both recursions
are unique and use only the recovered lower part and its construction order.
\end{proof}

We finish by fixing the two further canonical parameters needed for almost
disjoint coding and for the coding areas.  Canonical fine-structural mice
satisfy $\GCH$, and hence
\[
   M_n\models
   \bigl|({}^{<\omega_1}2)^{M_n}\bigr|=\omega_1.
\]

\begin{definition}
\label{definition-canonical-coding-parameters-Mn}
The \emph{canonical $M_n$-almost disjoint family} is
\[
   \mathcal D^n
   =\langle d^n_\xi\mid\xi<\omega_1\rangle,
\]
where, recursively, $d^n_\xi$ is the $<_{M_n}$-least infinite subset of
$\omega$ which is almost disjoint from every $d^n_\zeta$ with
$\zeta<\xi$.

The \emph{canonical $M_n$-area map} is the $<_{M_n}$-least bijection
\[
   \varrho_n:({}^{<\omega_1}2)^{M_n}
      \longrightarrow\omega_1.
\]
Equivalently, after identifying countable binary sequences with countable
subsets of $\omega_1$ by the fixed coding, we may regard $\varrho_n$ as a
bijection
\[
   \varrho_n:([\omega_1]^\omega)^{M_n}
      \longrightarrow\omega_1.
\]
\end{definition}

At stage $\xi<\omega_1$ of the first recursion there are only countably many
previous members of $\mathcal D^n$, so the required infinite set exists.  The
second object exists by the displayed cardinal computation.  Both are
uniformly definable from $M_n|\omega_1$ and the canonical
construction order.  We use the letter $\varrho_n$, rather than $\rho$, in
order not to confuse the area map with the ordinal indexing the allowability
hierarchy.

For a local lower-part model $m$ which carries out the same canonical
recursions, write
\[
   \vec D^{\,m},\qquad
   \vec S^{\,m},\qquad
   \mathcal D^m,
   \qquad\text{and}\qquad
   \varrho_m
\]
for its versions of the four objects just defined.  The following coherence
statement is the form needed in the later model-code localization arguments.

\begin{lemma}
\label{lemma-condensation-coherence-canonical-Mn-parameters}
Let $\mathcal J^{M_n}_\theta$ be an $\omega$-sound initial segment containing
\[
   \vec D^{\,n},\quad
   \vec S^{\,n},\quad
   \mathcal D^n,\quad
   \varrho_n,
\]
and let
\[
   X\prec\mathcal J^{M_n}_\theta
\]
be countable and contain these parameters.  Put
\[
   \lambda=X\cap\omega_1
\]
and let $\pi:\bar X\to X$ be the inverse of the transitive collapse.  Assume
that the hull is chosen as in
Theorem~\ref{theorem-condensation-lower-part-hulls-Mn}.  Then
$\bar X\trianglelefteq M_n$, $\omega_1^{\bar X}=\lambda$, and the collapse of
the four global parameters is exactly the corresponding system computed in
$\bar X$.  In particular,
\[
   \vec D^{\,\bar X}=\vec D^{\,n}\restriction\lambda,
   \qquad
   \vec S^{\,\bar X}=\vec S^{\,n}\restriction\lambda,
\]
\[
   \mathcal D^{\bar X}
      =\langle d^n_\xi\mid\xi<\lambda\rangle,
\]
and, for every
$s\in({}^{<\lambda}2)^{\bar X}$,
\[
   \varrho_{\bar X}(s)=\varrho_n(s)
\]
whenever the objects are identified through the collapse.
\end{lemma}

\begin{proof}
By Theorem~\ref{theorem-condensation-lower-part-hulls-Mn}, the transitive
collapse $\bar X$ is an initial segment of $M_n$.  The definitions of the
diamond sequence, the simultaneous Suslin-tree sequence, the almost disjoint
family, and the area map are all first-order definitions from the canonical
construction order.  Elementarity therefore sends each global object to the
object obtained by carrying out the same definition in the collapse.  Since
$\operatorname{crit}(\pi)=\lambda$, the collapse fixes every ordinal below
$\lambda$ and every subset of such an ordinal which belongs to $\bar X$.
The displayed identities follow.
\end{proof}

\begin{corollary}
\label{corollary-uniform-availability-canonical-Mn-parameters}
Let $N=M_n[G]$ be a small $M_n$-generic extension.  From the unique set
$M_n|\omega_1$ satisfying $\Theta_n$, the model $N$ uniformly
recovers
\[
   \vec D^{\,n},\qquad
   \vec S^{\,n},\qquad
   \mathcal D^n,
   \qquad\text{and}\qquad
   \varrho_n.
\]
Moreover, whenever a countable elementary hull is taken in the later
localization argument, its transitive collapse computes the corresponding local
versions of all four objects, as described in
Lemma~\ref{lemma-condensation-coherence-canonical-Mn-parameters}.
\end{corollary}

\begin{proof}
The recovery of $M_n|\omega_1$ is
Lemma~\ref{lemma-correctness-recovery-Mn-omega1}.  The uniform definitions of
the four parameters were established above.  The final assertion is exactly
the condensation-coherence lemma.
\end{proof}



\subsection[The preparatory branch extension Wn]{The preparatory branch extension $W_n$}
\label{subsection-preparatory-branch-extension-Mn}

We now pass to the ground model over which the coding iteration will be
performed.  Throughout this subsection, $\omega_1$ denotes
$\omega_1^{M_n}$.  Let
\[
   \vec S^{\,n}
   =\langle S^n_\xi\mid\xi<\omega_1\rangle
\]
be the canonical independent sequence of Suslin trees from
Lemma~\ref{lemma-canonical-independent-Mn-Suslin-sequence}.  All products in
this subsection are computed in $M_n$.

\begin{definition}
\label{definition-branch-products-over-Mn}
For $A\subseteq\omega_1$ with $A\in M_n$, put
\[
   \operatorname{Br}^n_A
   :=\prod_{\xi\in A}^{\mathrm{fs}}S^n_\xi.
\]
We write
\[
   \operatorname{Br}_n
   :=\operatorname{Br}^n_{\omega_1}
   =\prod_{\xi<\omega_1}^{\mathrm{fs}}S^n_\xi.
\]
Let $H\subseteq\operatorname{Br}_n$ be $M_n$-generic.  For
$\xi<\omega_1$, define
\[
   b_\xi
   :=\bigcup\{p(\xi):p\in H\text{ and }
                     \xi\in\operatorname{dom}(p)\}.
\]
For $A\subseteq\omega_1$ in $M_n$, let
\[
   H_A:=H\cap\operatorname{Br}^n_A.
\]
Finally, put
\[
   W_n:=M_n[H].
\]
\end{definition}

Thus $b_\xi$ is the branch added through $S^n_\xi$ by the $\xi$-th
coordinate.  As before, to \emph{destroy} a tree $S^n_\xi$ means to add a
cofinal generic branch through it.

\begin{lemma}
\label{lemma-branch-product-calculus-Mn}
Let $A,B\in M_n$ satisfy
\[
   A\subseteq B\subseteq\omega_1.
\]
Then the following hold in $M_n$.
\begin{enumerate}
   \item $\operatorname{Br}^n_A$ is c.c.c.
   \item $\operatorname{Br}^n_A$ is a complete subforcing of
         $\operatorname{Br}^n_B$, and
         \[
            \operatorname{Br}^n_B
            \cong
            \operatorname{Br}^n_A
            \times
            \operatorname{Br}^n_{B\setminus A}.
         \]
   \item If $\xi\notin A$, then
         \[
            \operatorname{Br}^n_A\Vdash
            ``S^n_\xi\text{ is a Suslin tree}."
         \]
\end{enumerate}
\end{lemma}

\begin{proof}
The sequence $\vec S^{\,n}$ is independent in $M_n$, so every nonempty
finite product of its members is a Suslin tree.  Apply
Lemma~\ref{finite-support-branch-products} inside $M_n$.  The second clause is
the coordinatewise factorization of the finite-support product.  For the last
clause, note that
\[
   \operatorname{Br}^n_A\times S^n_\xi
   \cong
   \operatorname{Br}^n_{A\cup\{\xi\}},
\]
which is c.c.c. by the first clause.  The product characterization of
preservation of a Suslin tree therefore gives the conclusion.
\end{proof}

\begin{theorem}
\label{theorem-properties-of-Wn}
The model $W_n=M_n[H]$ has the following properties.
\begin{enumerate}
   \item The forcing $\operatorname{Br}_n$ is c.c.c. and has size
         $\aleph_1$ in $M_n$.  Consequently, $W_n$ preserves every cardinal
         and cofinality of $M_n$, and
         \[
            W_n\models\CH.
         \]

   \item Every $S^n_\xi$ is destroyed in $W_n$, as witnessed by the branch
         $b_\xi$.

   \item If $A\subseteq\omega_1$ belongs to $M_n$, then $H_A$ is
         $\operatorname{Br}^n_A$-generic over $M_n$, and, for every
         $\xi<\omega_1$,
         \[
            M_n[H_A]\models
            ``S^n_\xi\text{ has a cofinal branch}"
            \quad\Longleftrightarrow\quad
            \xi\in A.
         \]

   \item Every real in $W_n$ belongs to $M_n[H_A]$ for some countable set
         $A\subseteq\omega_1$ with $A\in M_n$.

   \item More generally, every member of $H(\omega_1)^{W_n}$ belongs to
         $M_n[H_A]$ for some countable set
         $A\subseteq\omega_1$ with $A\in M_n$.
\end{enumerate}
\end{theorem}

\begin{proof}
This is the argument of Theorem~\ref{properties-of-W}, carried out inside
$M_n$, with $\vec S$, $\operatorname{Br}$, and $L$ replaced by
$\vec S^{\,n}$, $\operatorname{Br}_n$, and $M_n$, respectively.  
\end{proof}

The last two clauses are the branch-coordinate localization needed by the
coding argument.  We record a form in which a prescribed countable set of
coordinates may simultaneously be retained.

\begin{corollary}
\label{corollary-countable-branch-supports-Mn}
Let $a\in H(\omega_1)^{W_n}$, and let
$C\in[\omega_1]^\omega\cap M_n$.  Then there is
$A\in[\omega_1]^\omega\cap M_n$ such that
\[
   C\subseteq A,
   \qquad
   a\in M_n[H_A],
\]
and
\[
   M_n[H_A]\models
   ``S^n_\xi\text{ is a Suslin tree}"
   \qquad(\xi\in\omega_1\setminus A).
\]
In particular, no cofinal branch through any $S^n_\xi$ with
$\xi\notin A$ belongs to $M_n[H_A]$.
\end{corollary}

\begin{proof}
By Theorem~\ref{theorem-properties-of-Wn}(5), choose a countable branch support
$A_0\in M_n$ for $a$, and put
\[
   A=A_0\cup C.
\]
Then $A$ is countable and belongs to $M_n$.  Since
$\operatorname{Br}^n_{A_0}$ is a complete subforcing of
$\operatorname{Br}^n_A$, we have $a\in M_n[H_A]$.  The assertion about the
unused trees follows from
Lemma~\ref{lemma-branch-product-calculus-Mn}(3).
\end{proof}

\begin{corollary}
\label{corollary-canonical-parameters-branch-subextensions-Mn}
Let $A\subseteq\omega_1$ belong to $M_n$.  Then $M_n[H_A]$, and in
particular $W_n$, is a small $M_n$-generic extension in the sense of
Definition~\ref{definition-small-Mn-generic-extension}.  Consequently,
$M_n[H_A]$ uniformly recovers the true objects
\[
   M_n|\omega_1,
   \qquad
   \vec D^{\,n},
   \qquad
   \vec S^{\,n},
   \qquad
   \mathcal D^n,
   \qquad
   \varrho_n.
\]
The recovered objects are independent of $A$ and agree with the objects fixed
in $M_n$.
\end{corollary}

\begin{proof}
By Lemma~\ref{lemma-branch-product-calculus-Mn}, the forcing
$\operatorname{Br}^n_A$ is c.c.c., cardinal preserving, and has size at most
$\aleph_1$ in $M_n$.  It is therefore a small forcing in the sense of
Definition~\ref{definition-small-Mn-generic-extension}.  Apply
Lemma~\ref{lemma-correctness-recovery-Mn-omega1} and
Corollary~\ref{corollary-uniform-availability-canonical-Mn-parameters}.
\end{proof}

From now on, $W_n$ is the ground model of the coding and uniformization
iterations.  The full branch family
\[
   \langle b_\xi\mid\xi<\omega_1\rangle
\]
is available in $W_n$, while
Theorem~\ref{theorem-properties-of-Wn} and
Corollary~\ref{corollary-countable-branch-supports-Mn} allow any countable
collection of parameters to be localized to an inner branch extension which
omits all but countably many of these branches.  This tension between global
availability and local omission is the mechanism used below to rule out
unintended codes.

\subsection[The adapted Mn-coding mechanism]{The adapted $M_n$-coding mechanism}
\label{subsection-adapted-Mn-coding-mechanism}

We now adapt the coding forcing from Section~2 to the branch model $W_n$.
The canonical objects used in the definition are
\[
   M_n|\omega_1,
   \qquad
   \vec S^{\,n},
   \qquad
   \mathcal D^n,
   \qquad
   \varrho_n,
\]
and, by
Corollary~\ref{corollary-canonical-parameters-branch-subextensions-Mn},
they are computed uniformly and correctly in every branch subextension which
occurs in the localization arguments.  Throughout this subsection,
$\omega_1$ denotes the common first uncountable cardinal of $M_n$ and $W_n$.

\subsubsection*{The reservoir of coding areas}

\begin{definition}
\label{definition-Mn-Cohen-reservoir}
Put
\[
   \mathbb C_{M_n}(\omega_1)
   :=({}^{<\omega_1}2)^{M_n},
\]
ordered by end-extension.  Thus a condition is a countable binary sequence
belonging to $M_n$.  We use a tagged copy
$\mathbb C_{M_n}(\omega_1)_\eta$ for each $\eta<\omega_1$ and put
\[
   \forceC^{\mathrm{area}}_n
   :=\prod_{\eta<\omega_1}^{\mathrm{cs}}
      \mathbb C_{M_n}(\omega_1)_\eta.
\]
For $E\subseteq\omega_1$, let
\[
   (\forceC^{\mathrm{area}}_n)_E
   :=\prod_{\eta\in E}^{\mathrm{cs}}
      \mathbb C_{M_n}(\omega_1)_\eta.
\]
The canonical name for the generic subset of $\omega_1$ at coordinate
$\eta$ is denoted by $\dot g_\eta$.
\end{definition}

The product presentation is retained because its coordinates supply the
coding areas used by the later finite-support iteration.  As a forcing, the
reservoir is equivalent to a single copy of the $M_n$-computed
$\omega_1$-Cohen forcing.

\begin{lemma}
\label{lemma-basic-properties-Mn-Cohen-reservoir}
The following hold.
\begin{enumerate}
   \item In $M_n$, the forcing $\forceC^{\mathrm{area}}_n$ is
         $\sigma$-closed and has size $\aleph_1$.

   \item There is an
         $M_n$-definable isomorphism
         \[
            \forceC^{\mathrm{area}}_n
            \cong \mathbb C_{M_n}(\omega_1).
         \]
         The same holds for every product on a set of $\aleph_1$ many
         coordinates and after every permutation of the coordinates which
         belongs to $M_n$.

   \item In $W_n$ the forcing $\forceC^{\mathrm{area}}_n$ is
         $\omega$-distributive.  Consequently it adds no reals over $W_n$,
         preserves $\omega_1$, and preserves $\CH$.

   \item If $A,E\subseteq\omega_1$ belong to $M_n$, then the forcing
         $(\forceC^{\mathrm{area}}_n)_E$ may be commuted in front of
         $\operatorname{Br}^n_A$.  In the resulting reservoir extension,
         every tree $S^n_\xi$ which was Suslin before the branch forcing
         remains Suslin.
\end{enumerate}
\end{lemma}


\begin{proof}
This is the reservoir argument from
Subsection~\ref{subsection-allowable-forcings}, carried out with $M_n$ in
place of $L$. 
\end{proof}

\begin{definition}
\label{definition-Mn-Cohen-coding-area}
Let $g\subseteq\omega_1$ be
$\mathbb C_{M_n}(\omega_1)$-generic over a model containing the canonical
area map $\varrho_n$.  The \emph{coding area determined by $g$} is
\[
   h(g)
   :=\{\varrho_n(g\restriction\xi):\xi<\omega_1\}.
\]
If $K\subseteq\forceC^{\mathrm{area}}_n$ is generic, we write
\[
   g_\eta:=\dot g_\eta^K,
   \qquad
   h_\eta:=h(g_\eta)
\]
for the generic and the coding area at coordinate $\eta$.
\end{definition}

Every proper initial segment of $g$ belongs to $M_n$.  Moreover, since the
forcing is $\sigma$-closed in $M_n$, for every $\alpha<\omega_1$ the set
$h(g)\cap\alpha$ belongs to $M_n$.  Thus a coding area is a new unbounded
subset of $\omega_1$ all of whose bounded initial segments are already
available in the canonical lower part.

\begin{lemma}
\label{lemma-separation-Mn-coding-areas}
Let $g$ and $g'$ be distinct $M_n$-Cohen subsets of $\omega_1$.  Then
\[
   h(g)\cap h(g')
\]
is bounded in $\omega_1$.  In particular, the family
\[
   \langle h_\eta\mid\eta<\omega_1\rangle
\]
provided by a reservoir generic is pairwise almost disjoint modulo bounded
sets.
\end{lemma}


\begin{proof}
This is the same separation argument as in the $L$-construction, with the
area map $\varrho_L$ replaced by $\varrho_n$.  
\end{proof}

We retain the notation introduced in Section~2.  For $u\in2^\omega$,
$\gamma<\omega_1$, and $k<\omega$, put
\[
   \tau_u(\gamma,k)=
   \begin{cases}
      \omega\gamma+2k+1,&k\in u,\\
      \omega\gamma+2k,&k\notin u,
   \end{cases}
\]
and, for $h\subseteq\omega_1$, put
\[
   \operatorname{Sel}(h,u)
   :=\{\tau_u(\gamma,k):\gamma\in h,\ k<\omega\}.
\]
Thus the $k$-th bit of $u$ is written at the $\gamma$-th block by selecting
one of the two trees
\[
   S^n_{\omega\gamma+2k},
   \qquad
   S^n_{\omega\gamma+2k+1}.
\]

The following elementary strengthening of
Lemma~\ref{lemma-separation-Mn-coding-areas} is the combinatorial form used in
the later no-unwanted-codes argument.

\begin{lemma}
\label{lemma-countably-many-Mn-coding-areas}
Let $I$ be countable, let
$\langle h_i\mid i\in I\rangle$ be coding areas obtained from distinct
reservoir coordinates, let $u_i\in2^\omega$, and let
$B\in[\omega_1]^\omega$.  Put
\[
   U:=B\cup\bigcup_{i\in I}\operatorname{Sel}(h_i,u_i).
\]
Then:
\begin{enumerate}
   \item If $h\subseteq\omega_1$ has size $\omega_1$ and
         $h\cap h_i$ is bounded for every $i\in I$, then there are
         $\omega_1$ many $\gamma\in h$ such that
         \[
            [\omega\gamma,\omega\gamma+\omega)\cap U=\varnothing.
         \]

   \item Suppose that $i_0\in I$, that $h\cap h_{i_0}$ is unbounded, and
         that $u\neq u_{i_0}$.  If $k<\omega$ satisfies
         $u(k)\neq u_{i_0}(k)$, then there are $\omega_1$ many
         $\gamma\in h\cap h_{i_0}$ such that
         \[
            \tau_u(\gamma,k)\notin U.
         \]
\end{enumerate}
\end{lemma}


\begin{proof}
This is the countable version of the preceding separation argument, exactly as
in the $L$-case.  By Lemma~\ref{lemma-separation-Mn-coding-areas}, distinct
coding areas have bounded intersection.  Since $I$ is countable, one bound
controls all pairwise intersections among the areas under consideration.  Also
exclude the countable set of blocks meeting $B$.

In case (1), any sufficiently large $\gamma\in h$ outside these exceptional
sets has its whole $\omega$-block disjoint from $U$.  In case (2), choose
such a sufficiently large $\gamma\in h\cap h_{i_0}$; the inequality
$u(k)\neq u_{i_0}(k)$ means that, at the $k$-th pair of the $\gamma$-block,
$u$ selects the tree opposite to the one selected by $u_{i_0}$.  Thus
$\tau_u(\gamma,k)\notin U$.  In both cases there are $\omega_1$ many possible
choices of $\gamma$.
\end{proof}

\subsubsection*{Writing a real into the branch sequence}

Let $u\in2^\omega$ belong to a forcing extension of $W_n$, let $g$ be a fresh
$M_n$-Cohen generic, and put $h=h(g)$.  The selected branch sequence is
\[
   \mathcal B^n_{g,u}
   :=\left\langle
      b_{\tau_u(\gamma,k)}
      \,\middle|\,
      \gamma\in h,\ k<\omega
   \right\rangle.
\]
We regard this as an indexed set of triples
\[
   (\gamma,k,b_{\tau_u(\gamma,k)}).
\]
Consequently the index and the tree through which the displayed branch runs
recover the bit $u(k)$.  This formulation uses the distinguished selected
branch sequence; it does not make any assertion about branches through the
opposite member of the pair in the ambient universe.

Put
\[
   m^n_\infty:=M_n|\omega_1.
\]
Choose canonically a transitive model
\[
   \mathcal N_{g,u}
\]
of size $\aleph_1$ such that:
\begin{enumerate}
   \item
   $\mathcal N_{g,u}\models
      \ZFCminus+``\aleph_1\text{ exists}"$
   and $\omega_1^{\mathcal N_{g,u}}=\omega_1$;

   \item
   $m^n_\infty$, $h$, and $\mathcal B^n_{g,u}$ belong to
   $\mathcal N_{g,u}$;

   \item using $m^n_\infty$, the model $\mathcal N_{g,u}$ computes
   $\vec S^{\,n}$, $\mathcal D^n$, and $\varrho_n$ correctly;

   \item $\mathcal N_{g,u}$ verifies that every displayed member of
   $\mathcal B^n_{g,u}$ is a cofinal branch through the corresponding tree
   and reads the same bit of $u$ at every block whose starting point belongs
   to $h$.
\end{enumerate}
For definiteness, take the least model of the form
\[
   L_\theta[m^n_\infty,h,\mathcal B^n_{g,u}]
\]
which satisfies the displayed theory and contains the indicated objects.
The usual reflection argument gives such a $\theta$ of cardinality
$\aleph_1$.  This choice supplies a canonical constructibility order and
therefore a canonical enumeration of $\mathcal N_{g,u}$ of length
$\omega_1$.

Let
\[
   X_{g,u}\subseteq\omega_1
\]
be the resulting canonical code for
\[
   (\mathcal N_{g,u},\in,
       m^n_\infty,h,\mathcal B^n_{g,u}).
\]
The code records a well-founded extensional relation on $\omega_1$ together
with constants for the three displayed objects.  Decoding $X_{g,u}$ and taking
the transitive collapse therefore recovers the entire tuple.

\subsubsection*{Localization of the model code}

We now apply to $X_{g,u}$ the localization procedure introduced in the
$L$-case.

For $Y\subseteq\omega_1$, put
\[
   O(Y):=\{\alpha<\omega_1:2\alpha+1\in Y\},
   \qquad
   E(Y):=\{\alpha<\omega_1:2\alpha\in Y\}.
\]
For every countable ordinal $\alpha$, let $\operatorname{wo}_n(\alpha)$ be
the $<_{M_n}$-least real coding a wellorder of $\omega$ of order type
$\alpha$.

Choose canonically an ordinal $\lambda$ of size $\aleph_1$ such that
\[
   L_\lambda[X_{g,u}]
   \models\ZFCminus+``\aleph_1\text{ exists}",
\]
and, inside this model, choose the canonical continuous increasing chain
\[
   \langle M_\xi\mid\xi<\omega_1\rangle
\]
of countable elementary submodels whose union is
$L_\lambda[X_{g,u}]$.  We require that all decoding parameters belong to
$M_0$ and that, putting
\[
   c_\xi:=M_\xi\cap\omega_1,
\]
one has
\[
   c_0>\omega,
   \qquad
   c_{\xi+1}>c_\xi+\omega.
\]
Let
\[
   C_{g,u}:=\{c_\xi:\xi<\omega_1\}.
\]

\begin{definition}\label{definition-Mn-localized-code}
The \emph{localization of $X_{g,u}$} is the unique set
\[
   Y_{g,u}\subseteq\omega_1
\]
with
\[
   O(Y_{g,u})=X_{g,u}
\]
and whose even part is defined as follows:
\begin{enumerate}
   \item $E(Y_{g,u})\cap\omega$ is the fixed real
         $\operatorname{wo}_n(c_0)$;
   \item $E(Y_{g,u})\cap[\omega,c_0)=\varnothing$;
   \item for every $\xi<\omega_1$,
         $E(Y_{g,u})\cap[c_\xi,c_\xi+\omega)$ is the translate of
         $\operatorname{wo}_n(c_{\xi+1})$ to that interval;
   \item for every $\xi<\omega_1$,
         \[
            E(Y_{g,u})\cap[c_\xi+\omega,c_{\xi+1})
            =\varnothing.
         \]
\end{enumerate}
\end{definition}
We also write
\[
   Y_n(u,g):=Y_{g,u}.
\]


All choices in the construction of $\mathcal N_{g,u}$, $X_{g,u}$, the hull
sequence, and $Y_{g,u}$ are made by the canonical orders fixed above.  Hence,
for forcing names $\dot u$ and $\dot g$, the notation
\[
   \dot Y_n(\dot u,\dot g)
\]
denotes a canonical forcing name.

\begin{lemma}
\label{lemma-locality-Mn-localization}
Let $M$ be a countable transitive model of
$\ZFCminus+``\aleph_1\text{ exists}"$ and put
$\alpha=\omega_1^M$.  Suppose that
\[
   Y_{g,u}\cap\alpha\in M.
\]
Then $\alpha\in C_{g,u}$.  Moreover, inside its version of
$L[Y_{g,u}\cap\alpha]$, the model $M$ recovers
\[
   X_{g,u}\cap\alpha
\]
and the well-founded extensional structure coded by this initial segment.
Its transitive collapse is the local structure obtained by collapsing the
corresponding elementary hull $M_\xi$ with $c_\xi=\alpha$.  In particular,
the local collapse contains the appropriate initial segment of
$m^n_\infty$, the local coding area, and the local selected branch sequence.
\end{lemma}

\begin{proof}
This is the same localization argument as in the $L$-case.  The odd
coordinates of $Y_{g,u}\cap\alpha$ give $X_{g,u}\cap\alpha$.  The even
coordinates force $\alpha\in C_{g,u}$: if $\alpha<c_0$, or if
$c_\xi<\alpha<c_{\xi+1}$, then, in the latter case using
$c_\xi+\omega<\alpha$, the model $M$ contains the real
$\operatorname{wo}_n(c_0)$, respectively $\operatorname{wo}_n(c_{\xi+1})$,
and hence a well-order of $\omega$ of type at least
$\alpha=\omega_1^M$, contradiction.  Thus $\alpha=c_\xi$ for some $\xi$.
By the choice of the elementary hull chain, $X_{g,u}\cap c_\xi$ is exactly
the code of the collapse of $M_\xi$.  The distinguished constants collapse to
the corresponding local initial segment of $m^n_\infty$, the local coding
area, and the local selected branch sequence.  The remaining assertions follow
from uniqueness of the Mostowski collapse and from
Lemma~\ref{lemma-condensation-coherence-canonical-Mn-parameters}.
\end{proof}

\subsubsection*{Almost disjoint coding and the forcing
$\operatorname{Code}_n$}

For $Y\subseteq\omega_1$, let
\[
   \mathbb A^n(Y)
   :=\mathbb A_{\mathcal D^n}(Y)
\]
be the Jensen--Solovay almost disjoint coding forcing relative to the
canonical family
\[
   \mathcal D^n=\langle d^n_\xi\mid\xi<\omega_1\rangle.
\]
Thus, if $G\subseteq\mathbb A^n(Y)$ is generic and
\[
   r_G:=\bigcup\{s:\exists R\ ((s,R)\in G)\},
\]
then
\[
   \xi\in Y
   \quad\Longleftrightarrow\quad
   r_G\cap d^n_\xi\text{ is finite}
\]
for every $\xi<\omega_1$.

\begin{definition}
\label{definition-Mn-coding-forcing}
Let $u\in2^\omega$ and let $\eta<\omega_1$.  The tagged two-step coding
forcing is
\[
   \operatorname{Code}_n(u,\eta)
   :=
   \mathbb C_{M_n}(\omega_1)_\eta
   *\dot{\mathbb A}^{\,n}
      \bigl(\dot Y_n(\check u,\dot g_\eta)\bigr).
\]
If $u$ is given by a name $\dot u$, we use the corresponding name
\[
   \operatorname{Code}_n(\dot u,\eta)
   :=
   \mathbb C_{M_n}(\omega_1)_\eta
   *\dot{\mathbb A}^{\,n}
      \bigl(\dot Y_n(\dot u,\dot g_\eta)\bigr).
\]
\end{definition}

The ordinal $\eta$ merely tags the copy of the Cohen forcing in the standalone
definition.  Its real role appears in an iteration over the global reservoir.
Once $K\subseteq\forceC^{\mathrm{area}}_n$ has been added, the coordinate
$g_\eta$ is already present, and a coding stage using the fresh area $\eta$
consists only of the c.c.c. forcing
\[
   \mathbb A^n\bigl(Y_n(u,g_\eta)\bigr).
\]
Thus, exactly as in the $L$-construction, the Cohen reservoir is added once,
while only the almost disjoint coding forcings are iterated with finite
support.

\begin{lemma}
\label{lemma-basic-properties-Mn-coding-forcing}
Let $u\in2^\omega$ belong to the model over which the coding is performed.
Then:
\begin{enumerate}
   \item the reservoir step adds no real over $W_n$;

   \item in the reservoir extension,
         \[
            \mathbb A^n(Y_n(u,g_\eta))
         \]
         is Knaster and has size $\aleph_1$;

   \item $\operatorname{Code}_n(u,\eta)$ preserves $\omega_1$ and $\CH$;

   \item the generic real $r$ added by the second step canonically decodes
         $Y_n(u,g_\eta)$ and hence the localized code of
         \[
            (\mathcal N_{g_\eta,u},\in,
             m^n_\infty,h_\eta,\mathcal B^n_{g_\eta,u});
         \]

   \item the branch coordinates occurring in the construction of the
         localized code are exactly
         \[
            \operatorname{Sel}(h_\eta,u).
         \]
\end{enumerate}
\end{lemma}

\begin{proof}
Clause (1) is
Lemma~\ref{lemma-basic-properties-Mn-Cohen-reservoir}(3).  Clause (2) follows
from Lemma~\ref{a.d.coding preserves Suslin trees}, applied to the canonical
almost disjoint family $\mathcal D^n$.  The forcing has size $\aleph_1$
because $\mathcal D^n$ has length $\omega_1$.

The first step preserves $\omega_1$ and $\CH$, and the second step is c.c.c.
of size $\aleph_1$; hence (3) follows.  Clause (4) is the defining property of
Jensen--Solovay almost disjoint coding together with the fixed localization
decoding procedure.  Finally, $Y_n(u,g_\eta)$ depends on the branch family only through
$\mathcal B^n_{g_\eta,u}$, whose members are precisely the branches indexed by
$\operatorname{Sel}(h_\eta,u)$.  This proves (5).
\end{proof}




\subsection{The projective coding predicates and exactness}
\label{subsection-projective-Mn-coding-predicates}

We now isolate the projective content of the forcing
$\operatorname{Code}_n$.  

\begin{lemma}
\label{lemma-coherence-lower-part-approximations-Mn}
Let $N$ be a small $M_n$-generic extension, let
$m\in\mathcal I_n^N$, and put
\[
   \alpha:=\omega_1^m.
\]
Then the canonical objects computed in $m$ agree with the corresponding
initial parts of the global objects.  In particular,
\[
   \vec S^{\,m}=\vec S^{\,n}\restriction\alpha,
   \qquad
   \mathcal D^m
      =\langle d^n_\xi\mid\xi<\alpha\rangle,
\]
and, under the canonical identification of the domains,
\[
   \varrho_m(s)=\varrho_n(s)
\]
for every $s\in({}^{<\alpha}2)^m$.
\end{lemma}

\begin{proof}
By Lemma~\ref{lemma-definable-Mn-initial-segments}, the premouse $m$ is a
proper initial segment of $M_n$.  The canonical parameters are obtained by
first-order recursions from the construction order of $M_n$, and the
restriction of that order to $m$ is the construction order computed by $m$.
The uniqueness of the recursions therefore gives the displayed identities.
Equivalently, one may take a countable hull containing the relevant global
parameters and apply
Lemma~\ref{lemma-condensation-coherence-canonical-Mn-parameters}.
\end{proof}

\subsubsection*{The local decoding predicate}

Let $M$ be a transitive model of
$\ZFCminus+``\aleph_1\text{ exists}"$, let $m\in M$ be a premouse, and suppose
that
\[
   \omega_1^m=\omega_1^M=:\alpha.
\]
Inside $M$, the notation
\[
   \operatorname{LocalCode}_n(Y,u,m)
\]
denotes the following first-order assertion about
$Y\subseteq\alpha$ and $u\in2^\omega$.

First, $Y$ is a correctly localized code in the sense of
Definition~\ref{definition-Mn-localized-code}.  Its odd part decodes a set
$X\subseteq\alpha$, its even part supplies the corresponding localization data,
and $X$ codes a well-founded extensional structure with distinguished
constants.  The transitive collapse of that structure is a model
$\bar{\mathcal N}$ of cardinality $\aleph_1^M$ such that
\[
   \bar{\mathcal N}\models
      \ZFCminus+``\aleph_1\text{ exists}",
   \qquad
   \omega_1^{\bar{\mathcal N}}=\alpha.
\]
The distinguished constants decode:
\begin{enumerate}
   \item a lower-part premouse $m^\ast$ with
         $\omega_1^{m^\ast}=\alpha$;
   \item an unbounded coding area $h\subseteq\alpha$;
   \item an indexed branch sequence
         \[
            \mathcal B
            =\langle c_{\gamma,k}
               \mid\gamma\in h,\ k<\omega\rangle.
         \]
\end{enumerate}
The inverse images of the members of $h$ under $\varrho_{m^\ast}$ form an
end-extension chain.  Its union is a sequence
$g\in{}^\alpha2$ such that
\[
   \forall\xi<\alpha\,
   (g\restriction\xi\in m^\ast),
   \qquad
   h=\{\varrho_{m^\ast}(g\restriction\xi):
        \xi<\alpha\}.
\]
The premice $m$ and $m^\ast$ compute the same local canonical parameters
\[
   \vec S^{\,m},\qquad
   \mathcal D^m,\qquad
   \varrho_m.
\]
Finally, for every $\gamma\in h$ and $k<\omega$, the object
$c_{\gamma,k}$ is a cofinal branch through
\[
   S^m_{\tau_u(\gamma,k)}.
\]
Thus the real obtained from the parity of the tree indices in the
distinguished branch sequence is exactly $u$.  We make no assertion about
branches through the opposite members of the pairs in the ambient universe;
in $W_n$ every tree in $\vec S^{\,n}$ already has a branch.

All clauses in this description are first-order over $M$.  In particular,
$\operatorname{LocalCode}_n$ does not quantify over an iteration strategy.
The projective complexity enters only through the external requirement that
$m$ belong to $\mathcal I_n$.

Let
\[
   \mathcal D^m
   =\langle d^m_\xi\mid\xi<\omega_1^m\rangle
\]
be the almost disjoint family computed in $m$.

\begin{definition}
\label{definition-Psi-Phi-Mn}
For reals $r,u$, let $\Psi_n(r,u)$ be the assertion
\begin{align*}
\forall M\,\forall m\Big(&
   M\text{ is a countable transitive model of }
      \ZFCminus+``\aleph_1\text{ exists}"\\
&\wedge\ r,u,m\in M\\
&\wedge\ m\in\mathcal I_n\\
&\wedge\ \omega_1^m=\omega_1^M\\
&\Longrightarrow
   \exists Y\in M\Big[
      Y\subseteq\omega_1^M\\
&\qquad\wedge\
      \forall\xi<\omega_1^M\,
      \bigl(\xi\in Y
      \longleftrightarrow
      r\cap d^m_\xi\text{ is finite}\bigr)\\
&\qquad\wedge\
      M\models\operatorname{LocalCode}_n(Y,u,m)
   \Big]\Big).
\end{align*}
We then put
\[
   \Phi_n(u)
   \quad\Longleftrightarrow\quad
   \exists r\,\Psi_n(r,u).
\]
When $\Phi_n(u)$ holds, we say that $u$ is \emph{coded into}
$\vec S^{\,n}$, or that $u$ is \emph{written into} $\vec S^{\,n}$.
\end{definition}

\begin{lemma}
\label{lemma-complexity-Psi-Phi-Mn}
The predicate $\Psi_n(r,u)$ is
$\boldsymbol{\Pi}^1_{n+2}$, uniformly in $r,u$, and $\Phi_n(u)$ is
$\boldsymbol{\Sigma}^1_{n+3}$.
\end{lemma}

\begin{proof}
We use reals to code the countable structures $M$ and $m$.  The assertions
that a real codes a well-founded countable transitive model of the displayed
fragment of set theory, that the indicated objects belong to it, and that it
satisfies the fixed first-order formula
$\operatorname{LocalCode}_n$ have complexity at most
$\boldsymbol{\Pi}^1_1$, or are arithmetic once well-foundedness has been
assumed.  By Lemma~\ref{lemma-definable-Mn-initial-segments}, the assertion
$m\in\mathcal I_n$ is $\boldsymbol{\Pi}^1_{n+1}$.  The existential
quantifier over $Y\in M$ is a quantifier over an element of the coded
countable structure and hence is arithmetic in its code.

After rewriting the implication as a disjunction, the matrix following the
two universal real quantifiers is $\boldsymbol{\Sigma}^1_{n+1}$.  Therefore
$\Psi_n$ is $\boldsymbol{\Pi}^1_{n+2}$.  Adding the existential real
quantifier for $r$ gives
$\Phi_n\in\boldsymbol{\Sigma}^1_{n+3}$.
\end{proof}

\subsubsection*{The global objects decoded by a witness}

\begin{lemma}
\label{lemma-inner-model-decoded-by-Phi-Mn}
Work in a small $M_n$-generic extension and suppose that
$\Psi_n(r,u)$ holds.  Put
\[
   m^n_\infty=M_n|\omega_1.
\]
Then the inner model $L[r,m^n_\infty]$ contains $u$ and canonically decodes
objects
\[
   Y_r,\quad X_r,\quad \mathcal N_r,
   \quad g_r,\quad h_r,\quad \mathcal B_r
\]
with the following properties:
\begin{enumerate}
   \item
   \[
      Y_r
      =\{\xi<\omega_1:
         r\cap d^n_\xi\text{ is finite}\};
   \]
   the odd part of $Y_r$ decodes $X_r$, and $\mathcal N_r$ is the
   transitive collapse of the structure coded by $X_r$;

   \item $g_r\in{}^{\omega_1}2$, every proper initial segment of $g_r$
   belongs to $m^n_\infty$, and
   \[
      h_r
      =\{\varrho_n(g_r\restriction\xi):
           \xi<\omega_1\};
   \]
   in particular, $|h_r|=\omega_1$;

   \item
   \[
      \mathcal B_r
      =\langle c^r_{\gamma,k}
         \mid\gamma\in h_r,\ k<\omega\rangle
   \]
   is an indexed family of actual cofinal branches and
   \[
      c^r_{\gamma,k}
      \text{ is a branch through }
      S^n_{\tau_u(\gamma,k)}
   \]
   for every $\gamma\in h_r$ and $k<\omega$.
\end{enumerate}
Consequently,
\[
   \operatorname{Sel}(h_r,u)
   \subseteq
   \{\xi<\omega_1:
      L[r,m^n_\infty]
      \text{ contains a cofinal branch through }S^n_\xi\}.
\]
\end{lemma}

\begin{proof}
The almost disjoint decoding of $r$ relative to the canonical family
$\mathcal D^n$ is unique, so $L[r,m^n_\infty]$ can form the displayed set
$Y_r$.  We claim that $L[r,m^n_\infty]$ can canonically decode from $Y_r$ the
objects
\[
   X_r,\quad \mathcal N_r,\quad g_r,\quad h_r,\quad \mathcal B_r
\]
appearing in the statement of the lemma.

Suppose otherwise.  Then a part of the first-order decoding fails:
the localization test fails, the odd part fails to code the required well-founded
extensional structure, one of the distinguished constants is missing, or one
of the branch clauses fails.  Choose a sufficiently large regular $\Theta$
and a countable elementary submodel
\[
   Z\prec H_\Theta
\]
which contains $r,u,m^n_\infty$ and a witness to the failure.  We choose $Z$
from the standard club of lower-part hulls for which the collapse of the
lower-part parameter belongs to $\mathcal I_n$.  Let
\[
   \pi:Z\longrightarrow\bar M
\]
be the transitive collapse and put
\[
   m:=\pi(m^n_\infty).
\]
By Lemmas~\ref{lemma-definable-Mn-initial-segments} and
\ref{lemma-condensation-coherence-canonical-Mn-parameters},
\[
   m\in\mathcal I_n,
   \qquad
   \omega_1^m=\omega_1^{\bar M},
\]
and $m$ computes the correct initial parts of
$\vec S^{\,n}$, $\mathcal D^n$, and $\varrho_n$.

Inside $\bar M$, the almost disjoint decoding of $r$ relative to
$\mathcal D^m$ is exactly
\[
   Y_r\cap\omega_1^{\bar M}.
\]
Indeed, for every $\xi<\omega_1^{\bar M}$,
Lemma~\ref{lemma-coherence-lower-part-approximations-Mn} gives
$d^m_\xi=d^n_\xi$.  The failure of the global decoding reflects to
$\bar M$, and hence
\[
   \bar M\not\models
   \operatorname{LocalCode}_n
      (Y_r\cap\omega_1^{\bar M},u,m).
\]
On the other hand, $\Psi_n(r,u)$ applies to $(\bar M,m)$ and produces a set
$Y\in\bar M$ which is almost-disjointly decoded from $r$ and satisfies
\[
   \bar M\models\operatorname{LocalCode}_n(Y,u,m).
\]
The almost disjoint decoding is unique, so
\[
   Y=Y_r\cap\omega_1^{\bar M},
\]
a contradiction.  Thus the global decoding succeeds.

Let $u^*$ be the real read in $L[r,m^n_\infty]$ from the parity of the tree
indices in the decoded branch sequence.  If $u^*\neq u$, choose
$k<\omega$ on which they differ and repeat the preceding countable-hull
argument, now including enough of the decoded data to determine $u^*(k)$.
The unique local decoding would read the same indexed branch family as coding
both $u^*(k)$ and $u(k)$, which is impossible.  Hence $u^*=u$, so
$u\in L[r,m^n_\infty]$.

The inverse images under $\varrho_n$ of the members of the decoded area form
an end-extension chain of proper initial segments.  Their union is
$g_r$, and injectivity of $\varrho_n$ gives $|h_r|=\omega_1$.  The remaining
conclusions follow directly from the successful global decoding.
\end{proof}

\begin{theorem}
\label{theorem-correctness-persistence-Code-Mn}
Let $N$ be a small $M_n$-generic extension containing a real $u$, and let
$r$ be the real added by $\operatorname{Code}_n(u,\eta)$ over $N$.  Let $g_\eta$ denote the Cohen generic added by the first step.  Then
\[
   N[g_\eta][r]\models\Psi_n(r,u),
\]
and hence
\[
   N[g_\eta][r]\models\Phi_n(u).
\]
Moreover, the same real $r$ witnesses $\Psi_n(r,u)$ in every later small
$M_n$-generic extension of $N[g_\eta][r]$ occurring in the construction.
\end{theorem}

\begin{proof}
Let $g_\eta$ be the Cohen generic at the first step and put
\[
   Y=Y_n(u,g_\eta).
\]
The real $r$ added at the second step almost-disjointly codes $Y$ relative to
$\mathcal D^n$.

Work in $N[g_\eta][r]$ and let $M,m$ satisfy the antecedent in the
definition of $\Psi_n(r,u)$.  Put $\alpha=\omega_1^M$.  By
Lemma~\ref{lemma-coherence-lower-part-approximations-Mn}, $m$ computes the
initial almost disjoint family and the remaining canonical parameters
correctly.  Therefore the set decoded from $r$ inside $M$ is
\[
   Y\cap\alpha.
\]
Lemma~\ref{lemma-locality-Mn-localization}, together with condensation
coherence, shows that $M$ decodes from this set the corresponding local
collapse of the model coded by $Y$.  Its distinguished area and branch
sequence satisfy the $u$-pattern.  Hence
\[
   M\models
   \operatorname{LocalCode}_n(Y\cap\alpha,u,m).
\]
This proves $\Psi_n(r,u)$.

The formula $\Psi_n(r,u)$ is $\Pi^1_{n+2}$ by
Lemma~\ref{lemma-complexity-Psi-Phi-Mn}.  The persistence assertion therefore
follows from
Theorem~\ref{theorem-small-two-step-projective-absoluteness}, applied to the
relevant intermediate and later small forcing extensions.  The witness $r$
is unchanged.
\end{proof}

\subsubsection*{Coding iterations and exactness}

The exactness argument is most conveniently proved already for the basic
coding iterations which will form level zero of the allowability hierarchy.

\begin{definition}
\label{definition-Mn-coding-iteration}
Let $\delta\leq\omega_1$.  An \emph{$M_n$-coding iteration over the
reservoir of length $\delta$} consists of an injective sequence
\[
   \vec\eta=\langle\eta_\beta\mid\beta<\delta\rangle\in M_n
\]
and a sequence of  names
\[
   \langle\dot u_\beta\mid\beta<\delta\rangle
\]
for reals.  Put
\[
   U_\beta:=\operatorname{ran}
      (\vec\eta\restriction\beta).
\]
Recursively define
\[
   \forceP_\beta
   =\forceC^{\mathrm{area}}_n*\dot{\forceR}_\beta
   \qquad(\beta\leq\delta)
\]
as follows.  The forcing $\dot{\forceR}_0$ is trivial.  At stage $\beta$,
the name $\dot u_\beta$ is required to be the canonical lift of a name over
\[
   (\forceC^{\mathrm{area}}_n)_{U_\beta}
      *\dot{\forceR}_\beta,
\]
and
\[
   \dot{\forceR}_{\beta+1}
   =\dot{\forceR}_\beta*
      \dot{\mathbb A}^{\,n}
      \bigl(\dot Y_n(\dot u_\beta,
                      \dot g_{\eta_\beta})\bigr).
\]
At limit stages, $\dot{\forceR}_\lambda$ is the finite-support limit.  If
$K*G$ is generic for $\forceP_\delta$, put
\[
   u_\beta:=\dot u_\beta^{K*G},
   \qquad
   h_\beta:=h(g_{\eta_\beta}).
\]
The reals $u_\beta$, $\beta<\delta$, are called the reals
\emph{intentionally coded} by the iteration.
\end{definition}

\begin{lemma}
\label{lemma-size-preservation-Mn-coding-iterations}
Every $M_n$-coding iteration of length at most $\omega_1$ has size
$\aleph_1$, preserves cardinals, and preserves $\CH$.  Over the reservoir
extension, its almost disjoint tail is c.c.c.
\end{lemma}

\begin{proof}
The reservoir has size $\aleph_1$, is $\omega$-distributive over $W_n$, and
preserves $\CH$ by
Lemma~\ref{lemma-basic-properties-Mn-Cohen-reservoir}.  Every tail factor is
Knaster and has size $\aleph_1$ by
Lemma~\ref{lemma-basic-properties-Mn-coding-forcing}.  A finite-support
iteration of at most $\omega_1$ many c.c.c. forcings is c.c.c.; under $\CH$
the iteration has size $\aleph_1$.  The asserted preservation follows.
\end{proof}

For $\xi<\omega_1$, put
\[
   W_{n,\xi}
   :=M_n[H_{\omega_1\setminus\{\xi\}}].
\]
This is a genuine branch subextension defined by a ground-model product.  It
contains every branch $b_\zeta$ with $\zeta\neq\xi$, while $S^n_\xi$ remains
Suslin in $W_{n,\xi}$.

\begin{lemma}
\label{lemma-localization-branch-omission-Mn}
Let $\forceP_\delta$ be an $M_n$-coding iteration over the reservoir, let
$K*G$ be generic over $W_n$, and let $\vec x$ be a finite tuple of reals in
$W_n[K][G]$.  There are countable sets
\[
   E\subseteq\omega_1,
   \qquad
   I\subseteq\delta,
   \qquad
   B\subseteq\omega_1,
\]
with $B\in M_n$, such that:
\begin{enumerate}
   \item $I$ is recursively closed in the following sense: whenever
         $\beta\in I$, the name $\dot u_\beta$ defining the
         $\beta$-th iterand is a name over the induced closed subiteration
         on $I\cap\beta$ together with the reservoir coordinates in $E$,
         and $\eta_\beta\in E$.

   \item the tuple $\vec x$ is read by the induced closed subiteration on
         $I$ together with the reservoir coordinates in $E$;

   \item for every
         \[
            \xi\notin
            B\cup\bigcup_{\beta\in I}
               \operatorname{Sel}(h_\beta,u_\beta),
         \]
         the induced forcing on $(E,I)$ is computed in the same way over
         $W_{n,\xi}$ as over $W_n$, and
         \[
            \vec x\in
            W_{n,\xi}[K\restriction E][G\restriction I];
         \]

   \item for every such $\xi$,
         \[
            W_{n,\xi}[K\restriction E][G\restriction I]
            \models
            ``S^n_\xi\text{ is a Suslin tree}."
         \]
\end{enumerate}
Here $G\restriction I$ denotes the generic sequence for the induced closed
subiteration.
\end{lemma}
\begin{proof}
This is the proof of Lemma~\ref{allowable-localization-lemma}, with $L$,
$\vec S$, $D$, and $\varrho_L$ replaced by $M_n$, $\vec S^{\,n}$,
$\mathcal D^n$, and $\varrho_n$.

We recall the points where the $M_n$-setup enters.  Nice names have countable
reservoir and iteration support, since the reservoir adds no reals and the
almost disjoint tail is c.c.c.  Closing these supports recursively under the
names defining the retained iterands gives countable sets $E$ and $I$ with
the displayed closure property, and Theorem~\ref{theorem-properties-of-Wn}
gives a countable branch support $B\in M_n$ for the ground-model parameters.

If
\[
   \xi\notin B\cup\bigcup_{\beta\in I}\operatorname{Sel}(h_\beta,u_\beta),
\]
then the same induced subiteration on $(E,I)$ is computed over
$W_{n,\xi}$: no retained iterand uses the branch $b_\xi$.  The tree
$S^n_\xi$ is Suslin in $W_{n,\xi}$ by
Lemma~\ref{lemma-branch-product-calculus-Mn}, and each retained almost
disjoint factor is Knaster, hence preserves it.  Thus the restricted generic
interprets the chosen names for $\vec x$ in
$W_{n,\xi}[K\restriction E][G\restriction I]$, and $S^n_\xi$ remains Suslin
there.
\end{proof}

\begin{theorem}
\label{theorem-no-unwanted-Mn-codes}
Let $\forceP_\delta$ be an $M_n$-coding iteration over the reservoir, where
$\delta\leq\omega_1$, and let $K*G$ be generic over $W_n$.  Then
\[
   W_n[K][G]\models\Phi_n(u)
   \quad\Longleftrightarrow\quad
   u=u_\beta\text{ for some }\beta<\delta.
\]
Thus $\Phi_n$ recognizes exactly the reals intentionally coded by the
iteration.
\end{theorem}

\begin{proof}
This is the no-unwanted-codes argument from the $L$-case, with
Lemma~\ref{allowable-localization-lemma} replaced by
Lemma~\ref{lemma-localization-branch-omission-Mn} and with the canonical
objects $L,\vec S,D,\varrho_L$ replaced by
$M_n,\vec S^{\,n},\mathcal D^n,\varrho_n$.

If $u=u_\beta$ is intentionally coded at stage $\beta$, then the real added by
the corresponding almost disjoint coding factor witnesses $\Phi_n(u)$.
The persistence of this witness under the remaining forcing follows from the
small-forcing projective absoluteness established above.

Conversely, suppose that $r$ witnesses $\Phi_n(u)$ in the final extension and
that $u$ is not intentionally coded.  Apply
Lemma~\ref{lemma-localization-branch-omission-Mn} to $r$ and $u$, and
take the union of the resulting supports.  Thus there are countable sets
$E,I,B$ such that $r,u$ belong to the localized extension determined by these
supports.  As in the $L$-case, the coding area decoded from $r$ is either
almost disjoint from all retained areas $h_\beta$ with $\beta\in I$, or it is
one of them.  In the first case choose a fresh block outside all retained
areas and outside the countable branch support $B$; in the second case use a
coordinate on which $u$ differs from the corresponding intentionally coded
real.  In both cases we obtain
\[
   \xi=\tau_u(\gamma,k)
   \notin
   B\cup\bigcup_{\beta\in I}\operatorname{Sel}(h_\beta,u_\beta).
\]
By Lemma~\ref{lemma-localization-branch-omission-Mn}, the tree
$S^n_\xi$ remains Suslin in the localized inner model.  On the other hand,
the witness $r$ decodes a code asserting that the selected branch through
$S^n_\xi$ exists.  Since $r$ belongs to the localized model, this gives a
cofinal branch through $S^n_\xi$ there, a contradiction.  Hence $u$ was
intentionally coded.

The length-$\omega_1$ case follows, as in the $L$-case, by observing that
nice names for the relevant reals have countable support and therefore appear
in a countable initial segment.
\end{proof}

Theorem~\ref{theorem-no-unwanted-Mn-codes} is the exactness statement needed
in the uniformization argument.  It completes the $M_n$-specific coding
input; the remaining fixed-point construction is transferred from the
$L$-case in the next subsection.


\subsection{Transfer of the fixed-point construction}
\label{subsection-Mn-transfer-fixed-point}

The preceding subsections isolate all points at which the proof over $L$
uses features of the constructible universe.  We now transfer the fixed-point
construction from Subsection~\ref{subsection-alpha-allowability} and the final
iteration from Section~\ref{section-L-final-uniformization}.  We spell out the
translation and the hypotheses used by the transfer, but do not repeat the
forcing-theoretic inductions.

Fix a recursive universal enumeration
\[
   \langle\varphi^n_m(v_0,v_1)\mid m<\omega\rangle
\]
of the $\Pi^1_{n+3}$ formulas with two free real variables, and put
\[
   A^n_m
   :=\{(x,y)\in(2^\omega)^2:\varphi^n_m(x,y)\}.
\]
As before, triples $(x,y,m)$ are identified with their fixed recursive real
codes.

\paragraph{The translation.}
Let $\mathfrak t_n$ denote the following simultaneous translation of the
objects used in the $L$-construction:
\[
\begin{array}{rcl@{\quad}rcl}
   L&\longmapsto&M_n,
   &W&\longmapsto&W_n,\\
   \operatorname{Br}&\longmapsto&\operatorname{Br}_n,
   &\vec S&\longmapsto&\vec S^{\,n},\\
   \mathbb C_L(\omega_1)&\longmapsto&\mathbb C_{M_n}(\omega_1),
   &\forceC^{\mathrm{area}}&\longmapsto&
      \forceC^{\mathrm{area}}_n,\\
   D&\longmapsto&\mathcal D^n,
   &\mathbb A_D&\longmapsto&\mathbb A^n,\\
   \text{the area map }\varrho_L&\longmapsto&\varrho_n,
   &\dot Y(\dot x,\dot z,\dot m,\dot g_\eta)&\longmapsto&
      \dot Y_n(\langle\dot x,\dot z,\dot m\rangle,\dot g_\eta),\\
   \Psi&\longmapsto&\Psi_n,
   &\Phi&\longmapsto&\Phi_n,\\
   A_m&\longmapsto&A^n_m,
   &<_L&\longmapsto&\prec_n.
\end{array}
\]
Here $\prec_n$ is the canonical order on $M_n$-codes.  Every object in
$H(\omega_2)^{W_n}$ is represented by its $<_{M_n}$-least
$\operatorname{Br}_n$-name in $M_n$, and codes are compared in the canonical
construction order of $M_n$.  The same convention is used for forcings,
complete embeddings, conditions, and names which occur in localized support
pairs.

The semantic least-key selection is translated at the same time.  If
$\mathfrak a=(E,A)$ is an admissible support pair in an $M_n$-coding
iteration, write
\[
   M^n_{\mathfrak a}:=W_n[K_E][G_A].
\]
For $z\in M^n_{\mathfrak a}$, where in the application
$z=\langle x,y,m\rangle$, a presentation of $z$ below $\mathfrak a$ is defined
as in the $L$-case, replacing allowable forcing by $M_n$-allowable forcing and
requiring the unused coding areas to be disjoint from those used by
$\mathfrak a$.  Thus a presentation determines a pointed $M_n$-allowable
presentation
\[
   d=(\forceP_{\mathfrak b}*\dot\forceR,
      (1_{\forceP_{\mathfrak b}},\dot p),\dot\sigma).
\]
Two such pointed presentations are identified if they have a common
$M_n$-allowable refinement below which the pushed-forward names are forced to
be equal.  Let
\[
   \pcode_n(d)
   :=
   \min_{\prec_n}\{\operatorname{code}(e):
      e\equiv^n_{\mathrm{pres}} d\}.
\]
Then
\[
   \mathcal O^n_{\mathfrak a}(z)
   :=\{\pcode_n(d_c):c\text{ is a presentation of }z
      \text{ below }\mathfrak a\}
\]
and
\[
   \key^n_{\mathfrak a}(z)
   :=\min_{\prec_n}\mathcal O^n_{\mathfrak a}(z).
\]
The proof of Lemma~\ref{lemma-semantic-orbit-support-independence} translates
verbatim: if $\mathfrak b\preccurlyeq\mathfrak a$ and
$z\in M^n_{\mathfrak b}$, then
\[
   \mathcal O^n_{\mathfrak b}(z)=\mathcal O^n_{\mathfrak a}(z)
   \quad\text{and}\quad
   \key^n_{\mathfrak b}(z)=\key^n_{\mathfrak a}(z).
\]
Consequently reservoir reindexing, product normalization, and allowable
pullback do not change the key.  The selection criterion in the
$M_n$-construction is therefore the least
$\key^n_{\mathfrak a}(\langle x,y,m\rangle)$, not the least local name for $y$.
The translation also has the following syntactic conventions.
\begin{enumerate}
   \item Requirements that an area assignment, reservoir permutation, or
         canonical reindexing belong to $L$ become the corresponding
         requirements over $M_n$.

   \item A stage which in the $L$-construction codes $(x,z,m)$ now uses the
         factor
         \[
            \mathbb A^n\!\left(
               Y_n(\langle x,z,m\rangle,g_\eta)
            \right),
         \]
         or the corresponding forcing name.  Thus all stage forcings are
         instances of the coding mechanism from
         Subsection~\ref{subsection-adapted-Mn-coding-mechanism}.

   \item The role of Shoenfield absoluteness for a
         $\Pi^1_2$ coding witness is played by
         Theorem~\ref{theorem-small-two-step-projective-absoluteness}, since
         $\Psi_n$ is $\Pi^1_{n+2}$.  The upward persistence of a forced
         nonmembership statement, which is $\Sigma^1_{n+3}$, is supplied by
         Corollary~\ref{corollary-persistence-Sigma-n-plus-3-witnesses}.

   \item The exactness assertion for the coding predicate is supplied by
         Theorem~\ref{theorem-no-unwanted-Mn-codes}.
\end{enumerate}

For completeness, we record the translated recursive notion.  If
\[
   \forceP_\beta
   =\forceC^{\mathrm{area}}_n*\dot{\forceR}_\beta
\]
is an initial segment of an $M_n$-coding iteration and $\vec\tau$ is a finite
tuple of stage-respecting names, an \emph{admissible support pair} for
$\vec\tau$ is a pair
\[
   (E,A)\in[U_\beta]^{\leq\omega}\times[\beta]^{\leq\omega}
\]
such that $E$ consists of reservoir coordinates, $A$ consists of iteration
stages, $A$ is closed under
the names defining the retained iterands, and the induced forcing on $(E,A)$
is a complete subforcing of $\forceP_\beta$ over which every member of
$\vec\tau$ is a name.  We use the $\prec_n$-least admissible pair.  Likewise,
\[
   \forceQ\succeq^n_{\mathrm{res}}\forceP
\]
means that, after an $M_n$-coded reindexing of the reservoir onto a block of
coordinates and an increasing reindexing of a closed set of stages,
$\forceP$ is isomorphic to a complete subiteration of $\forceQ$.  This is the
$\mathfrak t_n$-translation of the complete reservoir extension relation from
Subsection~\ref{subsection-alpha-allowability}.

\begin{definition}
\label{definition-Mn-alpha-allowability}
A $0$-allowable iteration over $W_n$ is an $M_n$-coding iteration over the
reservoir in the sense of
Definition~\ref{definition-Mn-coding-iteration}, paired with the canonical
name for the empty auxiliary set.

Assume that $0<\rho$ and that $\theta$-allowability has been defined for every
$\theta<\rho$.  Let
\[
   F=\langle(\dot x_\beta,\dot z_\beta,\dot m_\beta)
      \mid\beta<\delta\rangle\in W_n
\]
be stage-respecting bookkeeping, and let
\[
   \vec\eta=\langle\eta_\beta\mid\beta<\delta\rangle\in M_n
\]
be injective, as in Definition~\ref{definition-Mn-coding-iteration}.  Put
\[
   U_\beta:=\operatorname{ran}(\vec\eta\restriction\beta).
\]
At stage $\beta$, let $(E_\beta,A_\beta)$ be the $\prec_n$-least admissible
support pair for the local bookkeeping names
\[
   (\dot x_\beta,\dot z_\beta,\dot m_\beta),
\]
and require $E_\beta\subseteq U_\beta$.  Put
\[
   \mathfrak a_\beta:=(E_\beta,A_\beta).
\]
Thus the new area coordinate $\eta_\beta$ is not used by the local data.  If
$K*G_\beta$ is generic for the current initial segment, write
\[
   M^n_{\mathfrak a_\beta}=W_n[K_{E_\beta}][G_{A_\beta}]
\]
and let $x,z,m$ be the interpretations of the bookkeeping names in this local
model.

For $\theta<\rho$ and $y\in M^n_{\mathfrak a_\beta}\cap2^\omega$, say that
$y$ is $\theta$-persistent for $(x,m)$ over $\forceP_\beta$ if
\[
   (x,y)\in A^n_m
\]
in the current extension and, for every $M_n$-$\theta$-allowable iteration
$\forceQ\in W_n$ together with a witness to
\[
   \forceQ\succeq^n_{\mathrm{res}}\forceP_\beta,
\]
the quotient over the selected copy of $\forceP_\beta$ forces
\[
   (x,y)\in A^n_m.
\]
Let
\[
\begin{aligned}
   C^{n,\rho}_\beta
   :=\{(\theta,y)\in
      \rho\times(M^n_{\mathfrak a_\beta}\cap2^\omega):{}&
      y\text{ is }\theta\text{-persistent}\\
      &\text{for }(x,m)\text{ over }\forceP_\beta\}.
\end{aligned}
\]
If $C^{n,\rho}_\beta$ is nonempty, choose first
\[
   \theta_\beta
   :=\min\{\theta:\exists y\ ((\theta,y)\in C^{n,\rho}_\beta)\},
\]
and then choose $y_\beta$ with $(\theta_\beta,y_\beta)\in C^{n,\rho}_\beta$
such that
\[
   \key^n_{\mathfrak a_\beta}(\langle x,y_\beta,m\rangle)
\]
is $\prec_n$-least among all values of rank $\theta_\beta$.  Record
$(x,y_\beta,m,\theta_\beta)$ in the auxiliary name.  If $z\neq y_\beta$, put
$z_\beta^*=z$.  If $z=y_\beta$, let $z_\beta^*$ be a real
$w\neq y_\beta$ in $M^n_{\mathfrak a_\beta}$ for which
\[
   \key^n_{\mathfrak a_\beta}(\langle x,w,m\rangle)
\]
is $\prec_n$-least.  The next factor intentionally codes
$(x,z_\beta^*,m)$, using $\prec_n$-least stage-respecting names only to form
the forcing name.

If $C^{n,\rho}_\beta$ is empty, the next factor intentionally codes the
bookkeeping triple $(x,z,m)$ and no tentative value is added.  At limit stages
take the finite-support limit of the almost disjoint tail and the union of
the auxiliary names.  All choices are made from $\prec_n$ and are uniform in
the generic filter.

A forcing is \emph{$M_n$-$\rho$-allowable} if it is isomorphic to the final
forcing of an iteration obtained by this recursion.  We denote the class of
such forcings by
\[
   \mathcal A^{W_n}_\rho.
\]
\end{definition}

The definition is exactly the $\mathfrak t_n$-translation of
Definition~\ref{definition-alpha-allowability}.  In particular, it uses the
same stage-respecting support convention and the same complete-reservoir
notion of continuation.  A \emph{relative $M_n$-$\rho$-allowable tail over
$\forceP$} is a quotient $\forceQ/\dot G_{\forceP}$, together with its
presentation, where $\forceQ$ is $M_n$-$\rho$-allowable and
$\forceQ\succeq^n_{\mathrm{res}}\forceP$.

\begin{theorem}
\label{theorem-transfer-fixed-point-to-Mn}
The definitions and conclusions of
Subsection~\ref{subsection-alpha-allowability} and
Section~\ref{section-L-final-uniformization} remain valid after applying
$\mathfrak t_n$.  More explicitly, the following statements hold in $W_n$.
\begin{enumerate}
   \item If $\xi<\alpha$, then every $M_n$-$\alpha$-allowable forcing is
         $M_n$-$\xi$-allowable.  Hence
         \[
            \mathcal A^{W_n}_0\supseteq
            \mathcal A^{W_n}_1\supseteq\cdots\supseteq
            \mathcal A^{W_n}_\alpha\supseteq\cdots.
         \]

   \item Relative allowable tails may be appended to an allowable iteration.
         If $\forceP^0$ and $\forceP^1$ are $M_n$-$\alpha$-allowable, then
         \[
            \forceP^0\times\forceP^1
         \]
         is isomorphic to an $M_n$-$\alpha$-allowable forcing.

   \item Tentative values remain in their intended sections.  More precisely,
         suppose that $(\forceP,\dot I)$ is $M_n$-$\alpha$-allowable,
         $G$ is $\forceP$-generic over $W_n$, and
         \[
            (x,y,m,\xi)\in\dot I^G
            \qquad(\xi<\alpha).
         \]
         Then $W_n[G]\models(x,y)\in A^n_m$.  Moreover, if $\forceQ$ is
         $M_n$-$\xi$-allowable and is supplied with a witness to
         $\forceQ\succeq^n_{\mathrm{res}}\forceP$, then
         \[
            W_n[G]\models
            ``\forceQ/G\Vdash(x,y)\in A^n_m."
         \]
         The same conclusion holds for an $\omega_1$-length iteration
         whose proper initial segments are allowable.

   \item Every class $\mathcal A^{W_n}_\alpha$ is nonempty, and the decreasing
         hierarchy stabilizes.  If
         \[
            \alpha^n_0
            :=\min\left\{\alpha:
               \forall\beta\geq\alpha\,
               \bigl(\mathcal A^{W_n}_\beta
                     =\mathcal A^{W_n}_\alpha\bigr)
            \right\}
         \]
         and
         \[
            \alpha_n^\ast:=\alpha^n_0+1,
         \]
         then the $M_n$-$\alpha_n^\ast$-allowable forcings form the fixed
         class $\mathcal A^{W_n}_{\alpha^n_0}$.  We call them
         \emph{$M_n$-$\infty$-allowable}.  Consequently, if a local real
         $y\in A^n_{m,x}$ fails every persistence test of rank below
         $\alpha_n^\ast$, there is an $M_n$-$\infty$-allowable complete
         reservoir extension whose quotient forces $(x,y)\notin A^n_m$.

   \item There is an $\omega_1$-length forcing
         \[
            \forceP^n_{\omega_1}
            =\forceC^{\mathrm{area}}_n*
               \dot{\forceR}^{\,n}_{\omega_1}
         \]
         together with an auxiliary name $\dot I^n$ such that every countable
         initial segment is $M_n$-$\infty$-allowable.  The forcing has size
         $\aleph_1$, preserves cardinals and $\CH$, and every real in
         its generic extension belongs to a countable intermediate stage.
         Whenever $K*G_{\omega_1}$ is generic over $W_n$, the final model
         \[
            V_n^\ast:=W_n[K][G_{\omega_1}]
         \]
         satisfies
         \[
            \Phi_n(u)
            \quad\Longleftrightarrow\quad
            u\text{ was intentionally coded by }
            \forceP^n_{\omega_1}.
         \]

   \item Writing
         \[
            I^n:=(\dot I^n)^{K*G_{\omega_1}},
         \]
         exactly one of the following alternatives holds in $V_n^\ast$ for
         every real $x$ and every $m<\omega$:
         \begin{enumerate}
            \item if $(x,y,m,\xi)\in I^n$ for some $y$ and
                  $\xi<\alpha_n^\ast$, then there is a unique real $y_0$
                  such that
                  \[
                     (x,y_0)\in A^n_m
                     \quad\text{and}\quad
                     \neg\Phi_n(\langle x,y_0,m\rangle);
                  \]
                  every other member $y$ of the section satisfies
                  $\Phi_n(\langle x,y,m\rangle)$;

            \item if no tentative value for $(x,m)$ occurs in $I^n$, then
                  \[
                     A^n_{m,x}=\varnothing.
                  \]
         \end{enumerate}
\end{enumerate}
\end{theorem}

\begin{proof}
We indicate precisely where the $M_n$-analogues established above supply the
facts used in the $L$-construction.

First consider the forcing-theoretic part of the proof.  The reservoir
$\forceC^{\mathrm{area}}_n$ has size $\aleph_1$, is
$\omega$-distributive over $W_n$, and admits every coordinate permutation
belonging to $M_n$, by
Lemma~\ref{lemma-basic-properties-Mn-Cohen-reservoir}.  Each coding factor is
Knaster and has size $\aleph_1$, by
Lemma~\ref{lemma-basic-properties-Mn-coding-forcing}.  Admissible countable
supports exist by the same nice-name and recursive-closure argument used over
$L$: the reservoir adds no reals, the almost disjoint tail is c.c.c., and the
countable supports are recursively enlarged so that every retained iterand is
already named over the retained earlier stages and retained reservoir
coordinates.  Consequently all localized subiterations and quotient forcings
used in the original proof are available over $W_n$.

The proofs of shrinking, the allowable-tail lemma, and nonemptiness are purely
formal inductions on the rank and the length of the iteration.  Replacing the
bookkeeping instruction at a stage by the instruction actually used at a
higher rank produces the lower-rank construction exactly as in the proof of
Lemma~\ref{shrinkinglemma}.

For products, fix in $M_n$ a partition
\[
   \omega_1=B_0\mathbin{\dot\cup}B_1
\]
and bijections $e_i:\omega_1\to B_i$.  Reindex the reservoirs of the two
factors onto $B_0$ and $B_1$ and concatenate their finite-support tails.  This
normalizes the product over one reservoir.  The induction from the proof of
Lemma~\ref{alphaallowableclosedunderproducts} applies verbatim: a real
fails a persistence test in one factor precisely when it fails the
corresponding test in the normalized product, because complete reservoir
embeddings compose and because the induction hypothesis makes the product of
the other factor with an eliminating forcing allowable.  At stages where the
persistent-value case applies, the translated local supports are related by
an $M_n$-coded reservoir reindexing and an allowable pullback.  Hence the
local reals have the same persistence ranks, and the support-independence of
$\key^n$ gives the same semantic keys in the original factor and in the
normalized product.  The product construction therefore makes the same
stage choice.  This proves (1) and (2).

The proof that tentative values remain in their sections is likewise the
proof of Lemma~\ref{fvaluesremain2}.  The selected value is persistent at its
recorded rank, shrinking puts every relevant later initial segment into that
rank, and complete reservoir extensions compose.  For an $\omega_1$-length
iteration whose proper initial segments are allowable, suppose toward a contradiction that membership fails.
Write this $\Sigma^1_{n+3}$ statement as
$\exists z\,\theta(z,x,y)$ with $\theta$ $\Pi^1_{n+2}$.  A name for the
witness $z$ has countable support and hence occurs in a bounded initial
segment.  Theorem~\ref{theorem-small-two-step-projective-absoluteness}
reflects $\theta(z,x,y)$ from the final extension to that bounded stage,
contradicting persistence there.  This proves (3).

Every allowable forcing has size at most $\aleph_1$ and belongs, up to
isomorphism, to the set $H(\omega_2)^{W_n}$.  Thus there are only set many
isomorphism types.  The nonempty decreasing hierarchy must therefore
stabilize, giving (4).  The successor $\alpha_n^\ast=\alpha^n_0+1$ is used so
that the stage rule tests persistence against the fixed class itself.

We next transfer the final bookkeeping construction.  Fix in $W_n$ a map
\[
   F:\omega_1\longrightarrow H(\omega_1)^{W_n}
\]
which enumerates every stage-respecting system of names and every countable
admissible enlargement of its support unboundedly often.  Use the canonical
area assignment $\eta_\beta=\beta$, which belongs to $M_n$.  At a bookkeeping
task for $(x,z,m)$, do nothing if the task is not legitimate or if the current
extension does not satisfy $(x,z)\in A^n_m$.  Otherwise form the local
admissible support $\mathfrak a_\beta$ and apply the stage rule from
Definition~\ref{definition-Mn-alpha-allowability}: if some real passes a
persistence test below $\alpha_n^\ast$, choose the least rank and then the
real of $\prec_n$-least key
\[
   \key^n_{\mathfrak a_\beta}(\langle x,y,m\rangle).
\]
If no real passes, failure at the fixed rank $\alpha^n_0$ supplies an
$M_n$-$\infty$-allowable complete reservoir extension whose quotient forces
\[
   (x,z)\notin A^n_m.
\]
Restrict below the least condition forcing this statement and append the
resulting relative tail.  Such a tail may have countable length; concatenate
the countable blocks and reindex them.  Since a sum of $\omega_1$ countable
ordinals has order type at most $\omega_1$, the result has the displayed
presentation $\forceP^n_{\omega_1}$.  The induction just proved
shows that every countable initial segment is $M_n$-$\infty$-allowable.
Lemma~\ref{lemma-size-preservation-Mn-coding-iterations} gives the size and
preservation assertions, and
Theorem~\ref{theorem-no-unwanted-Mn-codes} gives the displayed equivalence
with intentional coding.  This proves (5).

It remains to check the final dichotomy.  By the size and preservation
statements already proved, every pair of successive intermediate models used
below forms a small two-step extension in the sense of
Definition~\ref{definition-small-Mn-generic-extension}.  We use the following
elementary consequence of
Theorem~\ref{theorem-small-two-step-projective-absoluteness}.  If
$N\subseteq N[H]$ is one of the relevant two-step extensions,
$x,y\in N$, and
\[
   N[H]\models(x,y)\in A^n_m,
\]
then $N\models(x,y)\in A^n_m$.  Indeed, write
\[
   \varphi^n_m(x,y)\equiv
   \forall r\,\vartheta(r,x,y)
\]
with $\vartheta$ $\Sigma^1_{n+2}$.  For every $r\in N$, the extension
satisfies $\vartheta(r,x,y)$, and
Theorem~\ref{theorem-small-two-step-projective-absoluteness} reflects this
statement to $N$.

Suppose first that the auxiliary set contains a tentative value for $(x,m)$.
Let $\xi_0<\alpha_n^\ast$ be the least rank occurring for $(x,m)$.  Among all
occurrences of rank $\xi_0$, choose one whose semantic key is $\prec_n$-least;
write this key as $\kappa_0$, let $y_0$ be the selected real, and let $\beta$
be the first stage at which such an occurrence is added.  Clause~(3) gives
$(x,y_0)\in A^n_m$ in the final model.

We use the following no-reversal observation.  Suppose that, at a comparable
stage with local support $\mathfrak a_\eta$, the construction considers the
same pair $(x,m)$ and selects $y'\neq y_0$ at rank $\xi'\leq\xi_0$.  If
$y_0\in M_{\mathfrak a_\eta}$, then support-independence of $\key^n$ and
stability of eliminability imply that $y_0$ is available there at rank
$\xi_0$.  Hence the stage rule could choose $y'$ only if either
$\xi'<\xi_0$, or $\xi'=\xi_0$ and
\[
   \key^n_{\mathfrak a_\eta}(\langle x,y',m\rangle)
   \prec_n
   \key^n_{\mathfrak a_\eta}(\langle x,y_0,m\rangle).
\]
By support-independence, the second key is $\kappa_0$, and the first key is
the key carried by the occurrence of $y'$ at stage $\eta$.  Thus both
alternatives contradict the choice of $(\xi_0,\kappa_0)$ among the occurrences
in the final auxiliary set.  No such reversal occurs.

It follows that the triple $\langle x,y_0,m\rangle$ is never intentionally
coded.  Before stage $\beta$, any coding of this triple would have resulted
from selecting a different value for $(x,m)$ with rank at most $\xi_0$, which
is excluded by the no-reversal observation.  From stage $\beta$ on, whenever
$y_0$ is visible, the same observation forces the construction to protect
$y_0$ rather than code it.

Let $y\neq y_0$ and suppose that $(x,y)\in A^n_m$ in the final model.  At a
sufficiently late bookkeeping stage, the local support contains $x$, $y$, and
$y_0$.  The preceding downward-absoluteness observation shows that this local
model already regards $(x,y)$ as a member of $A^n_m$.  The elimination case is
impossible, because the resulting $\Sigma^1_{n+3}$ nonmembership would persist
to the final model by
Corollary~\ref{corollary-persistence-Sigma-n-plus-3-witnesses}.  Hence the
persistent-value case applies; by no reversal it protects $y_0$, and therefore
intentionally codes $\langle x,y,m\rangle$.  Exactness of the coding predicate,
supplied by Theorem~\ref{theorem-no-unwanted-Mn-codes}, yields the first
alternative in (6).

Finally suppose that no tentative value for $(x,m)$ occurs and that, toward a
contradiction, $(x,y)\in A^n_m$ in the final model.  At a sufficiently late
bookkeeping stage the names for $x$ and $y$ occur together.  Downward
absoluteness puts $(x,y)$ in the section already at that stage.  The
persistent-value case would add a tentative value, so the elimination case
must apply.  It forces $(x,y)\notin A^n_m$, and this nonmembership persists to
the final extension by
Corollary~\ref{corollary-persistence-Sigma-n-plus-3-witnesses}, a
contradiction.  This proves (6) and completes the transfer.
\end{proof}

For $m<\omega$, define
\[
   U^n_m(x,y)
   \quad\Longleftrightarrow\quad
   (x,y)\in A^n_m
   \ \wedge\
   \neg\Phi_n(\langle x,y,m\rangle).
\]

\begin{corollary}
\label{corollary-Pi-n-plus-3-uniformization-final-model}
In $V_n^\ast$, the relation $U^n_m$ is a
$\Pi^1_{n+3}$ uniformization of $A^n_m$ for every $m<\omega$.
Consequently, $V_n^\ast$ satisfies the
$\Pi^1_{n+3}$-uniformization property.
\end{corollary}

\begin{proof}
By Lemma~\ref{lemma-complexity-Psi-Phi-Mn}, the predicate $\Phi_n$ is
$\Sigma^1_{n+3}$.  Hence $U^n_m$ is $\Pi^1_{n+3}$, and
Theorem~\ref{theorem-transfer-fixed-point-to-Mn}(6) gives the domain and
uniqueness requirements.  The enumeration
$\langle A^n_m\mid m<\omega\rangle$ may be chosen universal; real parameters
are absorbed into the first coordinate using the fixed recursive pairing
function.  Thus the conclusion holds for every boldface
$\Pi^1_{n+3}$ relation.
\end{proof}

\begin{theorem}
\label{theorem-force-Pi-n-plus-3-uniformization-over-Mn}
Let $1\leq n<\omega$, and assume that the canonical inner model $M_n$
exists.  Then there is a cardinal-preserving set-generic extension of $M_n$
which satisfies $\CH$ and the $\Pi^1_{n+3}$-uniformization property.
\end{theorem}

\begin{proof}
First force over $M_n$ with the branch product producing $W_n$, and then with
$\forceP^n_{\omega_1}$.  By
Theorem~\ref{theorem-properties-of-Wn} and
Theorem~\ref{theorem-transfer-fixed-point-to-Mn}(5), both steps preserve
cardinals and $\CH$.  The resulting model is $V_n^\ast$, and the
uniformization property follows from
Corollary~\ref{corollary-Pi-n-plus-3-uniformization-final-model}.
\end{proof}


\section{Further possible applications and open problems}
We close with two directions suggested by the construction.  First, the method
should have an analogue in the generalized Baire-space context at a regular
uncountable cardinal $\kappa$.  In particular, one may ask whether the
$\Pi^1_1$-uniformization problem for $\kappa^{<\kappa}$ can consistently have
a positive solution.

The methods in this paper leave the following question open:
\begin{question}
Can the $\Pi^1_n$-uniformization property be forced over $L$ for $n >3$?
\end{question}

The method introduced here is so far limited to local effects.  It would be
interesting to force a less local, or even global, behaviour:
\begin{question}
Given a pair $n,m\in\omega$ such that $n\neq m$ and $n\neq m+1$, can one force
a universe in which the $\Pi^1_n$- and the $\Pi^1_m$-uniformization properties
both hold?
\end{question}

Finally, it seems worthwhile to combine the present method with other forcing
constructions in order to obtain models with properties that are usually
consequences of determinacy assumptions.  We record one paradigmatic question:
\begin{question}
Assume some sufficiently big large cardinal notion which can exist in $L$.  Is there a model which satisfies the
$\Pi^1_3$-uniformization property and in which every projective set of reals is
Lebesgue measurable?
\end{question}

\bibliographystyle{plain}
\bibliography{references}
\end{document}

